% =====================================================================
%  Harald Høffding, PSYKOLOGISKE UNDERSØGELSER.  Det Kongelige Danske
%  Videnskabernes Selskabs Skrifter, 6. Række, historisk og philosophisk
%  Afdeling, III, no. 1.  Kjøbenhavn: Bianco Lunos Kgl. Hof-Bogtrykkeri
%  (F. Dreyer), 1889.  103 pp.
%  TRANSCRIPTION (Danish) of the 1889 print, complete (pp. 1-103).
%  Built 2026-10-09 by the e-book method; see below and RESIDUAL RISKS.
%
%  SOURCE: Google's scan of the whole Academy volume (Stanford copy, Google
%  Books hgMNAQAAIAAJ), ~/bibliotek/Høffding, Harald/
%  1889-1895-vidensk-selsk-skrifter-6r-hist-filos-III.pdf (SCANS.tsv).
%  PAGE MAP (pagemap.py): printed page = PDF page - 24; p. 1 = PDF 25 ...
%  p. 103 = PDF 127.  p. 1 title, p. 3 Indhold, text pp. 5-103.  Not
%  transcribed: folios, the series line at each sheet's foot („Vidensk.
%  Selsk. Skr., 6. Række, historisk og philosophisk Afd. III. 1.“ + sheet
%  number), Google's watermark, library stamps.  Every page turn is marked
%  "% --- p. N ---" + \opage{N} at the first word of the printed page; a
%  word divided over a page turn ends in %.
%
%  CONVENTIONS: letterspacing (the print's emphasis) -> \emph; italic
%  (chiefly the letters a, b, A, B of the formulas) -> \textit; subscripts
%  (a₁, p₂) -> \textsubscript; printed dashes -> ---; quotation marks «…»
%  as printed.  Footnotes at the page foot, numbered 1), 2) on each page;
%  \footnote[k] gives the note's number on its printed page.  A note that
%  runs over to the next page is not split.
%
%  METHOD (the e-book method, TRANSCRIPTION-PLAYBOOK §00 point 0).  No one
%  typed the text.  The base text -- words, punctuation, paragraphs, italics,
%  the footnote texts -- is the Lindhardt og Ringhof e-book
%  (Psykologiske_undersoegelser.epub, 2023: the 1889 text in the 1889
%  spelling, the 89 footnotes set as endnotes).  ebook-method/
%  psykologiske-undersoegelser/build.py aligned it letter by letter with
%  Google's text layer of the scan, split per page into text and footnote
%  lines by line pitch, and took every page turn from the layer.  Of 306
%  letter disagreements 200 were settled by rule; 106 letter, 24
%  punctuation and 11 paragraph disagreements were decided at the image
%  (decisions.tsv).  Letterspacing, which neither the e-book nor the layer
%  marks, was measured from the width of each word's box against a width
%  fitted per character: wide words are \emph by rule (calibrated on 60
%  words read at the image), the doubtful band (150 words) was read at the
%  image; 2-3 letter words are spaced only between two spaced words.
%  Edits that are not a single word are in patches.tsv (the stacked
%  formula „(a over A)“, pp. 12, 16, 21, 22; words the e-book glued).
%  The print corrects the e-book at: p. 14 Bekendthedskvalitet; p. 15 at;
%  p. 18 Berøringsassociationens; p. 19 Forvexling; p. 31 adskilte; p. 36
%  1; p. 41 Lighedsassociation; p. 44 sér; p. 53 Genkendelsesakter; p. 61
%  Opmærksomhedsovergang; p. 66 Vurderéevne; p. 69 Følelsestilstand; p. 84
%  naturligvis; p. 85 vigtigste; p. 88 Sanseánskuelse (accent uncertain);
%  notes: p. 9 sensus; p. 14 Aufl., zur; p. 46 Høreforestillingernes;
%  p. 65 Revue; p. 68 même, aperçoit; p. 96 183; and p. 31 „556)“, p. 23
%  „Philosophie IX)“, p. 50 „p 278“; pp. 43, 54, 89, 92 spaces and hyphens
%  the e-book dropped („kalder Ensartethed“, „Bevidstheds- eller“,
%  „dogmatisk-metafysiske“, „Anger- eller“).  PRINTED AS IS: p. 44
%  „Lighedassociation“; p. 54 n. „onother“ (Locke quotation); and, found by
%  the translators and read at the image: p. 8 „nogenside“, p. 9 „Fordel“
%  (for Forskel), p. 39 „han“ (for kan), p. 40 „savne“ (probably for sanse),
%  p. 47 „ønker“, p. 74 „at at“.  REPAIRED WHILE TRANSLATING (the
%  translators' odd spots, TRANSLATION-PLAYBOOK §6c): p. 33 «Lighedsteorien»,
%  p. 36 „V B, 1“, p. 47 n. „før“, p. 51 n. „Il“, p. 60 n. „Wundt's“, p. 73
%  „Punkt,“, p. 84 n. (a stray point), p. 95 „kunde“.
%
%  RESIDUAL RISKS: (1) a misprint that the e-book and the layer both
%  normalise is invisible to the method (an unattested-word sweep was run;
%  p. 30 was read whole against the image); (2) italics are the e-book's;
%  (3) on pp. 8, 32, 64 and 68 the print shows a round point before a
%  lower-case word where a comma is wanted -- set as commas (a comma whose
%  tail did not print), to be confirmed; (4) p. 15 „at [kalde“: the print
%  has an opening bracket-like mark with no closing one, kept as the
%  e-book's „[“; (5) letterspacing of short words is inferred, not read.
% =====================================================================
\documentclass[12pt,a4paper,openany]{book}
\usepackage[utf8]{inputenc}
\usepackage[T1]{fontenc}
\usepackage{textalpha}
\usepackage{libertinus}
\usepackage{libertinust1math}
\usepackage{graphicx}    % for \reflectbox (the old Danish ɔ: id-est mark)
\DeclareUnicodeCharacter{0254}{\reflectbox{c}}  % ɔ (U+0254) = reversed c
\usepackage[danish]{babel}
\usepackage[protrusion=true,expansion=true]{microtype}
\usepackage{setspace}
\usepackage[margin=1.2in]{geometry}
\usepackage{fancyhdr}
\setlength{\headheight}{28pt}
\addtolength{\topmargin}{-13.5pt}
\usepackage{hyperref}
\hypersetup{pdftitle={Psykologiske Undersøgelser}, pdfauthor={Harald Høffding}, hidelinks}
\pagestyle{fancy}
\fancyhf{}
\fancyhead[LE,RO]{\thepage}
\fancyhead[RE]{\textit{Harald Høffding: Psykologiske Undersøgelser}}
\fancyhead[LO]{\textit{1889}}
\setstretch{1.3}
\usepackage{marginnote}
\newcommand{\opage}[1]{\marginnote{\footnotesize\textit{[#1]}}}
% Footnotes as the print marks them, numbered per page: 1), 2), 3) ...
\renewcommand{\thefootnote}{\arabic{footnote})}
\newcommand{\tocline}[3]{\noindent\makebox[2.2em][r]{#1}\hspace{0.6em}#2\dotfill\ #3\par}

\begin{document}

% --- p. 1 --- (title)
\opage{1}
\begin{center}
\vspace*{3em}
{\Huge Psykologiske Undersøgelser.}\\[4em]
Af\\[3em]
{\large\textbf{Harald Høffding.}}\\[4em]
{\small Vidensk. Selsk. Skr., 6. Række, historisk og philosophisk Afd.\ \ III.\ 1.}\\[6em]
{\large Kjøbenhavn.}\\[2pt]
{\small Bianco Lunos Kgl. Hof-Bogtrykkeri (F. Dreyer).}\\[2pt]
{\small 1889.}
\end{center}
\clearpage

% --- p. 3 --- (contents)
% CHECK: in this copy the left edge of p. 3 is cut off in the gutter, so the
% first letters of each line (and any numerals before them) are not in the
% scan.  The titles below are supplied from the section heads in the text;
% the page numbers are legible.  Verify against HathiTrust njp.32101076896016
% (Princeton copy of the same volume) before calling the book done.
\opage{3}
\begin{center}{\large\textbf{Indhold.}}\end{center}
\medskip
\hfill{\small Side}\par
\tocline{}{Indledning}{5.}
\tocline{I.}{Umiddelbar Genkendelse}{8.}
\tocline{II.}{Berøringsassociationens Forudsætninger}{36.}
\tocline{III.}{Den frie Lighedsassociation}{41.}
\tocline{IV.}{Forestillingsassociation og Sammenligning}{59.}
\tocline{V.}{Om Begrebet sjælelig Aktivitet}{67.}
\tocline{VI.}{Sjælelig og fysisk Aktivitet}{81.}
\tocline{VII.}{Den sjælelige Aktivitets Lovmæssighed}{92.}
\clearpage

% ===== TEXT (pp. 5--103) =====
% --- p. 5 ---
\setcounter{footnote}{0}%
\opage{5}

\begin{center}\large Indledning.\end{center}

1. Læren om Forestillingernes Association indtager i den nyere Psykologi en
fremragende Plads. Under Bestræbelsen for at opdage bestemte Love for de
sjælelige Fænomeners Opstaaen og indbyrdes Forhold rettede Opmærksomheden sig
naturligt mod Forestillingerne, hvis Opstaaen og Forbindelser ere lettest at
paavise. Fornemmelsernes Betingelser og Forhold kunne ikke saa let findes ad
Selviagttagelsens Vej, men opdages ofte kun ved Hjælp af de fysiske Forhold,
til hvilke Fornemmelserne ere knyttede. Men Forestillingernes Optræden og
Forbindelse bero væsentligt paa Betingelser, som ere givne i selve
Bevidstheden. Hvad Følelse og Villie angaar, ere de ikke saa lette at
analysere; skønt enkelte Hovedfænomener fremtræde med slaaende Tydelighed, er
det dog forbundet med store Vanskeligheder for den uøvede Iagttagelse og
Reflexion at finde de enkelte Elementer og deres Forhold. Heraf er for en stor
Del den Ensidighed at forklare, med hvilken flere Retninger inden for
Psykologien have fremhævet Erkendelsen paa Bekostning af de andre Sider ved det
sjælelige Liv. En stor praktisk Betydning havde desuden Spørgsmaalet om Lovene
for Forestillingernes Forbindelse paa Grund af den uimodsigelige Indflydelse,
Forestillingerne have paa Følelse og Villie: ved Love for Forestillingernes
Forbindelse eller Association vilde man altsaa tillige vinde Indsigt i,
hvorledes Følelsens og Villiens Karakter kunde ændres. Baade den Grad og den
Retning, i hvilken det sjælelige Liv udvikler sig, bestemmes for en stor Del
ved Forestillingerne og deres indbyrdes Forbindelse. Man har gjort denne
Indflydelse større end den i Virkeligheden er; man har overset den stadige og
betydningsfulde Tilbagevirken, som fra Følelsens og Villiens Side øves paa
Forestillingernes Udvikling og Forbindelse; man har ikke været tilstrækkeligt
opmærksom paa de ejendommelige Forhold, som gælde for Følelsens og Villiens
Vedkommende.

Naar jeg i denne Afhandling, der skal drøfte forskellige i den nyeste Tid
omtvistede psykologiske Spørgsmaal, begynder med en Undersøgelse af nogle
Hovedpunkter i Associationslæren, er det ikke de nys nævnte Ensidigheder i
ældre Behandlinger af denne Lære, jeg vil vende mig imod. Denne Side af Sagen
haaber jeg at have belyst i min «Psykologi i Omrids paa Grundlag af Erfaring».
Derimod vil jeg søge at bestemme
% --- p. 6 ---
\setcounter{footnote}{0}%
\opage{6}hvad jeg vil kalde den \emph{nedre} og den \emph{øvre Grænse for
Begrebet om Forestillingernes Association}. --- Spørgsmaalet om dette Begrebs
nedre Grænse opstaar ved, at man ved Udtrykket Forestillingsassociation fra
først af forstod et saadant Forhold mellem Forestillinger, der hver for sig
traadte frem som frie og selvstændige Elementer i Bevidstheden, at naar den ene
var given, optraadte ogsaa den anden. Spørgsmaalet er, om Associationsbegrebet
ikke ogsaa med Rette finder Anvendelse i saadanne Tilfælde, hvor vi ikke
umiddelbart kunne skelne flere Elementer, men dog have Anledning til --- i
hvert Tilfælde teoretisk eller hypotetisk --- at opfatte Sagen saaledes, at
flere Elementer ere sammensmeltede, Elementer, som under andre Forhold vilde
gøre sig gældende frit, hvert for sig. Det er Genkendelsen, eller rettere en
vis Art af Genkendelse, som her bliver nærmere at undersøge. --- Spørgsmaalet
om den øvre Grænse opstaar, naar vi drøfte Forholdet mellem
Forestillingsassociationen og den sjælelige Aktivitet. Man har --- især i den
ældre engelske Skole og i den Herbartske Skole --- opfattet Associationen
saaledes, at den udelukkede en egentlig sjælelig Aktivitet, og naar man saa
tillige gjorde Associationen til Grundfænomenet i Psykologien, er det klart, at
Sjælelivet blev opfattet paa en højst ensidig Maade. Man gik ud fra, at
Associationen fandt Sted mellem selvstændige Elementer, lige som i den fysiske
Verden Tiltrækningen forbinder de selvstændige Elementer; denne Analogi fik en
uberettiget og overvældende Indflydelse paa den psykologiske Tankegang. Naar
andre psykologiske Retninger i Modsætning hertil --- uden at negte
Associationens faktiske Tilværelse --- hævdede det sjælelige Livs Enhed og
Aktivitet, saa bliver Spørgsmaalet: gives der to Arter af sjælelige Fænomener,
en, i hvilken den blotte Association raader, og en anden, i hvilken den
sjælelige Enhed og Aktivitet fremtræder, --- eller er maaske selve
Associationen kun en Form for den sjælelige Aktivitet, saa at Forskellen mellem
hvad vi kalde sjælelig Passivitet og Aktivitet kun er en Gradsforskel eller
beror paa Abstraktion?

Men Spørgsmaalet om Grænserne for Associationsbegrebets Anvendelse hænger igen
mere eller mindre tydeligt sammen med det Spørgsmaal, \emph{hvor mange Love for
Forestillingernes Association} der maa \emph{opstilles}. --- De ældste
Forskere, Platon og Aristoteles, som omtale Forestillingsassociationen, nævne
allerede flere Arter, af hvilke de vigtigste ere Association formedelst
Ligheden af Forestillingernes Indhold (Lighedsassociation) og Association
formedelst Sammentræf i Tid eller Rum af Forestillingernes Indhold
(Berøringsassociation). Den tredie Lov, Aristoteles nævner, Association
formedelst Kontrast, er man nu enig om ikke at betragte som en særegen Lov; de
Fænomener, den skulde passe paa, finde deres Forklaring ved de to andre Love%
% footnote 1) on p. 6
\footnote[1]{Smlgn. min Psykologi. 2. Udg. p. 185 f.}. Der har nu i den nyere
Psykologi gjort sig Forsøg gældende paa at reducere Lighedsassociationen til
Berørings%
% --- p. 7 ---
\setcounter{footnote}{0}%
\opage{7}association, saaledes at denne sidste blev den eneste Lov for
Forestillingsassociationen. Denne Teori fremtræder, hvad vel maa mærkes, i
forskellig Sammenhæng hos forskellige psykologiske Forskere. I den ældre
engelske Skole fremtræder den i Forbindelse med Forsøget paa at gøre
Forestillingsassociationen til hele Bevidsthedslivets Grundlov. Man reducerer
først alle Bevidsthedsprocesser til Associationer, og derved til Vexelforhold
mellem forskellige Elementer, i hvilke man tænker sig Bevidstheden opløst;
dernæst reducerer man al Association til at være en Forbindelse formedelst
udvortes Sammenhæng. Et saadant Standpunkt repræsenteres paa typisk Maade af
\emph{James Mill} i hans «Analysis of the human mind» (1829). Men denne sidste
Reduktion er ogsaa forsøgt af Forskere, som gøre en meget skarp Forskel mellem
lavere, mekaniske og højere, ideelle Sjæleprocesser, og som ved Siden af
Fornemmelse og Forestillingsassociation hævde en særegen Evne til Sammenlignen
eller Vurderen af Bevidsthedselementerne. Man ser her maaske endog en Fare i
Anerkendelsen af Lighedsassociationens Selvstændighed, da derved Fænomener, som
menes at vidne om udpræget sjælelig Aktivitet, synes at blive forklarede paa
rent mekanisk Maade. I hvert Tilfælde hævder man med Eftertryk Forskellen
mellem Associationen som en lavere og den sammenlignende og «forholdsættende»
Tænken som en højere Art af sjælelig Proces. Et saadant Standpunkt
repræsenteres af \emph{Hermann Lotze}.

Det er da af Interesse at undersøge Forholdet ogsaa fra denne Side. Men det
fører igen videre. Begrebet Aktivitet bruges i Psykologien alt for ofte paa en
noget dunkel Maade. Der skal da forsøges at give det en større Bestemthed end
det hidtil har havt. Et Forsøg i denne Retning vilde ligge saa meget des
nærmere, hvis det skulde vise sig, at Forestillingsassociationen helt igennem
er en Form for sjælelig Aktivitet; ti saa meget des større Trang vil der være
til at faa dette sidste Begreb undersøgt. Og derfra føres vi igen naturligt til
et andet stort Problem: hvilket er Forholdet mellem sjælelig Aktivitet og
fysisk Aktivitet? Endeligt vil det ogsaa tjene til fuld Belysning af den
sjælelige Aktivitets Natur at drøfte, hvor vidt den paa alle Punkter er
bestemte Love underlagt, eller om det skulde være berettiget at antage, at den
i større eller mindre Omfang er unddragen Aarsagslovens Herredømme.

2. Skønt denne Afhandling drøfter en hel Række forskellige Problemer, viser dog
den foregaaende Udvikling, at der er en vis Forbindelse mellem dem. Gangen i
Fremstillingen er ogsaa given ved det foregaaende. Jeg vil 1) først undersøge
den \emph{umiddelbare Genkendelse} og dens Forhold til
Forestillingsassociationen. Dernæst vil jeg 2) gaa over til at undersøge
\emph{Berøringsassociationens Forudsætninger} og forsøge at vise, at til disse
Forudsætninger hører ogsaa umiddelbar Genkendelse, der i hvert Tilfælde er en
med Lighedsassociation beslægtet Proces. Det bliver da 3) at afgøre, om
\emph{Lighedsassociationen} med Rette opstilles som en selvstændig Form for
Association
% --- p. 8 ---
\setcounter{footnote}{0}%
\opage{8}ved Siden af Berøringsassociationen. Efter at disse Hovedpunkter af
Associationslæren ere belyste, spørger jeg 4) om \emph{Forholdet mellem
Forestillingsassociation} (\emph{særligt Genkendelse og Lighedsassociation})
\emph{og sammenlignende Tankevirksomhed}, og kommer derved til at betragte 5)
den \emph{sjælelige Aktivitet} i det Hele og 6) \emph{dens Forhold} til den
\emph{fysiske Aktivitet}, med hvilken den erfaringsmæssigt er forbunden. Medens
dette Spørgsmaal fører os ud over det rent psykologiske Omraade til
Erkendelsesteorien og Metafysiken, vil det sidste Spørgsmaal, som skal
opkastes, 7) om \emph{sjælelig Aktivitets Lovbundethed}, føre os fra
Psykologien over i Etiken.

\begin{center}\large I. Umiddelbar Genkendelse.\end{center}

3. Lad os begynde med ganske simple og utvivlsomme Iagttagelser. Noget, jeg
ser, forekommer mig bekendt. Lad det være et Ansigt, eller for at tage noget
endnu simplere, et enkelt Ansigtstræk, et Blink med Øjet. Eller jeg ser paa
Aftenhimlen en Farvenuance, som vel er usædvanlig, men dog synes mig bekendt.
Eller der nævnes et fremmed Ord, jeg ikke kan oversætte, men hvis Lyd dog har
en bekendt Klang for mig. En spørger mig: «Har De nogenside %
% PRINTED AS IS: „nogenside“ for „nogensinde“ (p. 8)
været i Les Plans?» Navnet Les Plans er mig bekendt, og dog kan jeg ingen som
helst Forestilling forbinde med det. Eller lad os tage Exempler fra den indre
Erfaring. En vis organisk Fornemmelse, en vis Stemning indenfor Livsfølelsen,
som dukker frem hos mig, staar med et vist Fortrolighedens og Hjemlighedens
Præg for mig. Jeg har en aldeles umiddelbar Bevidsthed om, at jeg har havt dem
før, skønt jeg ikke kan knytte en eneste nærmere Bestemmelse til dem.

Nært beslægtede Tilfælde ere saadanne, i hvilke vi søge at fremkalde noget i
vor Erindring, uden at det lykkes os, medens vi strax, naar det ad anden Vej
fremtræder for os, genkende det, eller saadanne, i hvilke vi ikke ere sikre
paa, om vi have opfattet et enkelt Fænomen eller ikke, men blive sikre i vor
Sag ved at gennemgaa hele den Række af Fænomener, til hvilken det hører. Ogsaa
her faa vi et aldeles umiddelbart Indtryk af noget Bekendt. Paa en lærerig
Maade fremtræder dette ved \emph{Cattell's} Forsøg over, hvor mange Linier, Tal
eller Bogstaver der samtidigt kunne opfattes. «Naar», siger han%
% footnote 1) on p. 8
\footnote[1]{Ueber die Trägheit der Netzhaut und des Sehcentrums. (Wundt's
Philosophische Studien. III. Leipzig 1886) p. 127.}, «et vist Antal blev
overskredet, vare de opfattede Spor for svage og af for kort Varighed til, at
jeg deraf kunde gætte mig til de Objekter, der tjente som Indtryk. Uagtet det
nu var mig umuligt uden videre at angive disse Objekter, kunde jeg dog svare
rigtigt, naar jeg blev spurgt, om et vist Skrifttegn var der eller ikke. Det er
især Tilfældet ved lange
% --- p. 9 ---
\setcounter{footnote}{0}%
\opage{9}Sætninger; jeg har der den ejendommelige Følelse, som om jeg havde
kendt Sætningerne og igjen glemt dem». Hermed kan et pudsigt Exempel fra de
senere Aars talrige Experimenter med hypnotiserede Personer sammenstilles.
\emph{Charles Richet} hypnotiserede en af sine Venner og paanødte ham trods
nogen Modstræben den Forestilling, at han hed Durand, saa at han svarede med
dette Navn, naar man spurgte ham, hvad han hed. Naar Richet nu spurgte ham,
hvad han hed i vaagen Tilstand, var det ham ikke muligt at finde paa det, skønt
han vidste, at han «kendte» det. Men naar det blev nævnet for ham, kendte han
det strax og gentog det flere Gange med stort Velbehag%
% footnote 1) on p. 9
\footnote[1]{Le somnambulisme provoqué (i Skriftet: L'homme et l'intelligence.
Paris 1885). p. 540. --- Som Cattell's Forsøg synes at godtgøre, behøver man
ikke at have havt en tydelig, end sige med Opmærksomhed forbunden Opfattelse af
noget, for at det skal kunne genkendes. Det kan vel undertiden endog være
Spørgsmaal underkastet, om den Fornemmelse, der genkendes, første Gang har
været over Bevidsthedens Tærskel. Det gaar hermed som med de Fornemmelser, vi
slutte, at vi have havt, deraf at vi faa visse Kontrastfornemmelser; en saadan
Slutning er ikke aldeles sikker. Smlgn. min «Psykologi», 2. Udg. p. 87 og
foruden de der omtalte Fakta \emph{Machs} smukke Forsøg i hans «Beiträge zur
Analyse der Empfindungen». Jena 1886. p. 107. --- Allerede \emph{Ludvig Vives}
omtaler Genkendelsen af noget, der ikke blev opmærksomt opfattet. «Qvædam
recipit intelligentia illa simplex, qvæ est de iis, qvæ extrinsecus obveniunt
per videndi vel audiendi sensus, qvæ non \emph{animadversa} tradit statim
memoriæ. Attentio vero velut resurgens considerat in memoria et intelligit
interdum continuo post». De anima. Brugis 1538. p. 63. Hin «intelligentia
simplex» er aabenbart en ren Hypothese.}.

Hvad der er givet i saadanne Bevidsthedstilstande, som de her omtalte, er den
umiddelbare Opfattelse af Forskellen mellem noget bekendt og fortroligt og
noget nyt og fremmed. Denne Fordel %
% PRINTED AS IS: „Fordel“ for „Forskel“ (p. 9; the next line speaks of „Forskellen“)
er saa simpel og klar, at den lige saa lidt lader sig nærmere beskrive som f.
Ex. Forskellen mellem Lyst og Ulyst eller Forskellen mellem Gult og Blaat. Det
er en umiddelbar Kvalitetsforskel, vi her staa over for. Den ejendommelige
Kvalitet, med hvilken det Bekendte træder frem i Bevidstheden i Modsætning til
det Nye, vil jeg i det Følgende kalde \emph{Bekendthedskvaliteten}.
Spørgsmaalet er, hvorledes denne Kvalitet er at forklare, om den skal antages
for at være et simpelt og usammensat psykologisk Element, eller om den kan
føres tilbage til mere sammensatte Forhold.

Inden jeg forlader de anførte Tilfælde, som skulle danne Udgangspunkter for
Undersøgelsen, vil jeg endnu gøre opmærksom paa nogle Ejendommeligheder ved dem.

De Fænomener, som genkendes, ere her \emph{ikke usammensatte}, men dog saa
simple, at de \emph{samtidigt} træde frem for Bevidstheden. Et Blink af Øjet,
en Nuance paa Aftenskyen, en kort Række skrevne Bogstaver, et kort Ord, det er
Fænomener, ved hvilke der ikke kan være Tale om successiv Opfattelse. Forsøg
have desuden godtgjort, at vi samtidigt kunne opfatte 4 à 5 uforbundne
Synsindtryk (Linier, Tal, Bogstaver), og omtrent tre Gange saa mange, naar
Indtrykkene ere Dele af en som Helhed forud bekendt Forestilling. Tillige har
det vist sig, at Øvelse i denne Henseende udvider Bevidsthedens Om%
% --- p. 10 ---
\setcounter{footnote}{0}%
\opage{10}fang. Ord, ja endog Sætninger, optages i Bevidstheden som Helheder%
% footnote 1) on p. 10
\footnote[1]{\emph{Cattell} (i den anførte Afhandling) p. 127. Han siger
udtrykkeligt: «Als an mir selbst Versuche angestellt wurden, bemerkte ich, dass
die Eindrucke simultan in das Bewusstsein gelangten». --- \emph{Berger}: Ueber
den Einfluss der Uebung auf geistige Vorgänge (Wundt's Studien, V. 1888). p.
177.}. --- Det er den ene Ejendommelighed, som vel maa holdes fast, for at den
umiddelbare Genkendelse kan skelnes fra andre Arter af Genkendelse. Ved den
adskiller den umiddelbare Genkendelse sig paa den ene Side fra saadanne
Tilfælde, hvor aldeles usammensatte og i vor sædvanlige Erfaring ikke som
isolerede forekommende Fornemmelser skulle genkendes, paa den anden Side fra
saadanne Tilfælde, hvor Opfattelsen af Fænomenet, og som Følge deraf ogsaa
Genkendelsen maa ske ved en successiv psykologisk Proces.

Den anden Ejendommelighed er den, at Genkendelsen er \emph{uforberedt}, d. v.
s. der gaar ikke kort forud nogen vilkaarligt eller uvilkaarligt fremkaldt
Forestilling om det Fænomen, som er Genkendelsens Genstand. --- Ved nogle af de
anførte Tilfælde var der vel paa en Maade en Forberedelse, idet man havde søgt
at faa Forestillingen om Fænomenet frem, men da det ikke lykkedes, er denne
Forberedelse højst indirekte, skønt den derfor ingenlunde er betydningsløs%
% footnote 2) on p. 10
\footnote[2]{Smlgn. min Psykologi, p. 83 f. --- B. \emph{Erdmann} (Zur Theorie
der Apperception. Vierteljahrsschrift fur wissenschaftliche Philosophie. 1886.
p. 343) skelner derfor mellem «ein erregtes und ein unerregtes Unbewusste», af
hvilke det første danner Bevidsthedens umiddelbare Grundlag.}. Dette Spørgsmaal
om Forberedelse vil i det Hele vise sig at være meget mere indviklet end man i
Regelen har antaget. Der er her utallige Grader mellem den fuldstændige
Uforberedthed og den fuldstændige Forberedthed. De vigtigste Forskelligheder
ville vi i det Følgende faa Lejlighed til at fremdrage.

For det Tredie maa det endnu fremhæves, at Selviagttagelsen i de anførte
Tilfælde \emph{ikke viser ringeste Spor af andre Forestillinger}, \emph{som
vækkes} ved det \emph{genkendte Fænomen} og kunde antages at spille en Rolle
ved selve Genkendelsen. For saa vidt nogen altsaa vilde antage, at al
Genkendelse forudsætter saadanne Forestillinger, paahviler Bevisbyrden ham, og
kan den umiddelbare Genkendelse, saaledes som den fremtræder i de anførte
Tilfælde, forklares uden en saadan Antagelse, vil det være den eneste
videnskabelige Forklaring.

4. En Vej til den simpleste Forklaring vil vise sig, naar vi nærmere
sammenligne en ny Fornemmelse med en genkendt Fornemmelse. For Kortheds Skyld
bruger jeg i det nærmest Følgende Udtrykket Fornemmelse i Stedet for
Fornemmelseskomplex eller Fornemmelsesgruppe. Det er jo, som allerede
fremhævet, i Regelen ikke enkelte Fornemmelser, det her drejer sig om.

Ved en ny \emph{Fornemmelse} behøve vi her ikke at forstaa en \emph{absolut} ny
Fornemmelse, d. v. s. en saadan, som vi aldrig nogensinde før have havt.
Saadanne absolut
% --- p. 11 ---
\setcounter{footnote}{0}%
\opage{11}nye Fornemmelser kunne vi, naar vi holde os til de simpleste
Tilfælde, kun have ved Bevidsthedslivets første Begyndelse eller --- hvad
Almenfornemmelserne angaar --- ved Begyndelsen af nye Livsperioder
(Pubertetsaarene f. Ex.) eller ved nye Sygdommes Frembrud. En ny Fornemmelse
kunne vi ogsaa da siges at have, naar der er gaaet en vis, ikke for kort Tid,
siden vi sidst havde en lignende Fornemmelse. Naar jeg kommer hjem en Aften og
efter at have klædt mig af tænker over, hvor vidt jeg nu ogsaa har laaset min
Gangdør, vil jeg uden Ulejlighed kun kunne afgøre dette, dersom jeg kan
reproducere de Fornemmelser, der ledsage min Aflaasen af Døren, og \emph{finde}
dem friskere og tydeligere end saadanne reproducerede Fornemmelser, der svare
til tidligere Aflaasninger. Hvis jeg tror at have lukket Døren, uden i
Virkeligheden at have gjort det, er Grunden den, at Reproduktionerne af
tidligere levende Fornemmelser af denne Handling have frembragt hos mig Skinnet
af, at en Gentagelse af Handlingen fandt Sted i Aften. Der opstaar da en
Erindringsillusion. --- Men under sædvanlige Forhold er en Fornemmelse relativt
ny, naar den indtræder efter et vist Mellemrum. Hvor langt dette maa være, for
at Fornemmelsen vedblivende skal kunne kaldes ny, vil være forskelligt i de
forskellige Tilfælde og for de forskellige Sansemodaliteter. En Tone genkendtes
med absolut Sikkerhed efter to Sekunders Mellemrum; endnu ved fem Sekunders
Mellemrum mærkedes Tidens Indflydelse ikke tydeligt; derimod mærkedes ved 10 og
15 Sekunders Mellemrum en betydelig Tiltagen af urigtige Skøn, og ved 30
Sekunders Mellemrum var Sikkerheden ganske forsvunden%
% footnote 1) on p. 11
\footnote[1]{\emph{Wolfe}: Untersuchungen über das Tongedächtniss (Wundts
Studien III. Leipzig 1886). p. 522. --- Smlgn. de i min Psykologi p. 141
anførte Fakta, som vise, at Tidsmellemrummet ogsaa kan bevirke, at
Forskellighederne ikke opfattes, hvilket bekræftes ved de af Wolfe ib. p. 541
f. meddelte Resultater.}. Ved mere sammensatte Tilfælde vil Tidens Indflydelse
vistnok ikke saa snart gøre sig gældende, da der er flere Holdepunkter for
Opmærksomheden, idet nemlig Forholdet mellem Elementerne i Regelen spiller en
større Rolle end noget af de enkelte Elementer for sig alene.

Sammenligne vi nu en (relativt) ny \emph{Fornemmelse} med en gentagen og
genkendt, saa er der én Betingelse, de i hvert Tilfælde have fælles, den
nemlig, at et Indtryk er kommet til Hjernen (Bevidstheden), og at dette Indtryk
staar i et saadant Forhold til Hjernens (Bevidsthedens) Tilstand, at en vis
bestemt Fornemmelse (med en vis Kvalitet og Selvstændighed) kan fremtræde. Hvad
kan nu det være, som \emph{Genkendelsen} forudsætter ud herover?

Ved Besvarelsen af dette Spørgsmaal maa vi først og fremmest holde os til den
eneste Omstændighed, som er forskellig ved de to Tilfælde: i det ene er
Indtrykket (relativt) nyt, i det andet gentaget. Den eneste Virkning, som denne
Omstændighed kan have, er, at en Reproduktion bliver mulig. Men denne
Reproduktion behøver ikke at føre til, at det, som reproduceres, fremtræder som
selvstændigt Led i Bevidstheden, og i de foreliggende Tilfælde sker dette
heller ikke. Deres Ejendommelighed bestod bl. a. netop i deres tilsyne%
% --- p. 12 ---
\setcounter{footnote}{0}%
\opage{12}ladende usammensatte Karakter. Foruden det eller de genkendte Træk
findes i Bevidstheden ikke det mindste, som har med Genkendelsen at gøre. Ordet
«Les Plans» lyder bekendt, og denne Bekendthedskvalitet ved Klangen er det
\emph{hele Fænomen}. Uagtet jeg strax bliver opmærksom paa Fænomenets
psykologiske Interesse, kan jeg dog ikke opdage det allermindste Spor af
medvirkende Fornemmelser eller Forestillinger. Naar det nu alligevel staar
fast, at Bekendthedskvaliteten maa hænge sammen med, at jeg har hørt Ordet før,
kan denne Kvalitet kun forklares ved, at den tidligere Tilstand, i hvilken jeg
kom, da jeg første Gang hørte Ordet, paa en eller anden Maade kan gøre sig
gældende paa ny sammen med selve den gentagne Fornemmelse, uden dog at træde
frem som selvstændigt Element i Bevidstheden ved Siden af denne. Dette kunne vi
meget vel udtrykke ved Ordet Reproduktion, naar vi antage, at det nu kommende
Indtryk i Hjernen (Bevidstheden) vækker en Tilstand, som \emph{tillige}
bestemmes ved Eftervirkningen af det tidligere Indtryk. Og denne Eftervirkning,
kunne vi simpelt hen antage, bestaar i den større Lethed, med hvilken
Tilstanden ved Gentagelse indtræder. Vi appellere her til den Øvelsens Lov, som
gælder for alt organisk Væv, ikke mindst for Nerve- og Muskelvævet. Jo
hyppigere en Proces er gaaet for sig i et Organ, des lettere indtræder den paa
ny. Selv udenfor det organiske Omraade vil man kunne tale om Øvelse, idet jo
som bekendt musikalske Instrumenters Sangbund des lettere sættes i Svingninger,
jo mere de bruges. Ethvert nyt Forsøg, ethvert nyt Indtryk nyder godt af de
tidligere Indtryk, idet disse ligesom have banet Vejen for det. ---
\emph{Psykologisk} kunne vi udtrykke, hvad der sker ved Genkendelsen, ved
Formlen (\textit{A} + \textit{a}), eller maaske bedre $\left({a\atop
A}\right)$, hvor \textit{A} betyder Fornemmelsen, \textit{a} dens Reproduktion
(Forestillingen), og hvor der ved Parentesen betegnes, at disse to her ikke
træde frem som selvstændige Led i Bevidstheden, men \emph{tænkes teoretisk} som
Faktorer i det tilsyneladende usammensatte Fænomen. Berettigelsen til denne
Udtryksmaade ligger deri, at den Eftervirkning af \textit{A}, som gør
Genkendelse mulig, under andre Forhold vilde have gjort \textit{a'}s Fremtræden
som fri (selvstændig) Forestilling mulig, f. Ex. hvis \textit{B} ved
Berøringsassociation fremkaldte \textit{a.} Meningen er altsaa: det
\emph{samme}, som er en \emph{Betingelse} for, at \textit{B under visse Forhold
kan genkalde} \textit{a,} er \emph{ogsaa en Betingelse} for, at \textit{A} kan
\emph{genkendes}. Navnet «Les Plans» kunde maaske være blevet fremkaldt ved
Synet af eller Forestillingen om Stedet; men Navnet kan, naar jeg hører det,
forekomme mig bekendt, selv om det ikke kan fremkalde Forestillingen om Stedet.
I Bekendthedskvaliteten virker da i sidste Tilfælde det samme, som i første
Tilfælde vilde være traadt frem som en selvstændig Forestilling. --- Hvorledes
man \emph{fysiologisk} vil tænke sig det, der ved Øvelsen foregaar i
Organismens, her Hjernens Smaadele, bliver en Sag for sig, som vi her ikke
behøve at gaa nærmere ind paa. Den naturligste Antagelse vilde vel være den, at
der ved det første Indtryk bevirkes en Omlejring af Molekulerne, som efter
Indtrykkets Ophør igen afløses af den tidligere Tilstand, men saaledes, at denne
% --- p. 13 ---
\setcounter{footnote}{0}%
\opage{13}nu er mere usikker, \emph{lettere at bringe} ud af \emph{Ligevægt}.
Der kan for saa vidt siges at være tilvejebragt en vis Disposition til samme
Omlejring, saa at denne lettere gaar for sig, naar samme Indtryk kommer igen%
% footnote 1) on p. 13
\footnote[1]{Denne Antagelse er saa langt fra at være noget moderne, at den
allerede med stor Klarhed udtales af \emph{Charles Bonnet} i hans Essai
analytique sur l'âme (Copenhague 1760), p. 61, hvor det hedder: «La souplesse
ou la mobilité des fibres augmente par le retour des mèmes ébranlements. Le
\emph{sentiment attaché} à \emph{cette augmentation de souplesse} ou de
\emph{mobilité}, \emph{constitue la réminiscence}.»}. Genkendelsen eller
rettere Bekendthedskvaliteten danner da det psykologiske Korrelat til den
større Lethed, med hvilken en Ændring i vedkommende Hjernemolekulers Lejring
frembringes. Der frembyder sig ingen anden Omstændighed, ved hvilken den
ejendommelige Forskel mellem Fornemmelse og Genkendelse kan forklares. Nu er
Forskellen mellem større og mindre Omlejringsmulighed ganske vist kun en
Gradsforskjel, medens der psykologisk set er en Kvalitetsforskel mellem
Fornemmelse og Genkendelse. Men dette kan ikke undre os, naar vi betænke, at
til de psykologiske Kvaliteter svare der i den fysiske Verden stedse kun
Kvantitetsforskelle. Farvekvaliteterne have deres fysiske Korrelater i
kvantitative Forskelle mellem Ætersvingningerne, og sandsynligvis ogsaa deres
fysiologiske Korrelater i kvantitative Forskelle mellem Molekulernes Bevægelser
i Hjernens Synscentra. Vor Forklaring af Genkendelsesfænomenet i dets simpleste
Form er altsaa i god Analogi med hvad der synes at gælde for andre simple
psykiske Fænomener.

Jeg bemærker blot endnu, at den her fremstillede Teori er en videre Udvikling
af hvad jeg allerede har fremstillet i min Psykologi (2. Udg. p. 140---142).
Jeg har her søgt at give nærmere Bestemmelser og Forklaringer, som have vist
sig at være nødvendige.

5. Fornemmelse og Genkendelse kunne ofte rykke hinanden saa nær, at der ikke
kan gøres nogen afgørende Forskel mellem dem. Det hænger sammen med, at vi
neppe have nogen Fornemmelse, uden at en vis Opmærksomhed vendes imod den. Og
dette kommer igen af, at der med enhver Fornemmelse følger en Følelse af Lyst
eller Ulyst og derfor en Trang til Virksomhed og Bevægelse, som naturligt
rettes mod det, der fremkalder Fornemmelsen. Men den saaledes vakte
Opmærksomhed virker ikke aldeles kontinuerligt. Den slipper sin Genstand og
griber den kort efter igen, naar Indtrykket, og dermed Opmærksomhedens Motiv,
vedbliver at virke. Hvori man end vil søge Aarsagen dertil, saa er den
Opmærksomhed, vi rette mod et kontinuerligt Indtryk, diskontinuerlig og
periodisk. Tydeligst ses dette ved meget svage Sanseindtryk. Naar et Tryk paa
Enden af Fingeren kontinuerligt tiltager, mærke vi ikke denne kontinuerlige
Tiltagen; Opmærksomheden sammenholder forskellige Stadier af Indtrykkets
Stigning og derved blive vi os denne bevidst%
% footnote 2) on p. 13
\footnote[2]{\emph{Stanley Hall}, anført hos \emph{Ribot}: La Psychologie de
l'Attention. Paris 1889. p. 16.}. Ved meget smaa Sansefornemmelser, der ligge
lige over Bevidsthedens
% --- p. 14 ---
\setcounter{footnote}{0}%
\opage{14}Tærskel, viser Erfaring, at de i det ene Øjeblik fremtræde i
Bevidstheden, men i næste Øjeblik forsvinde; de ere i en stadig, rytmisk
Bevægelse, som, naar Indtrykket virker kontinuerligt, kun kan forklares ved
Opmærksomhedens Periodicitet. Naar et og samme Billede kan opfattes
forskelligt, vexle de forskellige Opfattelser gerne med hinanden efter visse
Mellemrum. Et Omrids af en Trappe (den saakaldte Schröderske Trappefigur) kan
man ogsaa «se» som en Mur, af hvilken Stene for neden ere faldne ud. Betragter
man opmærksomt og uafbrudt dette Billede, skifte Trappe og Mur periodisk med
hinanden. Det er ikke Billedet, vel heller ikke egentligt Fornemmelsen, men det
maa være Opmærksomheden, der skifter%
% footnote 1) on p. 14
\footnote[1]{N. \emph{Lange}: Beiträge zur Theorie der sinnlichen
Aufmerksamkeit und der activen Apperception (Wundt's Studien. IV. Leipzig
1888). p. 395. 406.}. Det synes, som om den fulde Kraftsum, Opmærksomhedsakten
behøver, ikke i hvert Øjeblik kan være til Stede, hvorfor der med visse
Mellemrum maa anvendes nogen Tid paa at erstatte den forbrugte Energi. Men idet
saaledes Opmærksomheden griber sin Genstand igen, efter at have sluppet den,
kommer stedse Reproduktionen af det foregaaende Tag i Genstanden til at spille
en Rolle med. Naar en opmærksom Betragten gør os en simpel Genstand klarere,
kan det kun ske ved, at saadan Reproduktion (opfattet paa den i 4 fremstillede
Maade) forstærker Fornemmelsen. Gennem denne Vexelvirkning af Fornemmelse og
Reproduktion, altsaa gennem en Række Genkendelsesakter, forløber den
tilsyneladende kontinuerlige Betragtning. Tydeligt ses dette, naar vi
\emph{vilkaarligt} søge at fastholde en bestemt Tydning af den omtalte
Trappefigur; vi søge alvorligt at sætte os ind i, at det skal være en Trappe,
hvilket sker ved vilkaarlig Reproduktion af tidligere sete Trapper%
% footnote 2) on p. 14
\footnote[2]{Smlgn. N. \emph{Lange} ib. p. 407. Se ogsaa Wundt: Physiologische
Psychologie. 3. Aufl. Leipzig 1887. II. p. 238. 253 f.}, som dog ikke komme til
at træde frem som selvstændige Bevidsthedselementer, men kun tjene til
Forstærkelse af Fornemmelsen. Ved \emph{uvilkaarlig} Opmærksomhed maa det
naturligvis være selve Indtrykket, der, saasnart den indre Betingelse,
tilstrækkelig Energi, er til Stede, stedse paa ny udløser Opmærksomheden og
Reproduktionen. --- Paa denne Maade væves der altsaa en større eller mindre
Række af Genkendelsesakter ind i vore tilsyneladende ganske simple og
usammensatte Erfaringer.

6. En umiddelbar Genkendelse kan lige saa vel finde Sted overfor en
Forestilling (eller en Forestillingsgruppe) som overfor en Fornemmelse (eller
en Fornemmelsesgruppe). Vi kunne genkende Erindringsbilleder og
Fantasibilleder. Naar i en vaagen Drøm et Billede uvilkaarligt skyder sig frem
for mig, kan det komme med aldeles lignende Bekendthedskvalitet som den, et
Landskab har, naar jeg ser det paa ny. Erindringsbilledet nyder da godt af de
tidligere Indtryk. Det er en ikke sjelden Erfaring, at vi ikke gen%
% --- p. 15 ---
\setcounter{footnote}{0}%
\opage{15}kende Erindringsbilleder, at de staa for os som fremmede, og at vi
først møjsommeligt maa overbevise os om, at de svare til virkelige Oplevelser.
I saadanne Tilfælde mangle de hin Kvalitet. Og selv om vi ere sikre paa, at vi
have oplevet det forestillede før, kunne vi være usikre om, \emph{hvor og
naar}. Ved Fantasibilleder er det ikke tidligere virkelige Oplevelser, men
forestillede Oplevelser, det drejer sig om. Jeg kan f. Ex. have vænnet mig til
at forestille mig en digterisk Figur paa en vis bestemt Maade, og jeg kender da
denne Skikkelse, naar min Fantasi uvilkaarligt fører den frem for mig. Det er
et Fænomen, som særligt vil være hyppigt hos instinktmæssigt og med levende
Fantasi producerende Digtere. Smlgn. \emph{Goethe's} berømte Linier:

Ihr naht euch wieder, schwankende Gestalten!

Die früh sich einst dem trüben Blick gezeigt.

Her virker samme Lov, som ved Genkendelse af noget i Virkeligheden Erfaret,
nemlig Øvelsens Lov. Den gælder lige saa vel for de Hjernefunktioner, til
hvilke Erindring og Fantasi er knyttet, som for den, til hvilken
Sansefornemmelse er knyttet.

7. Skønt den umiddelbare Genkendelse for Selviagttagelsen fremtræder som et
aldeles simpelt Fænomen, kunne vi dog altsaa \emph{teoretisk} betragte den som
sammensat. Der er et Moment af Erindring, af Eftervirkning fra Fortiden, som i
den umiddelbart forbinder sig med den i Øjeblikket vakte Fornemmelse. Hvad der
fremtræder i Bevidstheden er begge Momenters Produkt; deres Sammensmeltning er
foregaaet under Bevidsthedens Tærskel. Jeg har derfor i min Psykologi (V B, 1)
betegnet Fænomenet som \emph{bunden Erindring}, i Modsætning til saadanne
Erindringsfænomener, ved hvilke det, som reproduceres, fremtræder som
selvstændigt Element i Bevidstheden; der finder her en fri \emph{Erindring} Sted.

Hvis det er berettiget at kalde den umiddelbare Genkendelse et
Erindringsfænomen, maa det ogsaa være rigtigt at [kalde den et
Associationsfænomen. Det staar naturligvis Enhver frit for, hvorledes han vil
bruge Ordene; og hvis man ved Association kun vil forstaa en saadan Forbindelse
af Forestillinger, ved hvilken begge de forbundne Led gøre sig gældende paa fri
og selvstændig Maade i Bevidstheden, kan den umiddelbare Genkendelse ikke
kaldes Association. Saaledes beskriver \emph{David Hartley},
Associationsteoriens anden Grundlægger, tydeligt baade den umiddelbare
Genkendelse og den frie Lighedsassociation, skønt han bruger Ordet Association
blot om Berøringsassociationen%
% footnote 1) on p. 15
\footnote[1]{Observations of man. Londonerudgaven af 1791. I. p. 291 ff.}. Og
--- for at tage et Exempel fra den nyeste Tid --- \emph{James Ward}, der
forkaster den frie Lighedsassociation, anerkender dog den beskrevne umiddelbare
Samvirken af Fornemmelse
% --- p. 16 ---
\setcounter{footnote}{0}%
\opage{16}og Reproduktion, men erklærer, at det er Assimilation og ikke
Association. Skal det kaldes Association, er det automatisk Association%
% footnote 1) on p. 16
\footnote[1]{Art. «Psychology» i 9. Udgave af Encyclopædia Britannica. Vol. XX.
p. 60 b.}. --- Den umiddelbare Genkendelse (eller, som jeg ogsaa kalder den:
Perceptionen) er et Grænsefænomen, som har den tilsyneladende usammensatte
Karakter fælles med Fornemmelsen, og som har fælles med den frie Association,
at Reproduktion tydeligt nok spiller en Rolle med. Fra den frie Association
adskiller den sig ved den fuldstændige Sammensmeltning af Fornemmelse og
Reproduktion. Der er i den umiddelbare Genkendelse forbundet Elementer, som vi
ikke blot kunne \emph{tænke} os særskilte, men som ogsaa under andre
Betingelser kunne optræde hver for sig. (Smlgn. 4.) Genkendelsen indtræder ikke
altid i første Øjeblik. \textit{A} kan komme op som blot Fornemmelse og først
efter en lille Mellemtid blive til $\left({a\atop A}\right)$; og paa den anden
Side kan \textit{a} fremkaldes i Bevidstheden ved fri Association, altsaa som
selvstændigt Element. Den umiddelbare Genkendelse forholder sig til
Associationen som Tangenten til Sekanten. Ved Tangenten ere de to
Skæringspunkter med Cirklen smeltede sammen til ét; den er en Sekant, hvis
Skæringspunkter ligge uendeligt nær ved hinanden. Det maatte nærmest være
Lighedsassociationen, med hvilken man efter denne Analogi maatte sammenstille
den umiddelbare Genkendelse. --- Det Spørgsmaal, om der gives en virkelig
Lighedsassociation, ville vi senere gaa nærmere ind paa.

8. Lige som den umiddelbare Genkendelse i det Hele er et Exempel paa Øvelsens
psykologiske Indflydelse, saaledes vil den ogsaa ved fortsat Gentagelse
undergaa nye Forandringer. Den foregaar ved Gentagelse hurtigere og hurtigere,
og kan foregaa saa hurtigt, at den ikke kræver nogen særskilt Bevidsthedsakt,
men blot tjener til at udløse følgende Bevidsthedsakter. Den nærmer sig da
Bevidsthedens Tærskel, om den ikke overskrider den. Forsøg have vist (hvad
iøvrigt almindelig Erfaring allerede godtgør), at med Kendskabet til et Sprog
stiger den Hurtighed, med hvilken man paa dette Sprog kan lære Bogstaver, som
danne Ord, og Ord, som danne Sætninger. Den bekendte Erfaring, at det synes os,
som om Udlændinge tale hurtigere end vi selv, forklares herved. Tage vi derimod
Ord, som ikke danne Sætninger, og Bogstaver, som ikke danne Ord, saa fordobles
omtrent den Tid, der udfordres til Sætningen%
% footnote 2) on p. 16
\footnote[2]{\emph{Cattell}: Ueber die Zeit der Erkennung und Benennung von
Schriftzeichen, Bildern und Farben (Wundt's Studien. II. Leipzig 1885). p. 646
ff.}. Naar den fysiologiske Tid, den Tid, der medgaar til Modtagelse af et
Indtryk og til Besvarelse af det ved en Bevægelse, forkortes formedelst
Gentagelse og Øvelse, kan den hele Proces sandsynligvis nu forløbe gennem
underordnede Hjernecentra, uden at de højere Centra, til hvilke de kombinerende
Bevidsthedsprocesser ere knyttede, behøve at virke med. Indtrykket udløser da
umiddelbart Be%
% --- p. 17 ---
\setcounter{footnote}{0}%
\opage{17}vægelsen, uden at man er sig bevidst, at man har modtaget Indtrykket.
Genkendelsen fortoner sig saaledes over i det Ubevidste i Kraft af de samme
Aarsager, ved hvilke den fra først af blev til. Men den har været en nødvendig
Betingelse for, at vedkommende Processer nu kunne foregaa med saa stor
Sikkerhed og Hurtighed. Den er et første Stadium af Tillempelse og betinger de
følgende Tillempelsesstadiers Indtræden. Vi gøre i daglig Tale ikke altid
Forskel mellem disse forskellige Stadier. Vi sige f. Ex. om den øvede Musiker,
at han \emph{kender} Noderne, skønt han slet ikke behøver at anvende nogen
særskilt Bevidsthedsakt ved Synet af dem, da de umiddelbart kunne udløse
Berøringsassociationer eller Bevægelser. Sprogbrugen er i sin Ret; ti denne
\emph{sekundære} Ubevidsthed, der skyldes Øvelse, maa ikke uden videre stilles
sammen med den \emph{primære} Ubevidsthed, som er til Stede, naar ingen Indtryk
og Fornemmelser ere opstaaede. Og dertil kommer, at Grænsen mellem bevidst og
ubevidst «Kending» kan være vanskelig at drage. Naar en Genkendelsesakt staar
som Indleder af en Række betydningsfulde, affektvækkende Forestillinger, overse
vi den let, idet hele Interessen er vendt mod det, som følger efter. Hvad der
kun har et Minimum af Interesse forsvinder (især for den uøvede
Selviagttagelse) i Forhold til det, der sætter Interessen i stærkere Bevægelse,
og bevares ikke i Erindringen. Over Explosionen glemme vi Gnisten, der voldte
den. Dette hænger sammen med, at det aldeles overvejende er praktiske Hensyn,
som bestemme vore Forestillingers og vor Erkendelses Udvikling. Vi have ingen
praktisk Interesse af at lægge særligt Mærke til de Bevægelsesfornemmelser, der
spille saa stor en Rolle ved Opfattelse af Afstand, hvorimod det har en stor
Interesse at være sikker i Bedømmelsen af Afstanden mellem os og Tingene uden om os.

Naar Instinkter udløses ved blotte Indtryk og Fornemmelser, er det sikkert for
en stor Del sekundære, automatisk blevne Genkendelser, vi her have for os. Den
unge Slangegrib, som strax slaar ned paa en Slange (saa ivrigt, at man kan faa
den til at slaa ned paa en Tarm i Stedet), lægger vel ikke meget Mærke til,
hvad der fremkalder dens Anfald. Men Indtrykket af Slangeskikkelsen maa dog
ligesom sætte gamle Strenge i Bevægelse, maa kunne forplante sig ind i dens
Indre ad forud banede Veje, for at det skal have denne Virkning. Naar en Rotte
farer sammen ved et Katteskrig, selv om den har mistet den store Hjerne, saa er
det, fordi den har tillempet sig saaledes efter dette Indtryk, at det kan gøre
sin Virkning, selv om de højeste Centralorganer mangle. Ud fra disse højeste
Organer ere Bevægelsescentrene blevne saaledes indøvede, at de kunne sættes i
Funktion ved et Indtryk eller en Fornemmelse, som ellers ikke vilde være
tilstrækkelig hertil. Tillempelsen kommer her ikke blot de enkelte Individer,
men ogsaa Efterkommerne til Gode, hvilket f. Ex. kan ses af den Maade, paa
hvilken nyfødte Dyrs og Børns Næringsinstinkt ytrer sig.

% --- p. 18 ---
\setcounter{footnote}{0}%
\opage{18}9. Det maa vel mærkes, at der her kun er talt om \emph{umiddelbar}
Genkendelse, om Genkendelsesfænomenet i dets simpleste Form. Hvis den
Forklaring af dette Fænomen, som jeg i det Foregaaende har givet, er rigtig,
har det den Interesse, at det frembyder en Overgangsform mellem de blotte
Fornemmelser og de egentlige (frie) Forestillinger. Jeg har derfor i min
Psykologi (V B, 1) fremstillet det efter Afslutningen af Læren om
Forestillingerne. --- Men denne Interesse vilde være grundet paa Illusion,
dersom de Indvendinger vare berettigede, der i den nyeste Tid ere rettede mod
den fremstillede Opfattelse.

Genkendelse, har man ment, er aldrig noget aldeles simpelt Fænomen, men maa
stedse forklares efter den sædvanlige (frie) Forestillingsassociation. Og idet
man saa tillige har negtet, at der gives egentlig Lighedsassociation, maa det
være Berøringsassociationens Lov, der lægges til Grund. Den Proces, der skal
finde Sted ved al Genkendelse, maa da være følgende. Jeg nærmer mig til et Hus,
som jeg har set før, og fæster mit Blik paa en enkelt Del af det, f. Ex. et
Vindu. Dette Vindu vækker nu (ved Berøringsassociation) Forestillingen hos mig
om en tilgrænsende Del af Huset, f. Ex. et Ornament. Kalder jeg Synet af
Vinduet \textit{A,} Forestillingen om Ornamentet \textit{b} (idet store
Bogstaver stadigt bruges til at betegne Fornemmelser eller Fornemmelsesgrupper,
og smaa Bogstaver til at betegne Forestillinger eller Forestillingsgrupper),
saa har jeg altsaa nu \textit{A} + \textit{b.} Med Forestillingen om Ornamentet
opstaar tillige naturligt en Forventning om at faa det at se, og idet denne
Forventning opfyldes, naar Blikket forlader Vinduet og gaar videre hen ad
Huset, indtræder Fornemmelsen \textit{B.} Det er da nu, efter denne Teori,
Overensstemmelsen mellem den \emph{forud} reproducerede Forestilling \textit{b}
og den \emph{dernæst} indtrædende Fornemmelse \textit{B}, som betinger
Genkendelsen, eller hvori denne bestaar. Og naar denne Proces gentager sig ved
Husets forskellige Dele, sige vi tilsidst, at vi genkende Huset.

Overfor denne Teori vil jeg aller først bemærke, at den psykologiske Proces,
den skildrer, er aldeles naturlig og i vort virkelige Bevidsthedsliv stadigt
finder Anvendelse. Jeg har selv fremdraget den i min Psykologi (V B, 2 og 4);
men det er rigtignok ikke faldet mig ind, at den kunde bruges til Forklaring af
al Genkendelse og til at bestride, at der gaves \emph{umiddelbar} Genkendelse i
den anførte Betydning af Ordet. Af den Omstændighed, at der gives
\emph{successiv} Genkendelse, som gaar for sig paa den beskrevne Maade, kan jo
dog ikke uden videre udledes, at der \emph{ikke} kan gives \emph{umiddelbar}
Genkendelse.

Det er nemlig klart, at den anførte Teori ikke passer til de tidligere anførte
simple Genkendelsesfænomener, saadanne Tilfælde, i hvilke et eller flere Træk
fremtræde for os med Bekendthedens Præg, uden at der i Bevidstheden kan spores
nogen som helst anden dertil knyttet Forestilling. Disse simpleste Tilfælde af
Genkendelse kunne meget godt angaa \emph{forholdsvis} sammensatte Genstande;
Genkendelse af sammensatte Genstande er nemlig ikke det samme som successiv
Genkendelse, naar det blot er muligt at sammenfatte det givne Indhold ved én
Bevidsthedsakt. (Smlgn. 3.) Hovedsagen er, at Bekendtheds%
% --- p. 19 ---
\setcounter{footnote}{0}%
\opage{19}kvaliteten umiddelbart gør sig gældende og staar isoleret i vor
Bevidsthed, uden at vække nogen Forestilling. Der er da her ikke Tale om nogen
successiv Bevægelse af Iagttagelsen fra Del til Del. Det vilde være en aldeles
uberettiget Konstrueren at lægge en saadan successiv Proces ind i hine simple
Fænomener.

Men selv hvor en saadan Proces gaar for sig, nemlig overfor saadanne Genstande,
der ikke kunne sammenfattes ved én Bevidsthedsakt, fører den anførte Teori til
Strid mod Erfaringen. Den forudsætter nemlig, at af en sammensat Gjenstand
\textit{X} = \textit{A} + \textit{B + C} (hvor Bogstaverne betyde Genstandens
Dele eller Egenskaber) genkender man \emph{aldrig} den første Del, som
opfattes, men \emph{altid} først den anden eller tredie. \textit{A} kan først
genkendes, naar man senere fra en af de andre Dele vender tilbage til
\textit{A.} Viser Erfaring dog ikke, at man meget vel kan genkende allerede det
første Led?

10. Den anførte Paastand, at al Genkendelse beror paa, at en forud vakt
Forestilling efterfølges af den tilsvarende Fornemmelse, grunder sig paa en
\emph{Forvexling af selve den umiddelbare Genkendelse og denne Genkendelses
Verifikation}. Den umiddelbare Genkendelse, der sker i et Nu og er en ganske
simpel aandelig Proces, kan i og for sig ikke godtgøre sin egen Gyldighed.
Baade Sanseillusioner og Erindringsillusioner godtgøre tilstrækkeligt, at
Bekendthedskvaliteten kan være til Stede, uden at der foreligger en virkelig
Gentagelse af det Fænomen, vi tro at have for os. Skal Genkendelsens Rigtighed
prøves, kan det kun ske ved at \emph{undersøge de Forestillinger}, \emph{det
genkendte Fænomen vækker} hos mig. \emph{Dette er umuligt} i \emph{saadanne
Tilfælde} som de i 3 \emph{omtalte}, i \emph{hvilke noget stod som bekendt} og
dog \emph{ingen Associationer voldte}. \emph{Netop saadanne Tilfælde ere derfor
af Interesse som experimenta crucis til Prøvelse} af de \emph{forskellige
Teorier}. Men for det meste vil det Genkendte strax ved Berøringsassociation
vække Forestillinger om Træk, der hænge sammen med det, og man kan da prøve, om
disse medfølgende Træk ogsaa virkeligt ere til Stede. (Smlgn. min Psykologi V
B, 4 om den experimenterende Metode, ved hvilken Bevidstheden uvilkaarligt
underkaster de opdukkende Forestillinger en Ildprøve og derved lærer at drage
Grænsen mellem det Mulige og det Virkelige. Bevidstheden gaar her praktisk frem
paa samme Maade som Videnskaben: den opstiller en foreløbig Hypotese og søger
at verificere den.) Ved Opfattelsen af meget sammensatte Genstande kan
Verifikationen saaledes følge med selve Opfattelsen Skridt for Skridt. Men det
\emph{behøver} ikke at gaa saaledes. Det ses allerede deraf, at vi kunne have
opfattet en Række Fænomener, uden at det er os muligt frit at reproducere de
enkelte Led i Rækken, skønt vi \emph{kende} dem igen, saa snart de paa ny
frembyde sig. (Smlgn. 3.) Naar jeg læser en god Bog anden Gang, beror den
særegne Nydelse (som kan kaldes Genkendelsens Nydelse) sikkert ikke derpaa, at
jeg uafbrudt tænker paa, hvad der skal komme, og stadigt faar denne, ikke
synderligt hensigts%
% --- p. 20 ---
\setcounter{footnote}{0}%
\opage{20}mæssige%
% footnote 1) on p. 20
\footnote[1]{\emph{David Hume} vilde i hvert Tilfælde have sig saadanne
Anticipationer frabedt, og der er vist mange, som dele hans Smag. Smlgn. hans
«Essay on Simplicity and Refinement» (Essays. London 1777. Vol. I, p. 211): If
the merit of the composition lies in a point of wit, it may strike at first,
but the mind anticipates the thought in the \emph{second perusal}, and is no
longer affected by it. When I read an epigram of Martial, the first line
recalls the whole, and I \emph{have no pleasure in repeating in myself what} I
\emph{know already}. But each line, each word in Catullus has its merit; and I
am never tired with the perusal of him.» --- Det er aabenbart, at Hume ikke
hyldede Forventningsteorien. Hvis den gjaldt ubetinget, vilde Hume ikke have
læst nogen Bog anden Gang.} Forventning bekræftet. Men den beror paa, at hvert
enkelt Træk, som frembyder sig, lettere glider ind i min Bevidsthed eller
opstaar i denne med Bekendthedens Kvalitet. Vi dvæle ved de enkelte Træk; de
staa umiddelbart for os som gamle Bekendte. Der var mere Forventning første
Gang, vi læste Bogen; da hastede vi fremad, fordi vi vidste, at der kom noget
efter; nu er denne Nysgerrighed tilfredsstillet, og vi kunne trygt dvæle paa
hvert Sted, som om vi hørte hjemme der.

At Verifikationen, hvor den er mulig, oftest gaar uvilkaarligt for sig, kommer
af, at der, naar Bevidstheden eller Hjernen er i sund og kraftig Tilstand, er
en Trang til at anvende Energien til Forestillingsvirksomhed. Der gives en
(relativt) spontan Forestillingsvirksomhed, lige som der gives en (relativt)
spontan Bevægelse. Enhver opdukkende Fornemmelse eller Forestilling leder
Forestillingstrangen i en bestemt Retning efter Associationens Love. Enhver
Forestilling vil derfor, naar Bevidsthedens Energi ikke er svækket, sætte
Rækker af Forestillinger i Bevægelse. Derfor gaar umiddelbar Genkendelse ogsaa
let over til successiv Genkendelse gennem vakte og verificerede Forestillinger
og Genkendelser.

Der kan her undertiden opstaa Illusioner, som ikke kunne forklares efter
Forventningsteorien, men tvertimod synes bestemt at stride mod denne. Jeg kan
tro at kende et Ansigt, skønt det kun er et enkelt Træk, f. Ex. Øjet, der
ligner et mig bekendt Menneskes. Dette forklares let og simpelt, naar vi
antage, at den umiddelbare Genkendelse af det enkelte Træk har ladet de andre
Træk træde i Skygge for os. Men dersom Øjet havde vakt Forestillinger hos mig
om Pande, Mund o. s. v. hos det Menneske, jeg virkeligt kender, saa vilde jo
disse Forestillinger strax komme i Strid med Billederne af den Pande og Mund,
det nærværende Menneske har; Forventningen vilde blive skuffet; --- og
hvorledes var da Genkendelsesillusionen mulig? Begyndte jeg Opfattelsen af
Ansigtet med Øjet (\textit{A}), saa vilde dette vække f. Ex. Forestillinger om
Panden (\textit{b}) --- men naar der saa som næste Led i den successive
Opfattelse kommer --- ikke \textit{B}, men --- \textit{P}, saa synes den
indbildte Genkendelse efter den nævnte Teori slet ikke at kunne komme i Stand. ---

Verifikationen gaar ikke altid uvilkaarligt for sig; den maa undertiden, hvis
den i det Hele skal ske, foretages efter bestemt Forsæt og med vilkaarlig
Opmærksomhed. Vi sige da: «Dette er Vinduet \textit{A}; skal dette Vindu nu
være det, jeg har set paa Huset \textit{X}, maa Ornamentet \textit{B} findes
ved Siden af». Forventningen formuleres her som en Hypotese,
% --- p. 21 ---
\setcounter{footnote}{0}%
\opage{21}hvis Gyldighed hviler paa visse Betingelser, og Genkendelsen (hvis
den indtræder) formuleres som en Dom, ved hvilken Overensstemmelsen mellem
\textit{A} og\textit{ a}, Billedet af det nærværende Vindu og Forestillingen om
det tidligere sete Vindu, slaas fast. Ved denne Dom fremtræder det som sondret,
der i en umiddelbar Genkendelse vilde være umiddelbart sammensmeltet. Vi faa
\textit{A = a}, i Stedet for $\left({a\atop A}\right)$. I Genkendelsesdommen
sættes Fornemmelse og Forestilling i bestemt og bevidst Forhold til hinanden.
Selve den umiddelbare Genkendelsesakt vilde det ikke være rigtigt at kalde en
Dom. Man vilde derved (ligesom naar man vilde kalde Kontrastvirkninger og
Sanseillusioner for Slutninger) overse Forskellen mellem et reflekteret og et
aldeles umiddelbart Bevidsthedsfænomen.

I mange Tilfælde kan en Slutning drages af den logiske Sammenhæng, i hvilken et
Fænomen forekommer, uden at der synes at finde nogen ved Berøringsassociation
bevirket Genkendelse Sted. --- Jeg læste nyligt en lille Drengs tyske Pensum
med ham. Nogle Timer efter spurgte jeg ham: «Hvad betyder «Baarschaft»?» Dette
Ord var forekommet i hans Pensum, men det var ham nu aldeles fremmed, og han
kunde slet ikke huske, at han havde hørt det før. Men da jeg nævnte Sætningen:
«Er hatte alle seine Waaren in Baarschaft angelegt», vidste han strax, hvad
Ordet betød; men det var, fordi han \emph{sluttede} sig til dets Betydning af
Sammenhængen; Ordet klang ikke derfor mere \emph{bekendt} i hans Øren; dertil
behøvedes der endnu adskillige Gentagelser. --- Vi ville i det følgende faa
Brug for denne Distinktion mellem at kende et Ords Betydning (at kunne nævne et
andet Ord, der udtrykker det samme) og at genkende det i umiddelbar og egentlig
Betydning (saa at det fremtræder med Bekendthedskvaliteten som «Les Plans» i 3).

11. Vi kunne ikke undvære teoretiske Konstruktioner i Psykologien. En saadan
har jeg anvendt, naar jeg fremstiller den umiddelbare Genkendelse som en
Sammensmeltning af en Fornemmelse og en Forestilling: $\left({a\atop A}\right)$
eller (\textit{A} + \textit{a}). Men Forventningsteorien maa gøre Brug af langt
videre gaaende Konstruktioner, idet den i \emph{enhver} Genkendelsesakt, selv
den simpleste, maa skelne mellem fire forskellige Operationer: 1) \textit{A}'s
Fremtræden; 2) \textit{b}'s Fremtræden ved Association; 3) \textit{B}'s
Fremtræden; 4) Sammenligning af \textit{b} og \textit{B.} Hvis alt dette
virkeligt gik for sig, maatte vi dog utvivlsomt mærke noget til det. Vi maatte
i hvert Tilfælde mærke, at der foregik noget i Bevidstheden, selv om vi ikke
havde ganske Rede paa, hvad det var. (I hvert Tilfælde maa saadanne Forfattere
antage, at Processen mærkes, som afvise den modsatte Forklaring som grundet paa
«overfladisk» Iagttagelse.) Man kan nemlig meget godt have en Fornemmelse af,
at der er flere Bevidsthedselementer til Stede, end man kan opfatte med fuld
Opmærksomhed. Saaledes bemærker \emph{Wundt}%
% footnote 1) on p. 21
\footnote[1]{Physiologische Psychologie. 3. Aufl. II, p. 247.}: «Naar flere
Indtryk frembyde sig end der kan apperciperes, er man sig tydeligt
% --- p. 22 ---
\setcounter{footnote}{0}%
\opage{22}bevidst, at der var endnu andre Indtryk til Stede; men man er ikke
eller dog først ved Hjælp af en tydelig Succession i Stand til bestemt at gøre
sig dem nærværende». Hvad der gælder for Indtryk, maa ogsaa gælde for
Forestillinger. Men det er sikkert, at der er Tilfælde af Genkendelse, hvor der
slet ikke mærkes flere Bevidsthedselementer, end man kan overskue, og hvor
netop dette giver hele Tilstanden dens ejendommelige Præg.

Selv ved successiv Opfattelse af en kendt Genstand synes der ikke at være noget
i Vejen for, at Genkendelsen kan bestaa i en Række af umiddelbare
Genkendelsesakter efter Formlen $\left({a\atop A}\right)$, uden at forud vakte
Forestillinger \emph{behøve} at spille nogen Rolle. Den i 10 omtalte indbildte
Genkendelse peger afgjort i den Retning.

Det er i det Hele en aldeles vilkaarlig Paastand, at Genkendelse altid skulde
forudsætte en \emph{forud} vakt Forestilling. Al Genkendelse beror paa en
Forbindelse mellem en Forestilling og en Fornemmelse, men denne Forestilling
kan meget vel vækkes af selve Fornemmelsen. Ved Sanseillusioner er dette
aabenbart Tilfældet. Paa en Fodtur troede jeg engang at se en lille Sø i det
fjerne, men blev snart opmærksom paa, at det var et skinnende hvidt Hus, jeg
havde perciperet som blinkende Vand. Jeg havde kort i Forvejen set paa Kortet,
som angav en lille Sø i den Retning, i hvilken Huset laa; men jeg havde ikke
tænkt videre derover og søgte aldeles ikke efter Søen, da den pludseligt syntes
at vise sig derude. Illusionen var her aldeles momentan og havde saa liden
objektiv Grund i Indtrykket, at dersom ikke Forestillingen om Søen saa let
havde kunnet vækkes, vilde Illusionen neppe have indfundet sig.

Maaske vil man i Anledning af dette Tilfælde sige, at Forestillingen dog var
til Stede forud, selv om den ikke var bevidst. I det mindste nogle af de
Forfattere, der mene, at Forventningsteorien forklarer al Genkendelse, erklære
udtrykkeligt, at de ved de Forestillinger, som de antage gaa forud ved enhver
Genkendelse, mene «ubevidste Forestillinger» eller Nerveprocesser, til hvilke
der ingen Bevidsthedselementer svare%
% footnote 1) on p. 22
\footnote[1]{Saaledes \emph{Arne Löchen} (Spørgsmaal vedkommende de afasiske
Sygdomme. Kristiania 1888. p. 128 f.) og \emph{Alfr}. \emph{Lehmann} (Om
Genkendelse. Det kgl. danske Videnskabernes Selskabs Skrifter, 6. Række, II, 2.
p. 15). Derimod faar man af \emph{Kroman's} Fremstilling i «Kortfattet Tænke-
og Sjælelære». 2. Udg. København 1888. p. 146 ff. det Indtryk, at han tænker
sig hele den sammensatte Proces som bevidst. Der tages i hvert Tilfælde intet
Forbehold.}. Saa snart man opfatter de medvirkende Forestillinger paa denne
Maade, hører den hele Strid egentlig op. Ti Forestillinger, der ikke ere over
Bevidsthedens Tærskel, existere kun som Forestillingsdispositioner, og det
afgørende Spørgsmaal er, hvorledes disse Dispositioner vækkes til Aktualitet.
Her paastaar jeg nu, at det kan ske og ofte sker ved den tilsvarende
Fornemmelses Gentagelse. Det følger af sig selv, at denne Gentagelse i sig selv
ingen Betydning vilde have for Genkendelsen, naar der ingen Disposition (Spor,
Residuum, Efterklang eller hvad man nu vil kalde det) var, som kunde vækkes.
\emph{Mellemrummet}, i \emph{hvilket}
% --- p. 23 ---
\setcounter{footnote}{0}%
\opage{23}\emph{Forestillingen blot har været til Stede som Disposition}
(«\emph{under Tærskelen}»), \emph{kan naturligvis være af meget forskellig
Længde}, og hele Stridsspørgsmaalet kan her let reducere sig til et rent
Gradsspørgsmaal. Naar jeg efter 20 Aars Forløb møder et Menneske igen og strax
genkender det, har Dispositionstiden (om jeg maa kalde den saa) været meget
lang. Ved den nys omtalte Sanseillusion var Dispositionstiden kun et
Kvarterstid omtrent. Ved Genkendelsen af Ornamentet (i 9) var Dispositionstiden
neppe et Minut. Selv om vi ikke vilkaarligt kunne fremkalde en Forestilling, og
selv om den end ikke ved i og for sig naturlige Anledninger kommer frem i
Bevidstheden, kan den dog ikke siges med Sikkerhed at være aldeles forsvunden
som Disposition%
% footnote 1) on p. 23
\footnote[1]{\emph{Leibniz} siger med Rette: «Je ne vois aucune nécessité, qui
nous oblige d'assurer, qu'il ne reste aucune trace d'une perception, quand il
n'y en a pas assez pour se souvenir qu'on l'a eue.» (Nouveaux Essais. I, 3, 24.
Bemærkningen er rettet mod \emph{Locke}, der negtede, at man kunde siges at
have en Forestilling, naar den ikke kunde fremkaldes i Erindringen.)}. Det kan
være, at der skal en ganske særegen Kombination af Omstændigheder til for at
vække Dispositionen til Aktualitet. Som vi i det Hele ere for tilbøjelige til
at udlede et Fænomen af en eneste Aarsag eller, maaske rettere, til at finde
Aarsagen i en eneste Omstændighed, saaledes lade vi os ogsaa let af vor
skematiske Opfattelsesmaade hilde i den Mening, at en Forestillingsassociation
kan bringes i Stand ud fra ét givet Bevidsthedselement. Hvad et enkelt Element
ikke vilde være i Stand til at udløse, vilde maaske flere kunne bevirke, naar
de fulgte saa hurtigt efter hverandre, at deres Virkninger kunde kumuleres. Vi
kunne fra flere forskellige Sider have faaet Stød i Retning af en vis
Forestillingsassociation, saa at denne derved bliver saaledes forberedt, at nu
en ubetydelig Anledning er nok til at faa Forestillingen op over Bevidsthedens
Tærskel%
% footnote 2) on p. 23
\footnote[2]{R. \emph{Wahle} har henledet Opmærksomheden paa saadanne
«supplerende og konkurrerende» Forhold ved Associationen i sin Afhandling
«Bemerkungen zur Beschreibung und Eintheilung der Ideenassociationen»
(Vierteljahrsschrift für wissenschaftliche Philosophie IX) p. 414---420.}. Ved
alle Genkendelser og Associationer virke sikkert saadanne Forhold med. Vort
Sjælelivs Traade ere saa inderligt sammenvævede eller sammenfiltrede, at de
Principer og Love, vi kunne formulere i vor Psykologi, kun kunne give de aller
tarveligste Omrids.

En \emph{aldeles uforberedt} Genkendelse gives der i strengeste Forstand ikke,
eftersom al Genkendelse forudsætter en Disposition, som stammer fra den første
Opfattelse af det Fænomen, der nu skal genkendes. Men der er alle mulige Grader
af Forberedelse. I nogle af de i 3 omtalte Tilfælde gik der forud for
Genkendelsen en vilkaarlig Søgen, som dog ikke førte til Maalet, men som dog
neppe har været aldeles uden Betydning, eftersom Erfaring hyppigt viser, at
selv om vi ikke finde den søgte Forestilling strax, kan den senere, uden ydre
Anledning, dukke op i Bevidstheden. En bestemtere Forberedelse have vi, naar
Forestillingen kort i Forvejen ved vilkaarlig eller uvilkaarlig Genkaldelse har
været i Bevidstheden, selv om vi ikke fastholde den uafbrudt, til Fænomenet
kommer igen (
% --- p. 24 ---
\setcounter{footnote}{0}%
\opage{24}hvad vi jo heller ikke formaa --- smlgn. 5). --- Men alt dette er kun
Gradsforskelligheder. Naar Forventningsteorien lader sin «Forventning» være
ubevidst, er den egentlig med det samme opgiven. Man vil ikke kunne negte, at
det er den \emph{følgende Fornemmelse}, \emph{der hjælper den}
«\emph{ubevidste}» \emph{Forestilling over Tærskelen}. Det af \textit{A vakte
ubevidste} \textit{b vilde for evigt forblive} i den \emph{ubevidste Verden},
\emph{hvis} \textit{B ikke} kom til og \emph{hjalp det frem}. Og da \textit{b}
og \textit{B} ved denne Lejlighed ikke altid fremtræde som selvstændige
Bevidsthedselementer, faa vi den ofte omtalte Sammensmeltning eller umiddelbare
Genkendelse%
% footnote 1) on p. 24
\footnote[1]{Allerede i min Psykologi p. 141 har jeg gjort opmærksom paa, at
Genkendelse i mange Tilfælde kun bliver mulig ved en Forberedelse, der ligger
et vist \emph{lille} Tidsrum forud.}.

Et Exempel vil klart vise dette. Jeg afskriver nogle Linier, som findes i en
Bog p. 270. Er Sideangivelsen nu ogsaa rigtig? Jeg ser efter i Bogen og finder,
at Siden, paa hvilken de citerede Linier staa, ganske rigtigt er p. 270. Idet
mit Øje falder paa Sidens Nummer i Bogen, faar jeg den Fornemmelse, der
udtrykkes ved Ordene «ganske rigtigt»; den oftere omtalte Bekendthedskvalitet
gør sig gældende. Jeg har dog i det lille Mellemrum ikke udtrykkeligt fastholdt
Tallet 270; tvertimod fór forskellige Forestillinger fra helt andre Regioner
hurtigt igennem Bevidstheden. Og jeg kom heller ikke, da Øjet faldt paa det
trykte Tal 270, til at tænke paa det Tal, jeg havde skrevet. Det trykte Tal
genkendtes umiddelbart, og hvis jeg ikke havde set det, vilde Forestillingen om
det skrevne Tal være vedbleven at være latent (med mindre den vilde blive
dragen frem af en anden Aarsag, f. Ex. af Ønsket om at verificere den).

Man finder med Hensyn til dette Forhold næsten hos alle psykologiske Forfattere
en Misforstaaelse, som hænger sammen med en Forvexling af den
\emph{psykologiske Forklaring} af Forestillingsvirksomheden med dennes
\emph{logiske Betydning}. Den Forestilling, som gør os en anden Forestilling
bekendt og forstaaelig, betragter man gerne som den aktive; den, der finder sin
Bestemmelse eller Karakterisering, betragtes som den passive. Efter Herbarts
Sprogbrug: hin kaldes den \emph{apperciperende}, denne den
\emph{apperciperede}. Naar jeg ser en Genstand for mig og siger: «det er en
Fisk!» saa kaldes Forestillingen «Fisk» den apperciperende Forestilling. Men
det er dog i Virkeligheden den øjeblikkelige Fornemmelsesgruppe, som sætter den
hele Proces i Gang; det er den, der vækker Almenforestillingen «Fisk»; det er
altsaa den, der er den aktive Part i den Proces, der foregaar. Og saaledes maa
Forholdet ogsaa opfattes ved enhver Genkendelse (for saa vidt vi her teoretisk
kunne sondre mellem to Elementer). Det følger af sig selv, at der ikke
tilkommer noget af Leddene absolut Aktivitet; den givne Fornemmelsesgruppes
Betydning er kun at \emph{udløse} Reproduktionen, at frigøre opsamlet Energi.
Men hin Sprogbrug synes ikke at være heldig, da den gaar ud fra det logiske
Udbytte, og ikke fra selve den psykiske Proces.
% --- p. 25 ---
\setcounter{footnote}{0}%
\opage{25}Kun ved den egentlige, med bevidst Søgen efter logisk Udbytte
forbundne Tænken begynde vi med det «apperciperende» (fortolkende og maalende)
Element, idet vi da have en foreløbig Forestilling om, hvad vi søge, eller
hvilke Fordringer det Søgte skal tilfredsstille.

Jeg tilføjer endnu, hvad der iøvrigt er tilstrækkeligt fremhævet ovenfor (4),
at Teorien om den umiddelbare Genkendelse egentlig slet ikke gør Brug af
«ubevidste Forestillinger». Hvad jeg \emph{teoretisk} udtrykker som
Sammensmeltning af en Fornemmelse og en Forestilling, er den Forandring, en
Fornemmelse kan undergaa ved Gentagelse. Det er Øvelsens Virkning, og intet
andet. --- \emph{Sekundært} kan dog (8) hele Genkendelsen blive en --- i hvert
Tilfælde saa godt som --- ubevidst Akt.

Det staar endnu tilbage at drøfte de Forsøg, man dels ved Benyttelse af
patologiske Tilfælde, dels ad experimental Vej har gjort paa at modbevise den
umiddelbare Genkendelses Tilværelse.

12. Dr. \emph{Arne Löchen} har i sit allerede nævnte Skrift villet modbevise
Antagelsen af umiddelbar Genkendelse gennem en Analyse af de mærkelige
patologiske Tilfælde, som kaldes Ordblindhed og Orddøvhed. Det er saadanne
Tilfælde, i hvilke Patienten vel \emph{opfatter} skrevne eller hørte Ord som
blotte Sansefænomener, men ikke \emph{forstaar} dem, d. v. s. ikke opfatter dem
som Ord eller Forestillingstegn. Ved Spørgsmaalet om, hvorledes saadanne
Tilfælde ere at forklare, staa to Opfattelser overfor hinanden. Den ene
(Kussmaul, Charcot o. fl.) mener, at Genkendelsen er falden bort, medens den
blotte Sansefornemmelse er uskadt (efter det tidligere brugte Udtryk:
\textit{A} opfattes, men kan ikke mere kombineres med \textit{a}). Den anden
(Stricker, Lichtheim) mener, at hvad der fattes, er Muligheden af at vække
saadanne associerede Forestillinger fra andre Sanseomraader, ved hvilke Ordenes
Betydning forstaas. Et synligt Ord forstaas nemlig efter denne Opfattelse
derved, at det vækker en Forestilling om den tilsvarende Lyd eller om den til
Udtalen nødvendige Bevægelse af Stemmemusklerne.

Imod den første Opfattelse strider ifølge Löchen for det første den
Omstændighed, at de Ordblinde synes at være i Stand til at kopiere. Denne Evne
beviser nemlig --- mener han --- at Synsforestillingerne (de optiske
Erindringsbilleder) intet have lidt. --- Denne Indvending forekommer mig at
have meget lidet at sige, naar man erindrer den sekundære Ubevidsthed, hvormed
indøvede Handlinger kunne foretages. Hos øvede Afskrivere finder der neppe mere
nogen bevidst Genkendelse af Skrifttegnene Sted; Bogstaverne gaa, saa at sige,
lige fra Øjnene ned i Fingrene. Denne erhvervede og indøvede Færdighed behøver
altsaa ingenlunde at falde bort, fordi Genkendelsesevnen er falden bort. Der er
ved Øvelsen etableret en Genvej, som meget godt kan være farbar, fordi den op%
% --- p. 26 ---
\setcounter{footnote}{0}%
\opage{26}rindelige Omvej er ufremkommelig. --- Dr. Löchen indrømmer da ogsaa
selv (p. 57) Muligheden af, «at Bogstaverne kunne forekomme fremmede, selv om
man kan kopiere dem».

Der er for det andet --- ifølge Forf. --- en vigtig Omstændighed, som Teorien
om umiddelbar Genkendelse har overset, nemlig den, at Ordblindhed i Regelen er
forbunden med større eller mindre Vanskelighed ved at finde Navnene paa
Genstandene. Den rigtige Forklaring vil, mener han, findes ved at gaa ud fra
denne Kendsgerning. Ordblindheden tænkes da at opstaa ved, at Vejen fra Synets
Sansningscentrum til Lydforestillingernes eller Bevægelsesforestillingernes
Centrum er beskadiget. Forudsætningen for denne Forklaring er den, at al
Genkendelse af et synligt Ord sker ved, at det synlige Tegn vækker den ved
Berøringsassociation dertil knyttede Forestilling om Ordets Lyd eller om den
Bevægelse, vi foretage med Stemmemusklerne, naar vi udtale det. Paa Vækkelsen
af denne Association beror «den subjektive Følelse af Genkendelse» (ved hvilket
Udtryk Löchen betegner, hvad jeg har kaldet Bekendthedskvalitet). Det er klart,
at hvis Genkendelse ikke bestaar i andet, maa den falde bort, naar Sammenhængen
mellem Syn og Hørelse eller Talebevægelse afbrydes.

Men Dr. Löchen har rigtignok taget sig det meget let med at godtgøre, at «den
subjektive Følelse af Genkendelse» kan forklares paa denne Maade. Det skal
efter hans Mening være tilstrækkeligt, at en Sansefornemmelse vækker
Forestillinger hos mig, ligegyldigt hvilke, saa er den genkendt. «Hvis vi traf
en Blomst», siger han (p. 103), «og denne fremkaldte et Artsnavn, der var
fejlagtigt, vilde Fremkaldelsen af Navnet visseligt være nok til, at vi
erklærede Blomsten for kendt; kommer imidlertid En, som har bedre Kundskab, og
fortæller os, at vi har taget Fejl af Navnet, erkender vi naturligvis, at vor
Genkendelse var fejlagtig... Men \emph{Genkendelsen} i og for sig \emph{ligger
visseligt} i \emph{Fremkaldelsen af Associationerne} og ikke i Stadfæstelsen af
dens Rigtighed. Her er jo kun Tale om, hvori den subjektive Følelse af
Genkendelsen ligger, ikke om dens objektive Sandhed.» Denne sidste Bemærkning
tiltræder jeg ganske; men jeg kan ikke billige den Anvendelse, Forf. gør af
den. Jeg har allerede bestemt skelnet mellem Genkendelsen og dens Verifikation
(10). Men den blotte Fremkaldelse af en Forestilling, ligegyldigt hvilken, er
ikke nok til Genkendelse. Den saaledes fremkaldte Forestilling kan jo selv
undertiden staa aldeles fremmed for os, hvis jeg nemlig har glemt den i den
Grad, at den ikke selv fremtræder med Bekendthedskvaliteten. Spørgsmaalet
rykkes saaledes kun længere ud. (Jeg ser her bort fra, hvad jeg senere kommer
til, om ikke Berøringsassociationen selv stedse forudsætter en Genkendelse.)
--- Det synes, som Dr. \emph{Löchen her skiller} sig fra
\emph{Forventningsteorien}, som han \emph{ellers slutter} sig til.
Forventningsteorien gik jo netop ud fra, at det ikke var nok til Genkendelse,
at en Forestilling vaktes, men søgte at vise, at Genkendelsen først indtraadte,
naar de ved de først opfattede Dele af Genstanden vakte Forestillinger
befandtes overensstemmende med de af de senere
% --- p. 27 ---
\setcounter{footnote}{0}%
\opage{27}fremtrædende Dele vakte Fornemmelser. Der finder efter
Forventningsteorien en uafbrudt Verifikation Sted. (Fejltagelse synes derfor
næsten at være udelukket ved denne Teori, som forvexler Genkendelsen med dens
Verifikation. Smlgn. 10.) --- Det er i Virkeligheden en tredie Teori, som her
træder frem for os, men en Teori, som synes at være uheldigere end nogen af de andre.

Skulde der ikke kunne finde Genkendelse Sted, uden at nogen fri
Berøringsassociation gør sig gældende? Skulde man ikke kunne genkende en Blomst
uden at belægge den med noget som helst enten rigtigt eller galt Navn? De i 3
anførte Tilfælde synes at tale derfor. Jeg kender to Smaapiger, hvis Navne jeg
stadigt forvexler og derfor stadigt er usikker paa, skønt jeg aldrig forvexler
dem selv, men sikkert genkender dem hver for sig, og jeg er mig ved
Genkendelsen af dem aldrig nogen som helst hjælpende Forestilling bevidst. ---
Forf. har ladet sig vildlede af den stærke Trang, vi i Regelen have, til fra en
given Fornemmelse eller Forestilling at gaa over til andre Forestillinger eller
til Bevægelse (bl. a. Tale). Den spontane Forestillings- og Bevægelsestrang
lader, som allerede omtalt (se 10), sjeldent den umiddelbare Genkendelse staa
isoleret og afsluttet, men fører ud over den til Verifikationer eller til nye
Genkendelsesakter. --- Det er et fortjenstligt Arbejde, Dr. Löchen har paataget
sig, at underkaste de herhen hørende patologiske Fænomener en psykologisk
Analyse, og jeg skylder hans rolige og objektive Fremstilling adskillig
Belæring. Men Beskæftigelsen med de patologiske Fænomener maa dog ikke aflede
Opmærksomheden fra de Fænomener, det normale Liv frembyder.

Hvad selve de patologiske Tilfælde angaar, er der Grund til at undre sig over,
at Forf. ikke har opsøgt saadanne Fænomener, der kunde give Oplysning om
Patientens Bevidsthedstilstand under Ordblindhed. Det var værdt at undersøge,
om der ikke ved Ordblindhed (ved Orddøvhed synes Forf. at indrømme det for de
fleste Tilfældes Vedkommende) finder en anden Opfattelse af Tegnene Sted end
tidligere, om disse ikke staa i et andet Lys, med en anden Kvalitet, end naar
de genkendes. I det mindste ét Tilfælde, i hvilket en intelligent Patient har
skildret sin Tilstand under Lidelsen, peger afgjort i den Retning. Den bekendte
Læge \emph{Lordat} fra Montpellier led en Tidlang af Ordblindhed og Orddøvhed.
Sin Tilstand beskrev han nærmere paa følgende Maade: «En perdant le souvenir de
la signification des mots entendus, j'avais perdu celui de leurs signes
visibles. La syntaxe avait disparu avec les mots, \emph{l'alphabet seul m'était
resté}, mais la jonction des lettres pour la formation des mots était une étude
à faire. Lorsque je voulus jeter un \emph{coup} d'œil sur le livre que je
lisais, quand ma maladie m'avait atteint, je me vis dans l'impossibilité d'en
lire le titre. Il m'a fallu épeler lentement la plupart des mots.» Lordat
skildrer, siger \emph{Kussmaul}%
% footnote 1) on p. 27
\footnote[1]{Die Störungen der Sprache. Leipzig 1877. p. 175 f. ---
Udhævelserne ere af mig.}, fra hvem jeg henter dette Træk, «med dyb Bevægelse
% --- p. 28 ---
\setcounter{footnote}{0}%
\opage{28}det lykkelige Øjeblik, da han efter nogle sørgelige Uger lod sit Blik
svæve hen gennem sit Bibliothek, og da \emph{uventet} paa et Hjørne de Ord paa
Ryggen af en Foliant: «Hippocratis opera» \emph{igen lyste} ham i \emph{Møde}.»
Denne Skildring synes ikke at stemme med Dr. Löchens Paastand (p. 119) om, at
den Ordblinde ser Bogstaverne lige som et normalt Menneske. Hvad Lordat klager
over, er, at han ser Bogstaver, men ikke Ord. Og dette stemmer fortræffeligt
med hvad ovenfor er anført, at Ord og Sætninger genkendes lettere end
Sammenstillinger af Bogstaver uden Betydning, og at Indtrykkene, naar de
saaledes hørte sammen, kom \emph{simultant} til Bevidsthed (3). Lordat's
Beskrivelse baade af hvad han mistede ved Sygdommen, og af hvad der igen
traadte frem for ham ved dens Ophør, viser, at det netop var Evnen til
umiddelbar Genkendelse, han havde mistet. Der foreligger her et godt Exempel
paa uforberedt Genkendelse (3, 5). Det er ikke let at forstaa, hvorledes Dr.
Löchen (p. 119) kan faa ud, at de Ordblinde se Bogstaverne fuldstændigt paa
samme Maade som sunde Mennesker: det gør dog en ikke ringe Forskel, om jeg blot
ser dem som en Hob af Figurer, eller de «lyse mig i Møde», som Titelen paa
Lordat's Yndlingsbog gjorde for ham i Helbredelsens Øjeblik, med
Bekendthedskvalitetens hele Varme. ---

Det er ikke saa, som Dr. Löchen (p. 11) synes at antage, at ifølge Teorien om
umiddelbar Genkendelse kommer Sansefornemmelsen altid først, og efter et
Tidsmellemrum Erindringsbilledet. Teorien gaar netop ud paa, at vi ved
umiddelbar Genkendelse ikke skelne mellem Fornemmelse og Reproduktion, men at
disse to Virksomheder kun træde frem for os i deres fælles Produkt, det med
Bekendthedskvaliteten forsynede Bevidsthedselement. Det er heller ikke
nødvendigt at antage, at Sansning og Erindring høre til to forskellige
Hjernecentre. Det synes at være saa; men Teorien kan bestaa uden det.

--- Idet jeg tager Afsked med Dr. Löchen, maa jeg endnu anføre, at han i
Modsætning til andre Forfattere, som i denne Sag have udtalt sig med stor
Selvtillid, erklærer, at han «paa Videnskabens nuværende Standpunkt har trot at
maatte indtage en noget forsigtigere Holdning».

13. I sin allerede nævnte Afhandling (Om Genkendelse. Forsøg paa en
experimental Verifikation af Forestillingsassociationernes Teori.
Videnskabernes Selskabs Skrifter. 6. Række. II) har Dr. \emph{Alfr}.
\emph{Lehmann} trot at kunne godtgøre ad experimental Vej, at der ingen
umiddelbar Genkendelse i den ofte angivne Betydning gives.

En saadan Tro maa allerede paa Forhaand vække nogen Forundring. Dersom der
gives forskellige Arter eller Former af Genkendelse, hvorledes vil man saa
egentligt ved Forsøg, der altid maa finde Sted \emph{under visse bestemte
Forhold}, kunne godtgøre, at Genkendelse ikke kan finde Sted paa anden Maade
\emph{under andre Forhold}? Det er dog en klar og simpel Regel for al
experimental Forsken, at dens Resultater kun gælde for
% --- p. 29 ---
\setcounter{footnote}{0}%
\opage{29}de Forhold, under hvilke Forsøgene ere anstillede, og at
Berettigelsen til at udvide dem til \emph{alle} Forhold i hvert Tilfælde maa
særligt godtgøres. Man kan ikke uden videre slutte fra en hypotetisk Dom til en
kategorisk Dom; det kunde jo nemlig være, at Betingelserne vare det afgørende.
Det er Grundfejlen ved Dr. Lehmanns Afhandling, at han ikke har indsét
Forholdet mellem de to Teorier, der stode overfor hinanden. Den ene siger, at
al Genkendelse beror paa Forventninger, der vækkes ved Berøringsassociation og
derefter stadfæstes; den anden siger, at \emph{foruden} den (successive)
Genkendelse, der sker paa denne Maade, gives der ogsaa en umiddelbar
Genkendelse. Det er da klart nok, at man ikke gendriver denne sidste Teori og
ikke leverer noget Bevis for den første, fordi man viser, at i visse Tilfælde,
under visse bestemte Betingelser sker Genkendelsen ved Hjælp af forud vakte
Forestillinger, --- dersom disse Betingelser netop vare saadanne, som
udelukkede Muligheden af en umiddelbar Genkendelse! Saadanne Experimenter
bevise, hvad ingen har negtet, og modbevise --- Intet. --- Og med Hensyn til
Dr. Lehmann's Experimenter maa jeg gaa et Skridt videre. Ti de forudsætte ikke
en Gang Forventningsteorien som nødvendig Forklaring. Kun ved en Række af
Misforstaaelser og Fejlslutninger kan der tillægges dem en saadan Betydning.
Jeg skal søge at bevise, hvad jeg her har paastaaet.

Min Indsigelse gælder ikke selve Dr. Lehmann's Experimenter. De have deres
positive Værdi afset fra hans uheldige teoretiske Fortolkning af dem. De oplyse
forskellige interessante Punkter med Hensyn til Genkendelsens Sikkerhed og
Nøjagtighed under \emph{visse} Omstændigheder, men hvad de lære om
Genkendelsens psykologiske Natur er saare lidet, og i hvert Tilfælde ikke det,
Dr. Lehmann selv mener. Jo mere Betydning man tillægger den experimentale
Psykologi, jo mere man glæder sig over, at vor Viden om Sjælelivet ad denne Vej
kan vinde i Klarhed og Nøjagtighed, des betænkeligere maa man ogsaa være ved at
se den experimentale Forsken sammenblandet med forhastet Teoretiseren.
Psykologien har ikke frigjort sig fra Metafysikens Herredømme for at blive
behersket af en kortsynet Betragtningsmaade, der overser de sjælelige
Fænomeners Mangfoldighed og Forskellighed. Afhandlinger, som ere forsynede med
Tavler og Tabeller, imponere let ukritiske Læsere. Ikke blot for selve
Problemets Skyld, men ogsaa for den experimentale Psykologis egen Skyld har jeg
derfor foretaget en kritisk Prøvelse af den omtalte Afhandling, og skal her
forelægge, hvad jeg har fundet.

a. Dr. Lehmann anstillede en Række Forsøg paa følgende Maade. Der blev vist de
Iagttagere, med hvilke der experimenteredes, først en graa Skive, Normalskiven
(\textit{n}), dernæst efter et vist Mellemrum enten en lysere Skive
(\textit{l}) eller en mørkere (\textit{m}) eller selve Normalskiven igen, og
Iagttagerne skulde saa afgøre og markere, om den Skive, de sidst fik at se, var
lig Normalskiven eller lysere eller mørkere end denne. --- Allerede denne
Anordning af Forsøgene viser, at den ovenfor omtalte umiddelbare og uforberedte
Genkendelse var udelukket. Den Genkendelse, der under disse Forhold kunde finde
Sted, maatte være for%
% --- p. 30 ---
\setcounter{footnote}{0}%
\opage{30}beredt paa to Maader. For det første er der saa kort et Mellemrum
mellem den første og den anden Skives Fremtræden, at Forestillingen om den
første maa siges at være tilstede, om ikke i Bevidstheden, saa dog som en
Disposition, der ikke er langt fra Bevidsthedens Tærskel, og der behøver altsaa
ikke at anvendes saa megen Energi til dens Reproduktion, som hvis det Fænomen,
der skulde genkendes, kom aldeles uventet og uforberedt. For det andet ligger
det i Sagens Natur, at naar man stiller en Person op for at prøve hans
Genkendelsesevne i en vis bestemt Henseende, saa vil hans Opmærksomhed
uvilkaarligt ogsaa vende sig i den Retning; han staar endnu mere forberedt
overfor den anden Skive, end han f. Eks. er paa Gaden, naar han først har mødt
et Menneske og derpaa tilfældigvis ogsaa møder et Menneske, og han da
uvilkaarligt finder, at dette er ligt hint eller forskelligt fra det. Den, der
gaar hen ad Vejen og ved en Havelaage uventet ser en Dame, i hvis Mine der er
ham noget bekendt, foretager Genkendelsesakten under andre Forhold end
Iagttagerne i Dr. Lehmanns Laboratorium. Det samme gælder om Lordats
Overraskelse ved at kunne læse Titelen paa sin Bog. Det ligger i Sagens Natur,
at man ad experimental Vej ikke kan undersøge den umiddelbare Genkendelse. Man
kan ikke undlade at forberede den paa de to anførte Maader. Den umiddelbare
Genkendelse maa man undersøge ved at aflure det virkelige Liv Tilfælde og
Exempler, der kunne belyse Fænomenet. Det forholder sig hermed som med saa
mange af det uvilkaarlige Sjælelivs Fænomener: den experimentale Psykologi kan
kun højst indirekte kaste Lys over det, og der maa gaas meget varsomt frem med
Overførelsen af dens Resultater paa dette Omraade. Hvad der gælder for
Selviagttagelsen i det hele, at den ved at rette Opmærksomheden mod det
sjælelige Fænomen strax ændrer dette; det gælder ikke mindst for den
experimentale Psykologi, der stedse maa anbringe og præparere sine
Forsøgsgenstande (Iagttagerne) paa en vis bestemt Maade, som det uvilkaarlige
Sjæleliv ikke behøver at være bundet til%
% footnote 1) on p. 30
\footnote[1]{Paa karakteristisk Maade fremtræder Præparationen i en Ytring af
Dr. Lehmann p. 27: «I det Moment, da Genkendelsen skal [sic] foregaa, søger man
ved vilkaarlig [sic] Koncentration af Opmærksomheden at genkalde sig den
tidligere Fornemmelse». Ja, det er netop Sagen: man véd, at man \emph{skal}
genkende; derfor er man jo i Laboratoriet, og derfor er man stillet op overfor
Skiverne o. s. v. Som Følge af denne Præparation anstrenger man sig derfor
efter bedste Evne. Forsøgene vise derfor, hvad man ved Anvendelse af «fuld
Opmærksomhed» (smlgn. p. 42) kan udrette i Retning af Genkendelse under visse
givne Omstændigheder og overfor visse givne Genstande. Men de vise intet om
Genkendelsens Natur, hvor den ikke er anordnet i Forvejen.}. Det er altsaa
klart, at Dr. Lehmann's Forsøg, hvorledes man end vil fortolke deres
Resultater, aldeles intet kunne bevise om umiddelbar Genkendelse, da dennes
Mulighed paa Forhaand var udelukket.

b. Dr. Lehmann gaar ud fra, at naar der er Tale om Genkendelse af en
Fornemmelse, maa man \emph{naturligt} (naar man ikke «gaar ud fra særlige
teoretiske Forudsætninger») antage, at den Forestilling, ved hvis Sammenligning
med den givne Fornemmelse Genkendelsen bliver til, er given \emph{forud}. Han
beraaber sig her paa \emph{Wolfe}, der i sine «
% --- p. 31 ---
\setcounter{footnote}{0}%
\opage{31}Untersuchungen \emph{über} das Tongedächtniss» (Wundt Studien. III.
p. 556) siger: «Naar vi nærmere gaa ind paa Fremgangsmaaden ved Sammenligningen
af to ved et Mellemrum adskilte Toner, saa er det klart, at uden et
Erindringsbillede af den første Tone er en Sammenligning overhovedet umulig».
Men Dr. Lehmann maa, da han læste Wolfe, selv være «gaaet ud fra særlige
teoretiske Forudsætninger»; ellers vilde han ved at læse to Sider videre i
Wolfes Afhandling have set, at denne virkeligt uhildede Forsker udtrykkeligt
siger (p. 558), at det «som bekendt» ikke er nødvendigt at beholde et Billede
af den første Tone i Bevidstheden, «da den \emph{anden Tone strax fremkalder et
Billede} af den \emph{første}». Med andre Ord: Wolfe gaar ud fra den modsatte
Opfattelse af den, Dr. Lehmann lader ham fremsætte som den ene rimelige%
% footnote 1) on p. 31
\footnote[1]{I sin Afhandling «Beiträge zur Theorie der sinnlichen
Aufmerksamkeit und der activen Apperception» (Wundts Studien IV. p. 407)
antager N. \emph{Lange} paa analog Maade, at det ved opmærksom Betragten af en
Genstand er Fornemmelsen, der udløser det tilsvarende Erindringsbillede og
derved forstærkes. --- Smlgn. ovenfor 5.}. Naar nu den anden Tone saaledes
fremkalder Forestillingen om den første, saa kan det ikke ske ved
Berøringsassociation, da de to Toner ikke ere forekomne sammen. Det kan kun ske
ved den Lighedsassociation, Dr. Lehmann forkaster, men som her anviser den
eneste mulige Forklaring. --- Jeg anfører kun dette for dertil at knytte den
Bemærkning, at selv om al Genkendelse foregik paa den Maade, Dr. Lehmann har
ladet den foregaa i sit Laboratorium, nemlig \emph{forberedt}, saa vilde derved
Forventnings- eller Berøringsteorien ikke være bevist. Ti Billedet af Skiven
\textit{n} holder sig ikke i Bevidstheden i hele Mellemtiden, indtil \textit{l}
eller \textit{m} eller selve \textit{n} kommer; det bliver fremkaldt paa ny, og
at det bliver fremkaldt, kan dog kun skyldes den anden Skives Fremtræden.
Reproduktionen sker altsaa \emph{efter} eller ved Indtrædelsen af det Fænomen,
som skal genkendes, medens Forventnings- eller Berøringsteorien forudsætter, at
den gaar \emph{forud}. (Smlgn. iøvrigt 11.)

c. Den Art Genkendelse, som foreligger ved de Lehmannske Forsøg, afviger nu
ikke blot fra den umiddelbare Genkendelse, Erfaringen kan lære os at kende, ved
at være forberedt og halvt vilkaarlig, men ogsaa ved at angaa saa fine Nuancer,
saa simple Fornemmelser, som kun ved særlig Anspændelse af Opmærksomheden kunne
genkendes. Dette bemærker Dr. Lehmann selv. Men han overser, at der mellem
saadanne fine og simple Fænomener, der kræve en given Maalestok for at kunne
genkendes, og de mere sammensatte Fænomener, hvis Opfattelse, og altsaa ogsaa
hvis Genkendelse maa ske successivt, ligger en Klasse af lidet sammensatte
Fænomener, der kunne opfattes simultant og genkendes i et Øjeblik. (Smlgn.
3---4.) Jo mere vi gaa i det Enkelte, eller jo mere vi gaa i det Indviklede,
des mere vidtløftig og vanskelig bliver Genkendelsesakten. Vi blive strax klare
over, at det er en rød Farve, vi se paa Postbudets Frakke, men kun ved en mere
% --- p. 32 ---
\setcounter{footnote}{0}%
\opage{32}indviklet Tankeproces kunne vi afgøre, om dette Røde er ligt med
eller forskelligt fra det Røde, som Gardister eller kongelige Lakajer bære%
% footnote 1) on p. 32
\footnote[1]{Smlgn. \emph{Spencer}: Principles of Psychology. § 115.}. Fordi
Genkendelsen af Gardistrødt ikke er umiddelbar, kan Genkendelsen af Rødt i det
hele godt være det. Fordi Genkendelsen af en Række Bogstaver \emph{uden} Mening
ikke er umiddelbar, kan Genkendelsen af en Række Bogstaver med Mening godt være
det. --- Det maa stadigt fastholdes, at Striden staar mellem to Teorier, af
hvilke den ene antager al Genkendelse for forberedt og successiv, medens den
anden antager, (ikke: at al Genkendelse er umiddelbar, men:) at der foruden hin
forberedte og successive Genkendelse \emph{ogsaa} gives en umiddelbar Genkendelse.

d. Medens de to Teorier, som Dr. Lehmann viser, hver for sig meget vel kunne
forklare flere af de Resultater, der fremkom ved Forsøgene, skal der med Hensyn
til et enkelt Punkt findes en skarp Modsigelse mellem disse Resultater og de
Konsekventser, der maa drages af Teorien om umiddelbar Genkendelse. Jeg kan
efter gentaget Studium af Dr. Lehmann's Afhandling ikke se andet, end at han
her paa en paafaldende Maade siger sig selv imod, idet han et Sted af sin egen
Teori drager netop de samme Konsekventser, som han et andet Sted drager af den
modstaaende Teori og finder stridende mod Forsøgenes Resultater. Jeg er dog
noget usikker her, da Dr. Lehmann ikke udtrykker sig tilstrækkeligt klart. Selv
om jeg ikke skulde have Ret i, at Dr. Lehmann her hilder sig i en
Selvmodsigelse, gør det iøvrigt intet til Sagen; ti jeg mener tilstrækkeligt at
have godtgjort, at umiddelbar Genkendelse paa Forhaand var udelukket ved de
omtalte Forsøg, og at intet fornuftigt Menneske kunde falde paa at mene, at al
Genkendelse skulde være umiddelbar.

Foruden Normalskiven \textit{n} blev der ved Forsøgene, som ovenfor omtalt,
vist Iagttagerne to andre Skiver \textit{l} og \textit{m.} Hvert Forsøg
begyndte med, at \textit{n} blev vist, og efter et lille Tidsmellemrum blev saa
enten \textit{l} eller \textit{m} eller \textit{n} vist. \textit{n} blev altsaa
vist langt flere Gange end \textit{l} eller \textit{m.} Nu slutter Dr. Lehmann:
Dersom der gaves umiddelbar Genkendelse, maatte \textit{n} genkendes med langt
større Sikkerhed end \textit{l} eller \textit{m}, da den kommer langt hyppigere
igen; Teorien om umiddelbar Genkendelse bygger jo paa Øvelsens Lov; --- men
Forsøgene viste, at \textit{n} ikke genkendtes sikrere end \textit{m}
og\textit{ l,} --- og dermed er da hin Teori dødsdømt! Den modstaaende (ɔ:
efter Dr. L.'s Mening modstaaende) Teori skal derimod ikke komme i Strid med
Forsøgsresultaterne, ti, hedder det (p. 30), «Genkendelsen beror her paa en
Sammenligning med et Erindringsbillede, der er mer eller mindre udvisket og
uklart, og paa hvilket den stadige Gentagelse kun har den Indvirkning, at det
efter en vis Tids Forløb ikke er slet saa uklart, som det vilde være, hvis
Øvelsen havde været mindre». Dette er den Konsekvents, som p. 30 drages af
«Berøringsteorien». --- Ganske anderledes lyder den Konsekvents, som p. 15
drages af den samme Teori. Her hedder det, efterat det
% --- p. 33 ---
\setcounter{footnote}{0}%
\opage{33}er udviklet, hvorledes Berøringsteorien forklarer Genkendelsen af
successivt opfattede Genstande: «Af denne Fremstilling forstaas det let,
hvorfor en \emph{Forestilling genkendes desto lettere}, \emph{jo hyppigere} den
har \emph{været} i \emph{Bevidstheden}. Desto tydeligere maa nemlig
Erindringsbillederne antages at blive, og \emph{desto sikrere} vil man derfor
ogsaa kunne vurdere det sansede kongruent med de reproducerede Fornemmelser».
Jeg ser ikke rettere, end at der her af Berøringsteorien drages den samme
Konsekvents, som hist af, hvad Dr. L. kalder «Lighedsteorien», og hvis der er
en Modstrid med Forsøgsresultaterne, er den altsaa til Stede for begge Teoriers
Vedkommende. --- Denne Selvmodsigelse kan jeg ikke faa bort fra Dr. L.'s Afhandling.

Ogsaa en anden Mislighed synes der at være til Stede ved de Deduktioner, paa
hvilke Dr. L. bygger Fortolkningen af sine Forsøg. Hvis Genkendelsen var
umiddelbar, saa maatte, mener han, ved Forsøg med to Skiver, Normalskiven
\textit{n} og \textit{l}, \textit{n} genkendes lettere end \textit{l.} Men
dette følger ingenlunde. Her er jo nemlig ikke Valg mellem Genkendelse af
\textit{n} og Genkendelse af \textit{l}, men mellem at finde den anden Skive
(\textit{n} eller \textit{l)} lig med eller forskjellig fra den første
(\textit{n}). Og dette er en anden Sag. Det er jo dog ingenlunde givet, at
Letheden i at genkende \textit{n} ved Øvelse skulde tiltage mere end Evnen til
at finde Forskel mellem \textit{n} og \textit{l.} Forskelsopfattelsen, som
spiller saa stor en Rolle ved alle vore Fornemmelser, vilde meget vel kunne
voxe lige saa meget i Sikkerhed ved den mindre Øvelse, som Lighedsopfattelsen
ved den større Øvelse. For at opfatte Forskellen mellem \textit{l} og
\textit{n} behøver man ikke en saa fuldstændig og tydelig Reproduktion af
\textit{n}, som til at opfatte Identiteten af \textit{n} og \textit{n.} Det maa
jo tilmed erindres, at Genkendelse og Forskelsopfattelse her sker med «fuld
Opmærksomhed». --- Denne Betragtning vil ogsaa kunne anvendes paa Forsøgene med
alle tre Skiver; det kan være forholdsvis lettere, og derfor kræve mindre
Øvelse, at skelne end at identificere, at finde noget lysere eller mørkere, end
at finde det lige lyst. ---

e. I en anden Række Forsøg gaar Dr. L. ud paa at undersøge, hvilken Indflydelse
det har, naar den Ting, der skal genkendes, har et bestemt Navn. Han gaar ud
fra den daglige Sprogbrug, der betegner noget som kendt, naar man kan nævne
dets Navn. Gælder det nu at genkende en Række Nuancer f. Ex. af Graat, saa vise
Dr. L.'s Forsøg, at denne Genkendelse sker med større Nøjagtighed, naar vi have
Navne for alle Nuancerne. Den er nøjagtigere ved de fem Graatnuancer, som
Sproget har Navne for (hvidt, lysgraat, mellemgraat, mørkegraat, sort), end ved
sex eller ni Nuancer. Og da Dr. L. øvede sig i at associere de ni første Talord
med de ni Nuancer af Graat, viste der sig ogsaa større Sikkerhed ved
Genkendelsen, end da de skulde genkendes uden at være markerede hver for sig.
Heraf drager Dr. L. saa den Slutning, at Genkendelse oftest sker ved, at den
gentagne Fornemmelse fremkalder en Forestilling om dens Navn.

% --- p. 34 ---
\setcounter{footnote}{0}%
\opage{34}Ogsaa her maa jeg skelne mellem de i og for sig interessante Forsøg
og den i mine Øjne aldeles unaturlige og uberettigede Fortolkning, de gøres til
Genstand for. --- Afset fra, hvad jeg allerede har havt Lejlighed til at
indvende mod den Paastand, at Genkendelse er Et med Evne til at benævne (12),
har jeg saa mange Indvendinger at gøre mod denne Fortolkning, at det er bedst
at numerere dem.

1) Usikkerheden ved Genkendelsen maa nødvendigvis være større, naar det drejer
sig om sex eller ni Nuancer, end naar det kun drejer sig om fem, hvad enten vi
saa have Navne til dem eller ikke. Enhver Hypotese maa antage dette, og kun
langvarig Øvelse kan ændre noget derved.

2) Det er klart, at selve Indøvelsen af Forbindelsen mellem Farvenuance og Navn
er en Øvelse i Genkendelse. Det er da intet Under, at Genkendelsen bliver
sikrere ved de benævnte Nuancer end ved de ubenævnte; men deraf følger
rigtignok ikke, at «det er Farvenavnene, der gør Genkendelsen efter længere
Tids Forløb mulig». Benævnelsen er kun en Sidevirkning; Aarsagen baade til
Genkendelsen og til Benævnelsen er den hyppigere Gentagelse og den udtrykkelige
Opmærksomhed, man maa rette mod Farvenuancen for at knytte Navnet til den. (I
et følgende Afsnit af denne Afhandling vil jeg komme nærmere ind paa dette
Punkt og godtgøre, at al Berøringsassociation --- altsaa ogsaa den mellem
Fornemmelse og Navn --- forudsætter en Genkendelse af første Led.)

3) Naar Nuancen \textit{A} har vakt Ordforestillingen \textit{b}, saa sker der
i det Øjeblik, det kendte Navn udsiges eller levende forestilles, en
Genkendelsesakt for \textit{b}'s Vedkommende, i Anledning af hvilken hele
Stridsspørgsmaalet om umiddelbar Genkendelse kan rejses paa ny. Naar vi til
dagligt Brug saa ofte slaa os til Taals ved at kunne give en Ting et Navn
(«give den et Navn og lade den løbe») og mene deri at have en vis Garanti, saa
kan det simpelthen forklares ved, at Bekendthedskvaliteten kan være kraftigere,
saa at sige mere mættet, ved Navnet end ved selve Tingen. Der virker her det
samme, som overalt, hvor man forvexler Ord og Begreber. Den Tryghed, man føler
ved det kendte Navn, fortolker man uvilkaarligt som Kendskab til Tingen, og man
undersøger ikke, om det kendte Navn ikke skulde være knyttet til en ukendt
Ting. Det er en hyppig Erfaring, at af to paa hinanden følgende sjælelige
Fænomener kan det følgende drage saa godt som hele Opmærksomheden til sig og
komme til at bestemme hele Tilstanden. Dog kan det ogsaa ske, at den levende
Genkendelse af andet Led (her Navnet) fører til, at man igen vender
Opmærksomheden mod første Led (hvilket forudsætter Evne til baglængs
Reproduktion), og derved vil da en \emph{tydeligere} Genkendelse af dette finde
Sted, end der indtraadte i det første Øjeblik. Herved kunde da hin daglige
Sædvane, at kalde en Ting kendt, naar Navnet kan nævnes, finde sin
Berettigelse. Meningen kan dog ikke være, at en Ting skulde være kendt, fordi
man kan fremplapre et Navn, man ingen Betydning forbinder med. Ved Navnet
vækkes jo Associationer om lignende Fænomener som det foreliggende:
% --- p. 35 ---
\setcounter{footnote}{0}%
\opage{35}Navnet staar --- naar det er mere end Ordskvalder --- som
Stedfortræder for en hel Række Lighedsassociationer. (Dette Punkt vil jeg ogsaa
i det Følgende faa Lejlighed til at vende tilbage til.)

4) Man maa ikke forvexle den større Lethed i at \emph{markere} Genkendelsen med
en større Lethed i \emph{selve} Genkendelsen. Naar Associationen mellem Nuance
og Navn er indøvet, er ogsaa Evnen til at markere indøvet, og det er da intet
Under, at Forsøg med benævnelige Nuancer vise større Sikkerhed end saadanne,
hvor en tilsvarende Indøvelse ikke var sket. Denne Lethed i Markeringen synes
at være den mest udprægede Ejendommelighed ved denne anden Række af de
Lehmannske Forsøg.

5) Lejlighedsvis omtales (p. 42) Vigtigheden af, at Iagttagerne anvendte «fuld
Opmærksomhed». Det vilde have været interessant at faa noget mere at vide om
denne Opmærksomhed. Var den vilkaarlig eller uvilkaarlig? Og hvad gik den ud
paa? Hvis den gik ud paa at give den opfattede Nuance det \emph{rette} Navn,
maa den antages at bestaa i en Besindelse paa, \emph{hvad} det nu var, man saa,
--- en \emph{Identifikation}. Og da der nu (ifølge p. 39---41) var gjort alt
for at holde Indflydelsen af \emph{forud} vakte (bevidste eller ubevidste)
Forestillinger borte, mon saa ikke den af Dr. L. med saa stor Umag, men med saa
liden Virkning bestridte umiddelbare Genkendelse (Assimilation) skulde stikke
Hovedet op, som Nissen hos Manden, der vilde flytte fra ham? --- Om
Opmærksomheden og dens Forhold til Genkendelsen har jeg talt i 5. ---

14. Som Afslutning af hele denne Undersøgelse om umiddelbar Genkendelse vil jeg
give en kort \emph{systematisk Oversigt over de forskellige Arter af
Genkendelse}, der i det foregaaende ere blevne omtalte.

Genkendelsens forskellige Arter bestemmes dels ved Beskaffenheden af
Genkendelsens \emph{Genstand}, dels ved den Maade, paa hvilken Genkendelsen er
\emph{forberedt}. Vi tage hver Inddelingsgrund for sig; det vil dog naturligvis
vise sig, at visse Genstande kun kunne genkendes efter en vis Forberedelse, saa
at et Led af den ene Inddeling kan falde sammen med et Led af den anden.

A. \emph{Genkendelsen karakteriseret} ved sin \emph{Genstand}.

1) Genkendelse af et forholdsvis simpelt, ikke meget sammensat Fænomen, hvis
enkelte Træk kunne opfattes simultant: \emph{øjeblikkelig Genkendelse}.
(Psykologi, V B, 1. Kfr. nærværende Afhandling 3---8.)

2) Genkendelse af mere sammensatte Fænomener, der ikke kunne opfattes
simultant: \emph{successiv Genkendelse}. (Psykologi V B, 2 og 4, samt C, 1---2
og C, 7 b og c. Kfr. nærværende Afhandling 9---10.)

3) Genkendelse af Elementer og Nuancer, der ikke pleje at forekomme særskilt i
% --- p. 36 ---
\setcounter{footnote}{0}%
\opage{36}vor Erfaring, og derfor kræve særlige Hjælpemidler for at genkendes:
\emph{Specialgenkendelse}. (Psykologi V B, 1, p. 141. Kfr. nærv. Afhdl. 13 c og e.)

B. \emph{Genkendelsen karakteriseret ved Forberedelsen}.

1) \emph{Uforberedt Genkendelse}. Her forudsættes kun en Disposition til
Forestillingens Genkomst, men ikke nødvendigvis nogen i nær Fortid foregaaet
Vækkelse af den. --- Uforberedt øjeblikkelig Genkendelse kaldes
\emph{umiddelbar Genkendelse}.

2) Genkendelse kan være \emph{forberedt ved vilkaarlig}, \emph{skønt forgæves}
Søgen. Smlgn. Psykologi p. 83 f. Nærv. Afhdl. 3.

3) \emph{Uvilkaarligt forberedt Genkendelse} forudsætter enten a) at
Forestillingen om det, der genkendes, kort forud har været fremme i
Bevidstheden, uden at vi selv have fremkaldt den, saaledes naar en
\emph{Forventning} er vakt hos os (Psykologi V B, 4. Nærv. Afhdl. 9---10), en
Forventning, som igen kan være \textit{α}) \emph{bevidst} eller \textit{β})
\emph{ubevidst}, --- eller, b) at Forestillinger, som hver for sig ikke ere
tilstrækkelige til at muliggøre en Genkendelse, ved deres \emph{opsummerede
Virkninger} muliggøre en saadan. (Nærv. Afhdl. 11.)

4) \emph{Vilkaarligt forberedt Genkendelse} kan enten a) blot bestaa i, at man
véd, at man skal prøve sin Genkendelsesevne for en vis Art Genstande
(\emph{ubestemt Forberedelse}), eller b) tillige ved, at man véd, hvilken
Forestilling man skal sammenligne det indtrædende Fænomen med (\emph{bestemt
Forberedelse}). Der er i begge Tilfælde en vis vilkaarlig Opmærksomhed til
Stede (smlgn. Psykologi VII A, 5. Kfr. nærv. Afhdl. 13 a). Specialgenkendelsen
forudsætter i Regelen en vilkaarlig Forberedelse. ---

De forskellige Arter af Forberedelse kunne naturligvis ikke holdes skarpt ude
fra hverandre. Saaledes nærmer den ved bevidst Forventning forberedte
Genkendelse sig i meget høj Grad til den vilkaarligt forberedte Genkendelse.

\begin{center}\large II. Berøringsassociationens Forudsætninger.\end{center}

15. Der er, som tidligere omtalt, gjort Forsøg paa at vise, at al
Forestillingsassociation egentligt er Berøringsassociation, d. v. s. kommer i
Stand ved, at Forestillinger, som hyppigt have været til Stede i Bevidstheden
samtidigt eller i umiddelbar Rækkefølge, have en Tendents til at genfremkalde
hinanden. Et saadant Reduktionsforsøg vil være tilstrækkeligt gendrevet, dersom
det kan vises, at \emph{selve Berøringsassociationen} hviler paa en
\emph{Forudsætning}, \emph{den ikke selv indeholder Forklaringen} af. --- Lad
\textit{A} og \textit{B} være to Fornemmelser (eller Fornemmelsesgrupper), og
lad os antage, at vi et vist Antal Gange have havt \textit{A} og umiddelbart
derefter \textit{B.} Følger saa deraf uden videre, at \textit{A'}s Genopstaaen
i Bevidstheden vækker Forestillingen \textit{b}?

% --- p. 37 ---
\setcounter{footnote}{0}%
\opage{37}Berøringsassociation er (ligesom den umiddelbare Genkendelse) et
specielt Tilfælde af Øvelse. Øvelsen bevirker i dette Tilfælde ikke blot, at
den enkelte Forestilling lettere vækkes; hvad der her er det afgørende, er, at
der ved Gentagelsen danner sig en saa nøje Sammenhæng mellem \textit{A}
og\textit{ B}, at de i Virkeligheden udgøre en og samme Proces. Man kan
(hvilket i et senere Afsnit nærmere skal udvikles) ikke forstaa
Forestillingsassociationens Fænomen, naar man ikke fastholder, at Fornemmelser
og Forestillinger ikke ere sondrede Dele eller absolut selvstændige Elementer,
men sjælelige Virksomheder eller Funktioner. Hvad Øvelsen (som her med Rette er
bleven kaldet Medøvelse) her bevirker, er, at der ikke bliver noget Spring
eller Mellemled mellem de to Virksomheder \textit{A} og \textit{B.} Naar den
ene af disse (\textit{A}) indtræder ved, at det tilsvarende Fænomen er givet
paa ny, vil den aandelige Virksomhed --- uden at behøve nogen særlig Impuls,
altsaa uden at Fænomenet \textit{B} behøver at indtræde paa ny --- fortsætte
sig videre, saa at vi faa \textit{A + b}. Her er altsaa det mærkelige sket, at
hvad der (for Bevidstheden selv) først var et blot \emph{Successionsforhold}
(at \textit{B} afløser \textit{A}) afføder et \emph{Kausalitetsforhold}, idet
\textit{b} fortsætter \textit{A} som Virkningen sin Aarsag.

Men lad os se lidt nærmere paa, hvad heri ligger. Dersom Gentagelse og Øvelse
have en saadan Indflydelse paa Forholdet mellem \textit{A} og\textit{ B,} maa
de da ikke ogsaa have Indflydelse paa \textit{A} alene? \textit{A} og
\textit{B} kunne dog ikke være helt forskellige Vilkaar undergivne. At
\textit{B} reproduceres, kunne vi kun forklare ved at antage, at der i
Bevidsthed og Hjerne ved hyppig Gentagelse er bleven en vis Disposition eller
Tendens tilbage, som kan udløses, uden at selve Fænomenet behøver at være
givet. Men \textit{A} er jo blevet gentaget lige saa ofte som \textit{B}!
Altsaa maa samme Disposition være til Stede for \textit{A'}s som for
\textit{B'}s Vedkommende. Og skulde det saa være muligt, at Fænomenet
\textit{A} kunde være givet, uden at denne Disposition vækkes? Det synes endog
at være en uundgaaelig Følgeslutning, at \emph{Dispositionen til
Forestillingen} \textit{a} maa \emph{vækkes} i \emph{langt højere Grad},
\emph{naar selve} \textit{A indtræder}, \emph{end Dispositionen} til \textit{b
vækkes}, \emph{fordi} det fra \textit{B forskellige} \textit{A indtræder}. Det
er ikke godt at indse, hvorledes denne Konsekvents kan undgaas. Og den
bekræftes desuden ved Erfaring. Det har nemlig vist sig ved Forsøg, at
ensartede Forestillinger lettere og hurtigere forbinde sig end uensartede,
eller med andre Ord, at Assimilationsprocessen gaar hurtigere for sig end
Komplikationsprocessen%
% footnote 1) on p. 37
\footnote[1]{\emph{Tschisch}: Ueber die Zeitverhältnisse der Apperception
einfacher und zusammengesetzter Vorstellungen (Wundt's Studien II. Leipzig
1885. p. 625). --- Smlgn. allerede \emph{Charles Bonnet}: Essai analytique, p.
399: «S'il faut moins de temps pour contracter les déterminations qui
constituent la simple \emph{réminiscence}, que pour contracter
\emph{l'habitude} de reproduire une \emph{suite} quelconque, c'est que la
reproduction de cette suite tient à de plus grands changements que la simple
réminiscence.»}. Forbindelsen af \textit{A} og\textit{ b} er en Komplikation,
Forbindelsen af \textit{A} og \textit{a} en Assimilation.

Dersom Dispositionen til \textit{a} ikke kan vækkes, saa kan Dispositionen til
\textit{b} heller ikke
% --- p. 38 ---
\setcounter{footnote}{0}%
\opage{38}vækkes. Altsaa Betingelsen for, at \textit{A} skal vække \textit{b,}
er, at det vækker \textit{a.} Dette har man udtrykt saaledes, at al Association
forudsætter en Lighedsassociation. Men derved maa da vel mærkes, at \textit{A}
og \textit{a} ikke \emph{behøve} at træde frem i Bevidstheden som selvstændige
Elementer. De kunne umiddelbart sammensmelte eller assimileres, og vi faa da en
umiddelbar Genkendelse som nødvendig Indledning til Berøringsassociationen. Da
man ofte har lagt stor Vægt paa Forskellen mellem Assimilation og Association
(smlgn. 7), er det rigtigere at udtrykke sig saaledes, som jeg har gjort i min
Psykologi (p. 180): «\emph{Berøringsassociation forudsætter
Lighedsassociation}, i det \emph{mindste umiddelbar Genkendelse}». At selve den
umiddelbare Genkendelse ikke kan forklares ved Berøringsassociation, har jeg
søgt at vise i det foregaaende Afsnit, hvor det ogsaa er omtalt, at
Genkendelsen kan foregaa saa hurtigt, at den kan siges at ske ubevidst. Paa
Grund af den Hurtighed, med hvilken den første Del af Processen
(Assimilationen) i Regelen foregaar i Forhold til den anden Del af Processen
(Komplikationen), er det let at forstaa, at der ofte ikke lægges Mærke til den.
Det er da for saa vidt intet Under, at der er Psykologer, der nægte, at den i
det hele taget foregaar.

16. Denne Opfattelse bekræftes ved bestemte Erfaringer. Ti hvor det første Led
i et Berøringsforhold ikke kan genkendes, indtræder heller ikke det andet Led.
Det er Modstykket til saadanne Tilfælde, i hvilke vi kunne kende det første
Led, men ikke genkalde os andet Led. (Se 3.) Jeg sidder i min Stue og hører
Rumlen paa Gaden; strax indtræder Forestillingen om en Omnibus. Men Betingelsen
herfor er, at jeg genkender netop \emph{denne} ejendommelige rumlende Lyd.
Forvexler jeg den med den Lyd, Rullen af tomme Tønder frembringer, eller med
Torden, kommer Forestillingen om Omnibus'en ikke. Hvis jeg føler mig usikker,
lytter jeg opmærksomt efter, og først den med \emph{Opmærksomhed og klar
Bevidsthed foretagne Identifikation} af den nærværende Fornemmelse (\textit{A})
og Forestillingen om den fra tidligere Erfaringer bekendte Lyd (\textit{a}),
lader da Forestillingen om Omnibus'en (\textit{b}) fremtræde med Sikkerhed.
Berøringsassociationen beror altsaa her paa bevidst Identifikation.

Men \emph{denne bevidste Identifikation forudsætter en Lighedsassociation}.
Hvad skulde nemlig foranledige mig til at sammenligne den Lyd, jeg nu hører,
med Forestillingen om den Lyd, jeg før har hørt, uden netop Ligheden mellem
dem? (Smlgn. 13 b.) Hvor tydeligt jeg bliver mig Lighedsforholdet bevidst, gør
intet til Sagen; her kan tænkes alle mulige Grader lige ned til den
fuldstændige og umiddelbare Sammensmeltning, som vi kalde umiddelbar
Genkendelse, og som vi allerede (7) have betegnet som en Grænseform for
Lighedsassociation. --- En Mellemform mellem den umiddelbare Genkendelse og den
bevidste Identifikation have vi, hvor et Fænomen vækker vor Opmærksomhed, og
hvor --- uden at der sker en selvstændig Reproduktion --- den blotte
opmærksomme Fastholden af
% --- p. 39 ---
\setcounter{footnote}{0}%
\opage{39}Indtrykket og Fordybelse i det, er nok til, at Berøringsassociationen
indtræder. Saaledes naar vi ville fremkalde et Navn i vor Hukommelse. Vi
forestille os da Tingen saa levende som muligt med Bortvenden af Opmærksomheden
fra alle andre Forestillinger; tilsidst kan da det søgte Navn springe frem.
Vandet strømmer ikke ud, før Stigboret er taget op, og dette kan koste stor
Anstrengelse.

Genkendelsen kan udeblive ved første Led, ikke blot naar der ikke er
tilstrækkelig Disposition til Vækkelse af \textit{a} (tilstrækkelig
Forberedthed) tilstede, men ogsaa naar Opmærksomheden afledes. Selv om
Forbindelsen mellem \textit{A} og \textit{B} er nok saa indøvet, kan \textit{A}
dog ved at forekomme i en Sammenhæng, i hvilket det ikke formaar at vække nogen
Interesse, forblive ugenkendt. Saaledes f. Ex. naar \textit{A} blot er
Gennemgangsled mellem de to hver for sig højst interesserende Iagttagelser
\textit{X} og \textit{Y. b} vil i dette Tilfælde ikke fremkaldes, skønt alle
Betingelser maatte være til Stede, dersom de havde Ret, der paastaa, at
Berøringsassociationen mellem \textit{A} og \textit{B} kun har to Betingelser:
1) \textit{A}'s og \textit{B'}s hyppige Samværen, og 2) \textit{A'}s Genindtræden%
% footnote 1) on p. 39
\footnote[1]{Smlgn. A. \emph{Lehmann}: Om Genkendelse p. 9.}. Begge Betingelser
ere i et Tilfælde som det beskrevne til Stede, uden at Associationen virker. De
stærkt følelsesbetonede Omgivelser have her umuliggjort den Dvælen af
Opmærksomheden ved \textit{A,} som betinger Berøringsassociationen i alle
saadanne Tilfælde, hvor Genkendelsen ikke er gaaet over til en sekundært
ubevidst Virksomhed. Det, at en vis Tid og Kraft maa anvendes ved
Berøringsassociationens første Led, tyder paa, at der her foregaar Processer af
mere sammensat Natur, end det ved første Øjekast kunde synes.

17. Dr. Lehmann anvender, for at godtgøre Urimeligheden af at antage flere end
de to omtalte Betingelser for Berøringsassociationen, en Sammenligning, som
netop kunde have ført ham til at faa Øje paa den rette Opfattelse. Idet han
bestrider, at Ligheden mellem det nu givne \textit{A} og det tidligere givne
\textit{A} skal spille nogen Rolle, naar \textit{A} ved Berøringsassociation
fremkalder \textit{b,} siger han (p. 9): «Forholdet han %
% PRINTED AS IS: „han“ for „kan“ (p. 39)
opfattes analogt med ethvert Kausalsammenhæng paa fysisk Omraade. Naar f. Ex.
en Lakstang befinder sig tilstrækkeligt nær ved en Stump Papir, og Lakstangen
derpaa gnides, saa flyver Papiret hen til Stangen. Denne Bevægelse indtræder,
hver Gang to Betingelser ere opfyldte, nemlig at Lakket er gnedet, og det ikke
befinder sig for langt fra Papiret. Paa ganske tilsvarende Maade kan man tænke
sig, at \textit{A} fremkalder \textit{b'}s Opdukken i Bevidstheden, naar de to
Betingelser ere opfyldte, at \textit{A} tidligere har været sammen med
\textit{B,} og at \textit{A} paany dukker op.» Hertil bemærker jeg følgende.
Til Lakstangens tilstrækkelige Nærhed af Papiret svarer den tilstrækkelige
Gentagelse af \textit{A} og \textit{B} i Forbindelse med hinanden. Men hvis
\textit{A} ikke kan genkendes, er det netop et Tegn paa, at Gentagelsen ikke
har været tilstrækkeligt
% --- p. 40 ---
\setcounter{footnote}{0}%
\opage{40}hyppig, eller at andre Omstændigheder i Øjeblikket hindre den i at
gøre sin Indflydelse gældende. (Smlgn. 15---16.) Hvis \textit{A} ikke træffer
paa en ved tilstrækkelig Gentagelse frembragt Disposition til Reproduktion af
\textit{A,} saa vil \textit{B} heller ikke blive reproduceret. Lige saa lidt
som en gneden Lakstang paa Mars tiltrækker et Stykke Papir paa Jorden, lige saa
lidt kan jeg ved nu at se en Mand, jeg for 40 Aar siden har set, komme til at
tænke paa hans Kone, som han den Gang havde under Armen. Og ligesom der er et
bestemt Punkt mellem Mars og Papirsstykket, hvor Tiltrækningen begynder at
virke, saaledes er der en bestemt Grad af Indøvelse, ved hvilken Synet af
Manden vil vække Billedet af hans Kone. Som hist Afstandsforholdet, saaledes
gaar her Successionsforholdet over til et Kausalitetsforhold. Men hvad Dr.
Lehmann har overset, er, at denne bestemte Grads Tilstedeværelse uundgaaeligt
ogsaa medfører en Genkendelsesakt, og at hvad der udelukker Genkendelsen derfor
ogsaa udelukker Berøringsassociationen. Her glipper Analogien i det anførte
Billede, og dette beviser derfor intet.

18. De engelske Psykologer, der (efter \emph{William Hamiltons} Mønster) have
fremhævet Lighedsassociationen som nødvendig Betingelse for
Berøringsassociationens Indtræden, have undertiden udtrykt sig paa en Maade,
der kunde lede til Misforstaaelse. De have saaledes sagt: Det er jo ikke det
\textit{A,} vi nu savne, %
% PRINTED AS IS: „savne“, probably for „sanse“ (p. 40)
som har været sammen med \textit{B,} men et tidligere \textit{A}; altsaa maa
det nu givne \textit{A} reproducere det tidligere \textit{A,} og det er da
først dette \textit{A,} der fremkalder \textit{b. Stuart Mill} f. Ex. siger (i
hans Note 35 i første Bind af hans Udgave af James Mill's «Analysis of the
human Mind»): «My present sensation could not remind me of those former
sensations unlike itself, unless by first reminding me of the sensations like
itself, which really did coexist with them». En Kritiker, der vil underlægge
den Teori, han bestrider, den mest bogstavelige og udvortes Opfattelse, vil her
ved at hænge sig i Ordene og overse, at alle psykologiske Betegnelser ere
Metaforer, have et let Spil. «Man tænke sig: Stuart Mill mener, at jeg, naar
jeg møder en Mand alene, som jeg tidligere altid har set sammen med hans Kone,
først maa tænke paa alle de tidligere Gange, jeg har set ham, før jeg kan komme
til at tænke paa hans Kone!» Før man tillægger en Mand som Mill en saadan
Mening, som, det er let at vise, er nærmest absurd, skulde man dog helst læse
ham nøjagtigt og ikke rive enkelte Udtalelser ud af Sammenhængen. Læser man
nemlig videre, finder man, at Mill udtrykkeligt betegner hin «Erindring om de
tidligere Fornemmelser» som en Forudsætning for Genkendelsen, en Forudsætning,
der ikke selv behøver at træde frem i Bevidstheden. Han siger: «This is
\emph{implied} in what we call recognizing a sensation as one which we have had
before».

Hvad man savner hos de engelske Psykologer paa dette Punkt, er for det første
en tilstrækkelig skarp Distinktion mellem Kendsgerning og Teori. Naar
Genkendelse forklares ved Reproduktion, saa er det (se 4), fordi den kun kan
forstaas ved den Antagelse.
% --- p. 41 ---
\setcounter{footnote}{0}%
\opage{41}at der ved den virker det samme som ved selvstændig og fri Erindring,
kun paa noget anden Maade. Men der maa skelnes mellem Genkendelse og den frie
Lighedsassociation, selv om de kunne tjene til at belyse hinanden, og selv om
hin kan opfattes som et Grænsetilfælde af denne. For det andet lægge de
engelske Forskere ikke tilstrækkelig Vægt paa, at det er Gentagelse og Øvelse,
der frembringe Bekendthedskvaliteten, som jo er det Faktum, der skal forklares,
og at disse Aarsager ved fortsat Virken maa føre til, at Genkendelsesakten
nærmer sig Bevidsthedstærskelen, ja overskrider denne og i hvert Tilfælde
vanskeligt bliver Genstand for Opmærksomhed, da det nye eller mindre bekendte
--- naar der ikke er forbunden en særlig Interesse ved Genkendelsen --- lettere
tildrager sig Opmærksomheden. Jo oftere \textit{A} har staaet som blot
Indledning til \textit{B,} des hurtigere vil der blive gaaet hen over det, og
dets Indtræden kan forblive aldeles ubemærket%
% footnote 1) on p. 41
\footnote[1]{R. \emph{Wahle} (Beschreibung und Eintheilung der
Ideenassociationen. Vierteljahrsschrift für wissenschaftliche Philosophie IX.
p. 429) bebrejder ogsaa Mill, at han ikke ytrer sig bestemt om, hvorvidt
Reproduktionen af \textit{A} fremtræder med \emph{fuld} Klarhed. For sit eget
Vedkommende tilføjer han: «Es scheint mir oft, dass es in ausserordentlich
schwacher Nuance doch erscheine, und dass es nur, weil 1) der Glanz des eben
präsenten, dem gleichen Elementes dasselbe verdunkelt, und weil 2) die
Aufmerksamkeit auf die eintretende sollicitirte Vorstellung abgelenkt wird,
übersehen oder schwerer entdeckt wird... Ob \emph{immer}, wage ich nicht zu
behaupten. Und vielleicht ist es nicht ganz ungünstig für den Fortschritt der
Philosophie, wenn ihre Bearbeiter es wagen, nicht zu behaupten.» Disse
Bemærkninger kan jeg efter min Erfaring ganske tiltræde.}. Det er i det hele en
Lov, at hvor en Bevidsthedsproces kan spares, der spares den. De med Bevidsthed
forbundne Funktioner staa i Livets Økonomi som Omveje, der maatte befares for
at naa Maalet; men ved fortsat Virksomhed og Øvelse bliver det muligt at naa
Maalet ad kortere Veje. Derved spares da den Tid og Kraft, enhver
Bevidsthedsvirksomhed lægger Beslag paa. Den fysiologiske Tid forkortes, og
Energien kan koncentreres paa nye Opgaver. Og selv om Bevidsthedens Tærskel
ikke overskrides, sker Overgangen fra et Bevidsthedselement til et lignende dog
hurtigere end Overgangen fra et Bevidsthedselement til et derfra forskelligt.
Assimilation foregaar hurtigere end Komplikation. (Se 15.) Naar vi overraskes
af levende Forestillinger, hvis Aarsager vi først ved udtrykkelig Efterforsken
finde i Indtryk, vi slet ikke have lagt Mærke til, maa det antages, at en
saadan hurtig eller ubevidst Assimilationsproces har dannet Overgangen%
% footnote 2) on p. 41
\footnote[2]{Smlgn. min Psykologi p. 181.}.

\begin{center}\large III. Den frie Lighedsassociation.\end{center}

19. Ved fri Lighedsassociation forstaar jeg en saadan Forbindelse mellem
selvstændige, hver for sig klart fremtrædende Forestillinger, som kun kommer i
Stand paa Grund af deres Lighed eller Slægtskab. Den umiddelbare Genkendelse
kan, som flere
% --- p. 42 ---
\setcounter{footnote}{0}%
\opage{42}Gange bemærket, opfattes som en Art Lighedsassociation, men den er
ikke fri Association, fordi Leddene ikke gøre sig selvstændigt gældende i
Bevidstheden. Ved den umiddelbare Genkendelse behandles Forestillingerne (eller
Fornemmelsen og Forestillingen), som om de vare aldeles identiske. Derfor er
det ikke sagt, at de ere det. Der maa skelnes mellem psykologisk Identitet og
logisk Identitet. Til logisk Identitet udfordres, at jeg i enhver Sammenhæng, i
hvilken en af de vedkommende Forestillinger indgaar, kan sætte den anden i
Stedet for den, uden at det gør nogen Forskel. Den logiske Identitet har
\emph{Leibniz} træffende udtrykt paa følgende Maade: Eadem sunt, qvorum unum
potest substitui alteri salva veritate%
% footnote 1) on p. 42
\footnote[1]{Non inelegans specimen demonstrandi in abstractis. (Opera
philosophica ed. Erdmann p. 94.)}. Psykologisk kommer det blot an paa, om
Forestillingerne kunne sammensmelte, assimileres; de behandles da som
identiske, skønt de objektivt set kunne være højst forskellige. Derfor gaar den
normale Perception saa ofte umærkeligt over i Sanseillusionen.

Hvad der gælder for den højeste Grad af Lighedsassociation, den, ved hvilken to
Elementer fuldstændigt kunne dække hinanden psykologisk, og som derfor kan
siges at bero paa \emph{Dækningslighed}, gælder ikke mindre for de fjernere
Lighedsforhold, som jeg til Forskel fra Ensheden eller Dækningsligheden kalder
Ensartethed og Analogi. Her findes der efter Sagens Natur endnu flere
Forskelselementer ved Siden af de beslægtede Elementer. Jeg plejer derfor at
udtrykke Lighedsassociationen (i alle dens Grader) ved Formlen
\textit{a}\textsubscript{1} \textit{ + a}\textsubscript{2}, idet jeg ved de
forskellige Indices vil betegne (ikke den Orden, i hvilken Forestillingerne ere
opstaaede, men) den Omstændighed, at der foruden Ligheder ogsaa stedse gør sig
Forskelselementer gældende, som kunne paaagtes af Bevidstheden i større eller
mindre Grad.

Ved Undersøgelsen af, hvorvidt der gives fri Lighedsassociation%
% footnote 2) on p. 42
\footnote[2]{Det benegtes bl. a. af \emph{James Mill} (Analysis of the human
Mind. Chap. 3), \emph{Lotze} (Drei Bücher der Metaphysik. Leipzig 1879. p. 526;
anderledes tidligere i Mikrokosmus. 2. Udg. I. p. 243) og Kroman (Kortfattet
Tænke- og Sjælelære. 2. Udg. p. 143), uden at dog nogen af disse Forfattere
nærmere undersøger Spørgsmaalet. --- Lotze synes (efter en Ytring anf. Skr. p.
527) at antage umiddelbar Genkendelse, skønt han ikke antager fri
Lighedsassociation.}, vil det være Ensartetheden og Analogien, vi maa holde os
til. Nogle af de Forskere, som bestride, at der gives fri Lighedsassociation,
anerkende nemlig, at der gives umiddelbar Genkendelse (Assimilation). (Smlgn. 7.)

Association formedelst Ensartethed antager jeg finder Sted, naar f. Ex. en
Farve fremkalder Forestillingen om en beslægtet Farve, en Form Forestillingen
om en lignende Form, et Portræt Forestillingen om Personen, det fremstiller,
Forestillingen om en Mand Forestillingen om en Mand med en lignende Karakter o.
s. v. Ensartethed er \emph{Kvalitetslighed}. Den adskiller sig fra
Dækningsligheden ved, at en Kongruents af Forestillingerne ikke er mulig; der
maa nødvendigvis finde en Overgang Sted fra den ene Forestilling til
% --- p. 43 ---
\setcounter{footnote}{0}%
\opage{43}den anden, naar Ligheden skal knytte Forbindelsen mellem dem. (Det
tyske Sprog er her saa heldigt at have to simple Udtryk, idet «Gleichheit»
betyder Dækningslighed, «Aehnlichkeit» Kvalitetslighed.) Ved Dækningslighed
behøver der ikke at finde en Overgang Sted; kun naar det kræver særlig
Opmærksomhed og Undersøgelse at konstatere Dækningsligheden (Identiteten), faa
vi en successiv Proces. (Smlgn. 10. 16.) --- Association formedelst Analogi
antager jeg finder Sted, hvor ikke en enkelt Del eller Egenskab, men Forholdet
mellem flere Dele eller Egenskaber frembyder Lighed og derved volder en
Forestillingsforbindelse. Exempler herpaa haves i digteriske Billeder og
Sammenstillinger, i Sprogmetaforer og visse Symboler. Den Lighed, som her
virker, kunne vi (til Forskel fra Dækningsligheden og Kvalitetsligheden) kalde
\emph{Forholdsligheden}.

20. Det vil nu være min Opgave at prøve, om man paa nogen Maade kan undgaa at
anerkende Lighedsassociation som et umiddelbart og simpelt sjæleligt Forhold.
Jeg vil gennemgaa de forskellige Veje, ad hvilke man har forsøgt eller kunde
have forsøgt at reducere al Lighedsassociation til Berøringsassociation. Men
først maa jeg dog gøre den Bemærkning, at jeg ikke antager, og aldrig har
antaget, at der gives noget Tilfælde, hvor Lighedsassociationen virker rent og
ublandet. Ved selve Formuleringen af Lighedsassociationen:
\textit{a}\textsubscript{1} \textit{ + a}\textsubscript{2} har jeg villet
udtrykke, at ved Siden af Lighedsforholdet vil Berøringsforholdet stedse gøre
sig gældende. Vi ville i hvert Tilfælde ikke i vor Erfaring kunne paavise noget
Exempel, hvor det ene af de to Forhold virkede aldeles alene. Ligheds- og
Berøringsassociation, tilligemed en Mellemform (Association mellem Del og
Helhed), har jeg derfor opstillet som foreløbige Love, der ved nærmere
Betragtning vise sig at være forskellige Sider eller Grænseformer for én eneste
Grundlov: Bevidsthedens Tendents til Genfrembringelse af hele sin tidligere
Tilstand (Helhedsloven). Uagtet al Berøringsassociation forudsætter en
Genkendelse, taber den derfor ikke sin selvstændige Betydning. «Ved enhver
Forestillingsforbindelse», har jeg sagt, «virke to Love med hinanden... I
Erfaringen kan der neppe paavises noget Tilfælde, hvor ikke Lighed og Berøring
virke sammen.»%
% footnote 1) on p. 43
\footnote[1]{Psykologi. 2. Udg. p. 181---184. --- Dette har imidlertid ikke
forhindret, at Dr. Lehmann i sin Afhandling adresserer Kritiken af den
\emph{rene} Lighedsassociation særligt mod min Bog. Han nævner ikke med et
eneste Ord, at jeg som Resultat af min Undersøgelse hverken antager Ligheds-
eller Berøringsassociation som aldeles oprindelige og selvstændige Love, men
fører dem begge tilbage til Helhedsloven. --- Jeg ser her endda bort fra, at
han citerer mig paa en saadan Maade, at hans Læsere maa tro, at jeg vil gjøre
al Association til Lighedsassociation.} Det gaar her som ved alle vore
psykologiske Udtryk og Begreber: vi maa ikke forvexle vore Abstraktioner, som
vi maa foretage for at kunne \emph{tænke} over Sagen fra alle Sider, med lige
saa skarpe Forskelligheder i selve Bevidstheden. Hvad det gælder om i det
følgende er da at paavise, at der er sjælelige Fænomener, som blive uforstaaelige,
% --- p. 44 ---
\setcounter{footnote}{0}%
\opage{44}naar vi ikke antage, at Lighedsforhold kan være en væsentlig
afgørende og uundgaaelig Betingelse for Forestillingers Association.

Saa vidt jeg sér, frembyder der sig kun \emph{fire} Maader, paa hvilke en
Reduktion af Lighedsassociation til Berøringsassociation vil kunne gennemføres.

a) Man kunde opfatte Sagen saaledes: Naar to Fænomener ligne hinanden, vil det
sige, at de have nogle Elementer (Dele eller Egenskaber) fælles, medens andre
ere forskellige. Lad det ene være \textit{A}\textsubscript{1} \textit{ =
m+n+p,} det andet \textit{A}\textsubscript{2} \textit{ = x+y+p.} Naar nu
\textit{A}\textsubscript{1} formedelst Lighed synes at vække Forestillingen om
\textit{A}\textsubscript{2}, sker da Overgangen ikke i Virkeligheden gennem det
fælles Element \textit{p,} saa at vi gennemløbe Rækken
\textit{m}+\textit{n}+\textit{p}+\textit{y}+\textit{x,} og derved umærkeligt
komme over fra \textit{A}\textsubscript{1} til \textit{A}\textsubscript{2}? Og
i bekræftende Fald vilde det jo dog være lutter Berøringsforhold, der virkede.

b) For saa vidt \textit{p}\textsubscript{1} (\textit{p} som Element i
\textit{A}\textsubscript{1}) og \textit{p}\textsubscript{2} (\textit{p} som
Element i \textit{A}\textsubscript{2}) skulde være saa forskellige, at de ikke
kunde identificeres, idet der kun fandtes Kvalitetslighed, ikke Dækningslighed
mellem dem, kunde man antage, at der fandt en umærkelig Gradation Sted, idet
\textit{p}\textsubscript{1} umærkeligt ændrede sin Beskaffenhed, indtil Nuancen
\textit{p}\textsubscript{2} forekom. Vi vilde da faa en Række Nuancer som
Mellemled mellem \textit{p}\textsubscript{1} og \textit{p}\textsubscript{2}
\textit{,} men ingen springende Overgang fra en Forestilling til en anden lignende.

c) Vore Forestillinger ere i Regelen nøje forbundne med sproglige Udtryk, og
denne Forbindelse er en Berøringsassociation. Beslægtede Forestillinger knyttes
gerne til samme Ord. Kunde Associationen mellem \textit{a}\textsubscript{1} og
\textit{a}\textsubscript{2} da ikke forklares af, at
\textit{a}\textsubscript{1} fremkalder Ordforestillingen \textit{n} og denne
igen \textit{a}\textsubscript{2} \textit{,} saa at Ordet danner et Mellemled
mellem de to Forestillinger, og «Ligheden» slet intet har at sige her?

d) Endeligt kunde man mene, at der bag den tilsyneladende Lighedassociation %
% PRINTED AS IS: „Lighedassociation“ (p. 44; elsewhere „Lighedsassociation“)
skjuler sig en mere sammensat Proces. Vi sammenligne uvilkaarligt Fænomener,
som enten ere samtidigt givne, eller givne i umiddelbar Succession, eller som
ved Berøringsassociation fremtræde for vor Forestillen. De Fænomener, vi finde
lige, stille vi sammen, og som Følge af denne ved Sammenligning tilvejebragte
Sammenstilling associeres de da siden efter i vor Bevidsthed.

21. a) Det vil ved Kvalitetslighed og Forholdslighed vise sig umuligt i alle
Tilfælde at dele Fænomenerne hvert for sig i Bestanddele, af hvilke nogle ere
identiske, andre forskjellige. Naar vi finde de forskellige Farvenuancer mere
eller mindre beslægtede, f. Ex. Gult mere ligt Orange end det er ligt Rødt,
naar derfor Personer, som aldeles intet kende til Optik, alligevel ordne
Farverne paa samme Maade, som de forekomme i Spektret%
% footnote 1) on p. 44
\footnote[1]{\emph{Wundt}: Die Empfindung des Lichts und der Farben (Studien
IV). p. 345 f. --- Hermed stemmer det, at man er længere Tid om at skelne to
Farver, der ere Naboer i Spektret, end saadanne, som ligge hinanden fjernere.
\emph{Cattell}: Ueber die Trägheit der Netzhaut und des Sehcentrums (Studien
III). p. 105.}, eller
% --- p. 45 ---
\setcounter{footnote}{0}%
\opage{45}naar en Ting med en vis Farve vækker Forestillingen om en anden Ting
med en beslægtet Farve, saa kan intet identisk Element paavises. Gult og
Orange, eller de forskellige Nuancer af Rødt have intet fælles Element, og dog
minde de om hverandre. Naar Homer sammenligner Synet af de græske Krigeres
straalende Vaaben med Synet af en fjern Skovbrand, saa have dog de to
Fænomeners Farver, Metalglansen og Ildebranden, forskellig Kvalitet, og
Digterens Fantasi maa gøre et Spring for at komme fra den ene til den anden.
Naar et Portræt minder mig om Manden, det forestiller, saa vil denne
Association kunne gaa for sig, selv om Farvenuancerne i Portrætet ikke ere de
samme som i Virkeligheden, og en Forvexling intet Øjeblik vilde være mulig.
Naar det engelske Sprog kalder en Sommerfugl butterfly, eller naar det danske
Sprog kalder en bekendt Blomst Smørblomst, er der næppe psykologisk
Kvalitetsidentitet, men kun Kvalitetslighed mellem Smørret og det gule Insekt,
der har givet hele Gruppen Navn, eller den gule Blomst. Endnu tydeligere viser
dette sig, naar vi holde os til andre Egenskaber end Farven. Naar Sproget har
Udtrykket «et Blad Papir», mon man da kan tænke sig, at Egenskaben Tyndhed for
Bevidstheden har dannet et \emph{frit} Overgangsled mellem Forestillingen om
Bladet paa en Plante og Forestillingen om et Stykke Papir? \emph{Imellem} disse
to Forestillinger skulde der altsaa ligge en Forestilling om Tyndhed, vel at
mærke \emph{hverken} om Plantebladets \emph{eller} Papirets Tyndhed, men om
Tyndhed i og for sig! Eller naar et Æble foran mig paa Bordet faar mig til at
tænke paa et Æble paa et Kobberstik af Adam og Eva, skulde saa Billedet af
Æblet i min Bevidsthed gøre Plads for Forestillingen om Æblets blotte Form, der
atter afløstes af Forestillingen om Kobberstikkets Æble, der maaske ikke engang
havde ganske samme Form som det virkelige Æble? Man kan ikke hævde denne
Opfattelse uden at forny den gamle Abstraktionsteori, der lader vore
Forestillingers «fælles» Elementer være tilstrækkelige til at udgøre nye,
selvstændige Forestillinger. Og allermest urimeligt bliver dette, naar man
lægger Mærke til den store Rolle, saadanne Lighedsassociationer spille hos
Børn. Et 14 Maaneder gammelt Barn kaldte Varmeapparatet i en Jernbanekupé
«Tiktak», fordi det mindede om et Ur. Samme Barn kaldte, da det var 20 Maaneder
gammelt, en rød Gummiballon først for Maane, saa --- efter nogen Betænkning ---
for Lampe%
% footnote 1) on p. 45
\footnote[1]{Dette sidste Exempel er hentet fra \emph{Marty}
(Vierteljahrsschrift für wissenschaftliche Philosophie III. p. 324 f.), som
benytter det i en anden Sammenhæng.}.

Formligheden danner paa en Maade Overgang fra Kvalitetsligheden til
Forholdsligheden, idet Formen beror paa eller udtrykker et Forhold mellem
Delene. Naar en Digters Fantasi fra Synet af de dansende Par i Balsalen gaar
over til Forestillingen om Klodernes Rotation i Himmelrummet, vil man ikke
kunne paavise nogen Bestanddel, der kan danne det \emph{selvstændige}
Overgangsled. Heller ikke naar Ordet grøn overføres til at betegne et uerfarent
ungt Menneske. Eller ved de mange Metaforer, ved hvilke materielle Udtryk gaa
% --- p. 46 ---
\setcounter{footnote}{0}%
\opage{46}over til at betegne aandelige Forhold og Virksomheder. F. Ex.
Overvejelse (smlgn. ogsaa det franske penser, der egentligt er det samme som
peser). De to Forestillinger: Vægten med Skaaler og Lodder, der bevæge sig op
og ned, --- og de hverandre bekæmpende Drifter og Tilskyndelser i et Menneskes
Sind, have ingen fælles Bestanddele. Berøringsteorien maa antage, at
Bevidstheden stedse paa abstrakt Vis tænker sig tertium comparationis som
Mellemled. Det gør man, naar man skal forklare et Billede eller en
Sammenligning for dem, der ikke forstaa dem; men for andre er det unødvendigt,
og først og fremmest har da den ofte vidunderlige Fantasi, der kundgør sig i
Sprogets Metaforer%
% footnote 1) on p. 46
\footnote[1]{Smlgn. om disse: \emph{Whitney}: Life and Growth of Language.
London 1875. p. 80 ff. og især: A. \emph{Darmesteter}: La vie des mots. Paris
1887. --- Selv om Distinktionen mellem Metaforer og Metonymier ikke i alle
Tilfælde er aldeles sikkert gennemført, saa vil det dog være en Umulighed at
reducere \emph{alle} Metaforer til Metonymier.} og i de digteriske
Sammenligninger, ikke behøvet denne Skolemestermetode.

Det synes saaledes at stride fuldstændigt mod psykologisk Erfaring, at
beslægtede Forestillinger hver for sig altid skulde deles i identiske og
forskellige Elementer%
% footnote 2) on p. 46
\footnote[2]{For Høreforestillingernes Vedkommende henviser jeg til
\emph{Stumpf}: Tonpsychologie. Leipzig 1883. I. p. 115 ff, der paa dette Punkt
er kommen til samme Resultat som jeg. --- At en saadan Deling \emph{undertiden}
kan ske, negter jeg ikke (smlgn. 22).}. Jeg tror at kunne paastaa, at ogsaa
denne Antagelse er et Exempel paa, at man har forvexlet den logiske Betragtning
med den psykologiske. (Se 11 Slutn.) Logisk kunne vi opsummere en Række
Kendemærker, der udgøre det ene Begrebs Indhold, og en anden Række Kendemærker,
der udgøre det andet Begrebs Indhold, og naar vi sammenligne disse to Rækker,
kunne vi finde, at deres Led dels ere identiske, dels forskellige. Men logiske
Abstraktioner maa ikke forvexles med virkelige Forestillinger. Til de logiske
Distinktioner svare ikke lige saa mange virkelige Forestillinger i Bevidstheden.

Og selv om man \emph{altid} kunde forklare Ligheden ved, at de to
Forestillinger dels havde identiske (\textit{p})\textit{,} dels havde
forskellige Bestanddele (i den ene \textit{m+n,} i den anden
\textit{x+y})\textit{,} saa vilde derved dog ikke al Association være reduceret
til Berøringsassociation. Fordi jeg inden for den ene Gruppe (\textit{m+n+p})
er kommen til \textit{p,} deraf følger ingenlunde, at der fortsættes med
\textit{y + x.} Det sker (som ovenfor vist, se 15---16) kun, naar der er en vis
Erindring om, at vi have havt \textit{p} før. Hvis \textit{p} staar for mig som
noget aldeles nyt, vil det ikke ske.

22. b) Selv om saaledes ikke altid en enkelt Bestanddel, der hører til begge
Forestillinger, er det \emph{frie} Mellemled mellem dem, kunde man tænke sig,
at en vis ændret, udvisket og utydelig Form af den første Forestilling kunde
danne Mellemleddet, idet den fælles Bestanddel gennemløb en Række Nuancer, ved
hvilken den første Forestilling
% --- p. 47 ---
\setcounter{footnote}{0}%
\opage{47}tilsidst viste sig at være forvandlet til den anden. Æblet, jeg ser
foran mig paa Bordet, er rødt; Æblet paa Kobberstikket er graat; og jeg kan
ikke have nogen Forestilling om den æblerunde Form uden \emph{nogen} som helst
Farve. Men maaske Farven kunde udviskes og staa hen i det ubestemte, og denne
ubestemte Forestilling saa dannede et Overgangsled mellem Forestillingen om det
røde og Forestillingen om det graa Æble.

Vi have sikkert ubestemte og foreløbige Forestillinger, baade hvad Kvalitet og
Kvantitet angaar. \emph{Berkeley}, som med saa stor Iver bekæmpede den gamle
Abstraktionslære i Psykologien og hævdede Nødvendigheden af, at enhver
Forestilling maatte have en individuel Bestemthed, var utilbøjelig til at
indrømme dette. Naar En lover mig «noget Godt», saa opstaar der, mener
Berkeley, ikke nogen ubestemt Forestilling om gode Ting; men Ordene vække ---
uden at nogen som helst Forestilling opstaar --- umiddelbart en Følelse af
Glæde, ligesom en Trusel umiddelbar vækker en Følelse af Frygt%
% footnote 1) on p. 47
\footnote[1]{Principles of Knowledge. Introduction. § 20.}. Der vil dog
sikkert, før man er bleven aldeles vant til Ordenes Lyd og Betydning opstaa en
Forestilling og ikke blot en Følelse. Man vil ved Løftet tænke paa noget, som
har Lighed med, hvad man mest ønker, %
% PRINTED AS IS: „ønker“ for „ønsker“ (p. 47)
og ved Truselen paa noget i Lighed med hvad man mest frygter. Der vil i hvert
Tilfælde findes mange Mellemformer mellem den klart bevidste Forestilling om
Løftets eller Truselens Betydning og den rent instinktmæssige
Følelsesbevægelse, ved hvilken det lovende eller truende Ord kun virker som
udløsende Impuls, uden at dertil knyttes nogen Forestilling%
% footnote 2) on p. 47
\footnote[2]{I sin interessante Afhandling «On some Omissions of Introspective
Psychology» (Mind. 1884) henviser W. \emph{James} til de foreløbige
Forestillinger, vi danne, før vi endnu have faaet Svar paa et Spørgsmaal eller
have naaet en bestemt Opfattelse af en Sag.} (smlgn. 8).

Betragte vi nu saadanne ubestemte og foreløbige Forestillinger nærmere, maa de
deles i to Klasser, alt efter som den ubestemte Egenskab er af extensiv eller
af kvalitativ Art. Hvor det drejer sig om Størrelsesforhold, kan jeg have en
Forestilling om et mindre Omfang og ændre den ved simpel Udvidelse, eller en
Forestilling om et større Omfang og ændre den ved simpel Bortskæren. I begge
Tilfælde kan jeg skelne mellem fælles og forskellige Dele, og Ubestemtheden
bestaar blot i, at jeg ikke véd, hvor jeg skal sætte Grænsen. Hvor det derimod
drejer sig om Kvaliteter, kan jeg nok ogsaa vælge en foreløbig Kvalitet, men
hvis den ikke er den rette, eller hvis Tvivlen indfinder sig, maa den ved et
nyt Valg ombyttes med en anden Kvalitet, uden at der kan være Tale om fælles
Bestanddele. Jeg kan ikke have nogen Forestilling om Farve uden ved foreløbigt
at tænke mig en vis bestemt Farve, som saa senere maaske (naar Svaret kommer)
maa ændres ved mindre eller større Spring paa Farveskalaen.

En mørk Genstand, jeg ser i det Fjerne, kan være en livløs Genstand, et Dyr
eller et Menneske. Det er en ubestemt Forestilling, jeg har om den, men jeg har
virkeligt en
% --- p. 48 ---
\setcounter{footnote}{0}%
\opage{48}Forestilling. Naar den nu kommer nærmere, ændres dens Størrelse ved
Udvidelse; den bliver mere end en sort Prik. Men de andre Egenskaber, som gøre,
at jeg genkender den som et Dyr eller et Menneske, kan jeg ikke paa denne Maade
øge til; de komme til ved kvalitative Ændringer.

Hvad selve disse kvalitative Ændringer angaar, kunde de tænkes at foregaa
kontinuerligt. Det vilde da egentligt slet ikke blive noget
Associationsfænomen, lige saa lidt som Efterbilledernes Varieren er
Association. Æblets røde Farve kunde tænkes at blive mindre og mindre mættet,
indtil vi kom til den graa Farve. Hvis man vilde paastaa, at der alligevel
fandt en Association Sted, idet Forestillingen om en Farvenuance fremkaldte
Forestillingen om en mindre mættet Nuance o. s. v., saa vilde dette (i det
mindste hos dem, der ikke vare fortrolige med Nuancernes Skala) ikke være nogen
Berøringsassociation, men en Association formedelst Kvalitetslighed. Der er
ingen Grund til at antage, at Nuancer eller Former, som have stor Lighed med
hverandre, derfor skulde forekomme hyppigt sammen.

Foregaa Ændringerne ikke kontinuerligt, kan det dog heller ikke være nogen
Berøringsassociation, der finder Sted. Ved hver ny Kvalitetsovergang maa det
være en beslægtet Kvalitet, der afløser den forestillede. Slægtskabet er det
eneste Motiv, som kunde tænkes for en saadan kvalitativ Varieren af Forestillingerne.

23. c) Den Vej, ad hvilken man hyppigst forsøger at reducere
Lighedsassociationen til Berøringsassociation, er sikkert den tredie af de
ovenanførte. Man mener, at Ordet, den fælles Betegnelse, tjener som Mellemled
mellem to beslægtede Forestillinger. Denne Overgang synes endog at ligge saa
nær for nogle, at naar de f. Ex. finde Overgangen fra Forestillingen om
Napoleon til Forestillingen om Alexander anført som Exempel paa
Lighedsassociation, gaa de uden videre ud fra, at Ordet Feltherre har dannet
Mellemled. Ja, man er i den Grad hildet i Antagelsen om Ordforestillingernes
Fremhersken, at man begynder sin Konstruktion af den psykologiske Proces, som
antages at være foregaaet, med \emph{Navnet} Napoleon. Man tænker sig, at dette
Navn reproducerer Forestillingen om en stor Feltherre, der har vundet mange
Slag og erobret mange Lande, og denne Forestilling reproducerer saa igjen
Forestillingen om Alexander%
% footnote 1) on p. 48
\footnote[1]{A. \emph{Lehmann}: Om Genkendelse p. 8.}. Det kunde her være af
Interesse at faa at vide, om «Forestilling om en stor Feltherre» betyder
Forestilling om en stor Feltherre i al \emph{Almindelighed} eller blot
Forestilling om den store Feltherre, som Napoleon var. I første Tilfælde vilde
man vende tilbage til den gamle Abstraktionsteori og tillægge os Evne til at
have virkelige abstrakte Forestillinger, medens vel de fleste nu til Dags vilde
være enige om, at naar de skulde tænke sig en Feltherre, maa de altid
% --- p. 49 ---
\setcounter{footnote}{0}%
\opage{49}forestille sig en bestemt (historisk eller opdigtet) Person. Vælge vi
nu en eller anden (fra Napoleon forskellig) Person som psykologisk Repræsentant
for Begrebet «stor Feltherre», saa er Overgangen fra Forestillingen om Napoleon
til Forestillingen om denne Repræsentant jo akkurat af samme Art som Overgangen
fra Napoleon til Alexander, og intet er forklaret. Det gør (ifølge 22) intet
til Sagen, hvor ubestemt man tænker sig Mellemleddet. --- I andet Tilfælde ---
hvis Forestillingen om en stor Feltherre var individualiseret ved Napoleon selv
--- vilde vi jo ikke være komne af Stedet. --- Ad ingen af de to Veje kan
Forklaringen altsaa naas. Hvad selve det anvendte Exempel, som er hentet fra
min Psykologi, angaar, kan jeg bevidne, at jeg støttede mig til en
Selviagttagelse, ved hvilken Ordforestillingen ingen Rolle spillede. Jeg er
vant til at forestille mig historiske Personligheder i konkret Skikkelse
(Napoleon f. Ex. i den bekendte grønne Uniform), og utydelige Billeder af
Handlinger og Tildragelser knytte sig til dem. Overgangen skete fra et Billede
til et andet.

Der er store individuelle Forskelligheder med Hensyn til, hvor stor en Rolle
Ordforestillingerne spille i de forskellige Individers Forestillingsliv. Der er
dem, hvis Erindring væsentligt er en Sagerindring, og dem, hvis Erindring
væsentligt er en Orderindring%
% footnote 1) on p. 49
\footnote[1]{Themistocles rerum, Hortensius verborum recordatione dicuntur
valuisse. L. \emph{Vives}: De anima. Brugis 1538. p. 56. --- Themistocles var
praktisk Statsmand, Hortensius var Taler!}. Og hos et og samme Individ spille
Ordforestillingerne under forskellige Forhold en forskellig Rolle. Naar man
lever meget med Bøger og har meget literært Arbejde, eller naar man taler meget
med andre Mennesker, ville Ordforestillinger naturligvis blive hyppige og
fremtrædende og ville knytte sig som Bestanddele til de fleste andre
Forestillinger. Men jo mere man, f. Ex. i en længere Ferie, frigør sig fra
disse Beskæftigelser, des mere træder Ordforestillingerne tilbage, og
Erindringens og Fantasiens Billeder bevæge sig umiddelbart mellem hverandre. En
omhyggelig og selvstændig Iagttager, Patologen \emph{Stricker}, der selv lægger
megen Vægt paa Ordforestillingerne, udtaler sig saaledes herom: «Naar jeg
arbejder meget i literær Retning, naar jeg former Afhandlinger eller Foredrag,
koster det mig umiddelbart efter Arbejdet Umage blot i nogle Minuter at være
uden Ord- eller Toneforestillinger. Derimod lykkes det mig umiddelbart efter en
Spadseretur i det Frie, eller efter at have beset Kunstværker, let at lade mig
beherske af Erindringerne om de Skikkelser, jeg har set. Lige saa let kan jeg
umiddelbart efter et behageligt Bad eller efter andre behagelige Sanseindtryk
nogle Minutter igennem ganske hengive mig til Erindringen om \emph{disse}
Indtryk, uden at forestille et eneste Ord eller en eneste Tone»%
% footnote 2) on p. 49
\footnote[2]{Studien über die Sprachvorstellungen. Wien 1882. p. 2.}. Den
kunstneriske og digteriske Fantasi bevæger sig gennemgaaende i Billeder, ikke i
Ord, og des mere, jo stærkere Bevægelse den er i. Det kan endog være med en vis
Resignation, at den store Digter bestemmer sig til at iklæde sine
Forestillinger Ordets fattigere, abstrakte Form. \emph{Goethe} siger saaledes
(i en Samtale med Eckermann) om sine Ballader: «Jeg havde
% --- p. 50 ---
\setcounter{footnote}{0}%
\opage{50}allerede i flere Aar havt dem i Hovedet som yndige Billeder, skønne
Drømme... Jeg besluttede mig nødigt til at sige dem Farvel ved at give dem et
Legeme i det utilstrækkelige tarvelige Ord»!

Det er heller ikke for alle Forestillingers Vedkommende, at Associationen med
de tilsvarende Ord er lige fast og uundgaaelig. Ord, som betegne individuelle
Fænomener, glemmes lettere end Betegnelser for abstrakte Forhold. Højst
mangelfuld Navnehukommelse behøver ikke at være et patologisk Fænomen, men kan
hænge sammen med, at man har liden Brug for Ordene, idet man væsentligt holder
sig til de konkrete Fænomener. Man har vist med Rette forklaret konkrete
Benævnelsers lettere Forglemmelse ved, at vi ved konkrete Fænomener ikke behøve
Ordet, da Erindrings- eller Fantasibilledet er os nok, medens abstrakte Forhold
vanskeligere fastholdes uden ved Ordtegnets Hjælp. Deraf følger, at Benævnelser
paa Abstraktioner ere ganske anderledes indøvede end Benævnelser paa konkrete
Fænomener%
% footnote 1) on p. 50
\footnote[1]{\emph{Kussmaul}: Die Störungen der Sprache. Leipzig 1877. p. 164.}.

Efter den nu almindelige Opfattelse af Sprogets Opstaaen og Udvikling%
% footnote 2) on p. 50
\footnote[2]{\emph{Madvig}: Om Sprogets Væsen, Udvikling og Liv. Kbhvn. 1843.
p. 5. \emph{Whitney}: Life and Growth of Language. London 1875. p 278---309.}
skyldes det væsentligt Trangen til Meddelelse. Først efter at det er udviklet
som Middel til Meddelelse, bliver det sekundært for den Enkelte Middel til
Fastholdelse og Skelnen af Forestillinger. Det er med Urette, naar især ældre
Tænkere vendte Forholdet om og lod Ordet være Mærke (nota) for den Enkelte i
hans Bevidsthed, før det blev et Tegn (signum), ved hvilket han gjorde sig
forstaaelig for andre%
% footnote 3) on p. 50
\footnote[3]{Saaledes \emph{Hobbes}: Logica. II, 3.}. Vi kunne --- som baade
Filologer og Patologer indskærpe%
% footnote 4) on p. 50
\footnote[4]{A. \emph{Darmesteter}: La vie des mots. p. 37. ---
\emph{Kussmaul}: Die Störungen der Sprache. p. 17.} --- meget vel have
Forestillinger uden Ord. Ovenfor er der i en anden Sammenhæng (12) allerede
gjort Brug af den Kendsgerning, at vi kunne genkende et Fænomen, uden at
erindre dets Navn; hvad der gælder for Genkendelse (bunden Erindring), gælder
ogsaa for den frie Erindring. Ordbetydningernes Udvikling viser desuden
tydeligt nok, at Forestillingsvirksomheden og Forestillingsforbindelsen kan ile
forud for Ordanvendelsen. Naar et Ord overføres fra at være Betegnelse for et
Fænomen til at være Betegnelse for et lignende Fænomen, forudsætter denne
Overgang en Forbindelse mellem Forestillingerne om de to Fænomener, en
Forbindelse, der kun kan være tilvejebragt ved Ligheden mellem dem. Ordet
Feltherre er dannet for at betegne saadanne Mennesker som Napoleon og
Alexander; selv om man vilde sige, at det er Ordforestillingen Feltherre, der
nu knytter Forestillingen om Napoleon og Forestillingen om Alexander sammen,
saa forudsætter dog denne Berøringsassociation en forudgaaende, primær
Bevidsthedsvirksomhed, ved hvilken Napoleon og Alexander (eller lignende Mænd)
første Gang sammenstilledes.
% --- p. 51 ---
\setcounter{footnote}{0}%
\opage{51}Hvorfra skulde ellers det fælles Ord skrive sig? --- Ved den
videnskabelige Klassifikation søger man at indrette Benævnelserne saaledes, at
de henpege paa den Plads, Fænomenet indtager i Forhold til andre Fænomener
efter sin Lighed eller sin Forskel overfor dem. Hvad der i Videnskaben sker med
fuld Bevidsthed og paa Grundlag af gennemført Analyse, det foretages paa et
tidligere Udviklingstrin i de grove Træk uvilkaarligt paa Grundlag af
Lighedsassociationer.

Imidlertid skal det ingenlunde negtes, at Ordforestillinger --- hvilke saa
disses nærmere psykologiske Beskaffenhed maatte være%
% footnote 1) on p. 51
\footnote[1]{Ogsaa her gøre individuelle Forskelligheder sig gældende. Hos
nogle (de fleste) er det lydende, hos andre det synlige Ord, hos andre igen
Ordet som Bevægelsesforestilling herskende. \emph{Stricker} har i sine «Studien
über die Sprachvorstellungen» lagt saa megen Vægt paa den rent motoriske Side
ved Ordforestillingerne, at en Kritiker i «Revue philosophique» (1886. II. p.
335) har formodet en Abnormitet hos ham: «Le cas de M. Stricker est peut-être
pathologique. Il y a peut-être des sourds chez ses ascendants. Son procédé
intellectuel est-il un fait d'atavisme?» I Stedet for at spørge, hvem der er
normal, og hvem der er abnorm, er det rigtigere at anerkende de individuelle
Forskelligheder. --- Smlgn. om dette Spørgsmaal: \emph{Stumpf}: Tonpsychologie.
I. p. 153 ff. \emph{Ribot}: Psychologie de l'Attention. Paris 1889. p. 84 ff.}
--- spille en overordentligt stor Rolle. Ordet bliver Meddelelsesmiddel, fordi
det selv fra først af er et naturligt Udslag af den Virkning, Fænomenet gør paa
os: Lyd eller Gestus ere Følger af den Følelsesbevægelse, vore Erfaringer
medføre. Sanseindtryk, Lyst- eller Ulystfølelse og lydeligt eller synligt
Udbrud, det er den primitive Række af sjælelige Elementer, som ethvert for
Individets Existents betydningsfuldt Fænomen vil fremkalde. Denne Række
forkortes efterhaanden, idet Udbruddet kan holdes tilbage, naar det ikke netop
skal bruges til Meddelelsesmiddel. Vore blotte Fornemmelser og Forestillinger
ere for saa vidt Kunstprodukter, som de oprindeligt kun vare Overgangsled til
Udbrud (eller Handling). En Hæmningsproces er Betingelsen for, at Fornemmelser
og Forestillinger kunne blive Hovedsagen ved Individets Forhold til de
Genstande, det staar overfor. Men skønt den primitive Sammenhæng mellem
Forestilling og Udbrud%
% footnote 2) on p. 51
\footnote[2]{Smlgn. Psykologi. p. 179 og 197.} saaledes hæmmes, forstærkes den
paa den anden Side ved Meddelelsestrangen, som bevirker, at Ordforestillingerne
ere vigtige Elementer ved saare mange af vore Forestillingsforbindelser, og ved
Trangen til at fæstne Forestillingerne om Tingenes fælles Egenskaber. Meget
ofte vil der derfor, foruden Lighedsassociationen ud fra Forestillingen om
Tingen, foregaa en Berøringsassociation ud fra Navnet paa Tingen. Som ovenfor
omtalt findes der neppe nogen ren Lighedsassociation, men begge Motiver virke
sammen. De kunne føre i samme, men ogsaa i forskellig Retning. Naar den blotte
Ordassociation faar Overvægt, og ikke kontrolleres af Lighedsassociationen, vil
det være et Tegn paa Tankeløshed, om ikke paa Sindssygdom.

Det kan i de enkelte Tilfælde være vanskeligt at afgøre, hvormeget der skyldes
Lighedsassociation ud fra den anskuelige Hovedforestilling, og hvor meget der skyldes
% --- p. 52 ---
\setcounter{footnote}{0}%
\opage{52}Berøringsassociation ud fra den simplere, mere mekaniske
Ordforestilling. --- Under et Besøg i en lille Provinsby gjorde jeg en Dag den
Bemærkning, at man sjeldent saa nogen gaa ind i en Butik. Næste Dag sad jeg
flere Mile derfra ved Kysten og saa Maagerne paa Fiskefangst. Jeg ytrede da til
min Ledsager, at man næsten aldrig saa nogen af dem faa fat i en Fisk. I samme
Øjeblik kom jeg til at tænke paa den lille Bys tomme Gader og ubesøgte Butiker,
og den karakteristiske Sammenstilling af de to Billeder (Maagerne, der ingen
Fisk fik, og Købmændene, der ingen Kunder fik) slog mig strax, uagtet det var
Ferietid og jeg ikke særligt beskæftigede mig med psykologiske Teorier. --- Jeg
er sikker paa, at Analogien mellem de to Fænomener spillede den største Rolle
ved Associationen, men kan ikke benegte, at jeg kan have brugt de samme Ord ved
begge Lejligheder. De i sig selv saa gængse Ord vilde dog ikke være
tilstrækkelige til at forklare Associationen; her kan kun være Tale om en
Medvirken. ---

Da der faktisk gives Forestillinger, som ikke behøve at være knyttede til Ord,
--- da samme Ords Anvendelse for beslægtede Fænomener forudsætter
Lighedsassociation mellem de tilsvarende Forestillinger, --- og da
Ordassociationen øjensynligt ofte kun er en medvirkende Faktor, --- saa vil det
være umuligt at forklare \emph{alle} Lighedsassociationer som
Berøringsassociationer, der komme i Stand ved Ordets Hjælp. ---

24. d) Hvad endeligt den fjerde Hypotese angaar%
% footnote 1) on p. 52
\footnote[1]{Den er bl. a. fremsat af \emph{Victor Brochard} i hans Afhandling:
La loi de similarité (Revue philosophique. 1880. I. p. 252 f.) og af
\emph{James Ward} i Art. «Psychology» i Encyclopaedia Britannica.
9\textsuperscript{th} ed. Vol. XX. p. 60 b.}, vil den kunne føre til en
Undersøgelse af Forholdet mellem Forestillingsassociation og Sammenligning.
Denne Undersøgelse vil jeg dog opsætte til næste Afsnit og her kun foreløbigt
bemærke, at der neppe er Grund til at antage nogen principiel Forskel mellem
den Bevidsthedsvirksomhed, der ytrer sig som Forestillingsassociation, og den,
der ytrer sig som fri Sammenligning. Derimod skal jeg her dvæle ved et andet
Punkt, som synes at være afgørende for denne Hypotese, og som tillige vil give
mig Anledning til at gaa noget nærmere ind paa Lighedsassociationens Psykologi
og Fysiologi.

Man har nemlig opkastet det Spørgsmaal: Hvad vil Lighed mellem en nærværende og
en ikke nærværende Forestilling egentligt sige? Hvorledes kan Ligheden gøre, at
en Forestilling dukker frem, naar Bevidstheden ikke forud har sammenlignet
denne Forestilling med den nærværende, som antages at fremkalde den?

Det er ganske vist en mærkelig Egenskab ved Bevidstheden, vi her staa overfor;
men er den egentligt mærkeligere end den, der lægger sig for Dagen i
Berøringsassociation og i al Erindring? Det er de højere Livsformers
Ejendommelighed, at de ere uafhængige af det enkelte Øjeblik og det enkelte
Sted og kunne udøve Funktioner, som pege langt ud
% --- p. 53 ---
\setcounter{footnote}{0}%
\opage{53}over Tidens og Rummets Begrænsning. I al Erindring og i al
Tankevirksomhed, i ethvert Instinkt og i enhver Drift lægger denne
Ejendommelighed, om man vil, denne Idealitet sig for Dagen. Ved
Berøringsassociationen forbinder Bevidstheden det, der er forskelligt i
Beskaffenhed, skønt det hører sammen i Tid og Rum. Denne «Sammenhøren» i Tid og
Rum ophæver dog ikke Forskelligheden mellem Forestillingerne; og det er kun en
Indbildning, naar man mener bedre at kunne forstaa, hvorledes \textit{a} kan
vække \textit{b}, end, hvorledes \textit{a}\textsubscript{1} kan vække
\textit{a}\textsubscript{2} \textit{.} Ved Lighedsassociationen er der Forskel
i Tid og Rum, men Beskaffenheden er identisk, ensartet eller analog; der
sammenstilles her ganske vist noget, som ikke før har været sammenstillet i
Bevidstheden; men skulde ikke Forestillingsindholdets Lighed kunne veje op mod
Nærheden i Tid og Rum? Ligesaa vel som man spørger om, hvorledes Bevidstheden
kan forestille sig noget, blot fordi det er ligt noget andet, naar den ikke i
Forvejen har sammenlignet begge Dele, ligesaa vel kunde man spørge, hvorledes
Bevidstheden kan forestille sig noget, blot fordi det tidligere er indtraadt
samtidigt med eller umiddelbart efter noget andet, naar den ikke først erindrer
sig denne tidligere Indtræden. Der er jo saare mange Tilfælde, hvor vi blive
overraskede over vore egne Berøringsassociationer og først ved nærmere
Undersøgelse blive klare over, at de ere Reproduktioner af noget, vi forhen
have oplevet; ofte maa man endda ty til vaagne eller virkelige Drømme for at
finde Forklaringen. --- Det vil i hvert Tilfælde være værdt at prøve, om man
ikke kunde vise, at det er lige saa «naturligt» at gaa over fra en Forestilling
til en lignende, som at gaa over fra en Forestilling til en derfra forskellig,
der er forekommet sammen med den.

Forestillingsassociationens Liv og Omfang beror paa den Energi, hvorover
Bevidstheden og Hjernen raade. Der kan være saa liden Energi til Stede, at
Associationsprocessen saa godt som stanser, eller at i hvert Tilfælde
Associationer, som ellers vilde staa til Rede, ikke ske. Ved Overanstrengelse
og Udmattelse kan man miste Kundskaber, man ellers disponerer over. Paa den
anden Side tager Forestillingsassociationen Fart, naar der er særligt megen
Energi til Raadighed, og naar den særligt let udløses. I Febertilstande, under
Indflydelse af Opium og i stærkt bevæget Stemning blive
Forestillingsassociationer og Genkendelsesakter mulige, som ellers ikke kunne
foregaa. I fysiologisk Henseende maa her mindes om den afgørende Betydning, som
Blodtilstrømningen og dens Svingninger have for Hjernens Virksomhed. Ifølge
Maudsley strømmer en Femtedel af hele Legemets Blodmængde gennem Hjernen, og
enhver Forminskelse eller Forøgelse faar strax Indflydelse paa Hjernevirksomheden.

Dersom der er rigelig Energi til Stede, vil der være en Trang til at anvende
den. Fjederen er stærkt spændt, og den aller ringeste Paavirkning er nok til at
faa den til at slaa ud. Ligesom Dyr og Mennesker i sund og kraftig Tilstand
uvilkaarligt bevæge sig blot for at faa Afløb for deres opsamlede Kraft, saa at
selv om der er ydre Anledning til
% --- p. 54 ---
\setcounter{footnote}{0}%
\opage{54}Bevægelse, denne ydre Aarsag i hvert Tilfælde er aldeles forsvindende
i Forhold til den indre Aarsag, saaledes er der ogsaa ofte en Trang til
Forestillingsvirksomhed, som ikke er rettet mod noget bestemt Maal, men som
fysiologisk ligeledes hænger sammen med Trangen til Udladning af den i Hjernen
opsparede Energi. Der kan tales om spontan Forestillingsvirksomhed, ligesom der
tales om spontane Bevægelser. Og ligesom den spontane Bevægelse naturligt slaar
ind paa de Retninger og Veje, Individet er vant til at gaa, saaledes vil der
ogsaa, hvad den spontane Forestillingsvirksomhed angaar, være visse
Forestillingsrækker, som særligt let reproduceres; det er den vante
Forestillingskreds, der uvilkaarligt indfinder sig, naar vi efter Søvnen igen
med saa stor Hurtighed orientere os i den vaagne Virkeligheds Verden. Den
\emph{spontane Virksomhed gaar} i den \emph{Retning}, \emph{hvor den mindste
Modstand mødes}, \emph{og hvor Udladningen altsaa lettest gaar} for sig.
Berøringsassociationen finder herved sin simple Forklaring. Men i dette samme
Princip ligger ogsaa Forklaringen af Lighedsassociationen. Dersom nemlig en vis
Forestilling eller Forestillingsgruppe er given, og dersom den omfattes med
stærk Interesse og altsaa søges fastholdt saa længe som muligt, samtidigt med
at den stærke Bevægelse i Bevidsthed og Hjerne søger Afløb, saa synes der her
at stilles to Opgaver eller at næres to Fornødenheder, der indbyrdes stride mod
hinanden: Interessen ved det forestillede fører til at holde det over
Bevidsthedens Tærskel og fordybe sig i det saa meget som muligt, men den stærke
Bevægelse fører til at gaa over til andre Forestillinger, hvortil kommer, at
(som allerede ovenfor vist, se 5.) Opmærksomheden ikke kan virke aldeles
kontinuerligt. Det gælder: paa én Gang saa meget som muligt at dvæle ved, eller
i hvert Tilfælde fjerne sig saa lidt som muligt fra den givne Forestilling, og
dog at lade Trangen til Forestillingsvirksomhed ske Fyldest%
% footnote 1) on p. 54
\footnote[1]{For at gøre Sagen mere indlysende har jeg tænkt mig et Tilfælde af
særligt rigelig Bevidstheds- eller Hjerneenergi. Men ogsaa under sædvanlige
Forhold er der saamegen Energi til Stede, og udløses den saa stadigt og
uundgaaeligt, at vi ikke formaa at stanse Forestillingernes Strøm i os, men maa
lade os nøje med ved Anspændelse af Opmærksomheden at danne Associationscentra,
om hvilke Strømmen kan kredse. Allerede \emph{Locke} bemærker: I would have any
one try whether he can keep one unvaried single idea in his mind, without any
other, for any considerable time together..... He will find it difficult to
keep all other ideas out of his mind, but that some, \emph{either of another
kind}, \emph{or various considerations of that idea} (each of which
considerations is a new idea) will constantly succeed one onother (Essay, II,
14, 13---14).}. Ved Lighedsassociationen tilfredsstilles begge Krav, og den
tilsyneladende modsigende Opgave løses. Bevidstheden kredser her uden om en
given Forestilling, idet den kun forlader den for at gaa over til ensartede
(kvalitetslige) eller analoge (forholdslige) Forestillinger. Trangen til
Forandring fyldestgøres; der komme \emph{andre} Forestillinger i Stedet for den
givne; men Forandringen bliver, paa Grund af Ensartetheden eller Analogien, den
\emph{mindst} mulige. Ogsaa her gælder det da, at Virksomhedstrangen fører i
den Retning, i hvilken der er mindst Modstand at overvinde.

% --- p. 55 ---
\setcounter{footnote}{0}%
\opage{55}Man vil ikke kunne forstaa den digteriske Fantasis Virken uden at
lægge en Betragtningsmaade som den fremførte til Grund. Digterens inderlige og
stærke Medfølelse fører ham, som episk og dramatisk Digter, til omhyggeligt og
trofast at forestille sig Genstanden i alle dens Enkeltheder og med alle dens
Træk. Men det er den samme Trang, der fører ham, som lyrisk Digter, til at
udtrykke sin Oplevelse og sin Tilstand ved at søge Paralleler og Analogier til
den. Han kan ikke udtrykke sin hele Stemning ved blot at dvæle ved selve
Forestillingen om Genstanden; og paa den anden Side vil han ikke slippe denne
Forestilling. Da gribes den Udvej at søge andre Forestillinger, der frembyde
lignende Egenskaber. Han betragter saa at sige hele Tilværelsen fra det
Synspunkt, den til hans Stemning knyttede Forestilling giver, og finder en
Genspejling af den overalt. I de homeriske Lignelser finder man
Tilfredsstillelsen for denne Trang til at lade den forestillede Tings
Ejendommelighed komme til sin fulde Ret --- netop ved at forestille sig en
anden Ting. Trangen til Symboler opstaar for en Del af denne Aarsag. (Kun for
en Del, ti mange Symboler have dannet sig rent historisk, altsaa ad
Berøringsassociationens Vej.) --- Lad mig tage et enkelt Exempel. Det hedder i
et gammelt lillerussisk Digt%
% footnote 1) on p. 55
\footnote[1]{Anført hos G. \emph{Brandes}: Indtryk fra Rusland. p. 277.}:

En Humleranke alene i Haven

slynged sig langs med Jorden.

En ung lille Pige blandt Menneskene

græd ganske bitterligt.

O grønne blomstrende Humle, sig,

hvorfor ranker du dig ikke opad?

O kære unge Pige, hvorfor

sørger du over din Skæbne?

Den grædende Pige og den langs Jorden krybende Ranke knyttes sammen, fordi
Stemningen trænger til mere end én Forestilling for at faa Luft, og fordi den
første Forestilling dog ikke ganske forlades, naar en analog Forestilling
træder i dens Sted. Sorgens Ubetvingelighed træder klarere frem, naar den
sammenstilles med Umuligheden for Humlen af at ranke sig opad uden Støtte. Der
virker her den samme Trang, som i en stærkt bevæget Stemning fører os til at
gentage samme Ord flere Gange. Forundring, Inderlighed og Følelse af Sagens
Vigtighed skaffe sig Luft paa denne Maade. (En stor, stor Mand. Jeg blev meget,
meget glad. Sandelig, sandelig jeg siger Eder.) Det hænger sammen med Følelsens
langsomme Udvikling og med dens Tendents til at brede sig i Bevidstheden%
% footnote 2) on p. 55
\footnote[2]{Smlgn. min «Formel Logik». 2. Udg. p. 27 f. og de der anførte
Afsnit af min Psykologi.}. Ogsaa Lighedsassociation er en Slags Gentagelse, en
Gentagelse med Variation, men en Variation, der ikke nøder til at ændre den
oprindelige Stemning.

% --- p. 56 ---
\setcounter{footnote}{0}%
\opage{56}Ogsaa ved den egentlige Tænkning spiller Lighedsassociation, Overgang
til beslægtede Forestillinger en væsentlig Rolle. Ved Behandlingen af et
Problem gaar Tænkeren ud fra en foreløbig Forestilling og søger at belyse den
ved saa mange Paralleler som muligt. En stor Opdagelse beror ofte paa en genial
Sammenstilling, som kun kommer i Stand formedelst Lighedsassociation. Forud for
den vilkaarlige og metodiske Gennemarbejden og Prøvelse af Forestillingernes
Forbindelse gaar der en uvilkaarlig Forestillingsvirksomhed, der tilfører de
Muligheder, som de bevidste Metoder kunne prøve. Uden en saadan uvilkaarlig
Proces som Grundlag, vilde Erkendelsen ikke komme af Stedet.

Videnskabernes Historie viser os f. Ex., hvorledes en uvilkaarlig
Klassifikation fører til «at benævne yderst \emph{lignende}, men i \emph{ingen
synlig Forbindelse staaende} Naturgenstande med samme Navn», og hvorledes
derved det første Skridt blev gjort til en videnskabelig Undersøgelse af
Fænomenernes Forhold%
% footnote 1) on p. 56
\footnote[1]{Smlgn. R. \emph{Körner}: Die logischen Grundlagen der Systematik
der Organismen (Wundt's Studien. II). p. 208.}. Tillige viser den
videnskabelige Begrebsdannelse, ved hvilken ensartede og analoge Fænomener
sammenfattes under én repræsenterende Forestilling, hen til det samme Motiv,
som virker ved al Lighedsassociation: Tendentsen til saa lidt som muligt at
ændre den nærværende Tilstand, til at fastholde det givne Synspunkt, uagtet der
gaas over til nye Forestillinger. Denne Tendents tilfredsstilles snart ved, at
vi tillempe vore tidligere Forestillinger efter de nye Erfaringer (i Stedet for
at danne helt nye Forestillinger), snart derved, at vi saa vidt muligt opfatte
de nye Erfaringer i Belysning af vore allerede givne Forestillinger.
Kontinuiteten i Erkendelsens Udvikling beror paa denne \emph{psykologiske}
«Sparsommelighedens Lov», som iøvrigt ikke altid fyldestgør, hvad den
\emph{metodologiske} Sparsommelighedslov fordrer, ti ofte kunne Fænomenerne
opfattes og forklares paa en simplere Maade end vi, ledede af Trangen til at
bevare Bevidsthedens Kontinuitet, foreløbigt ere i Stand til%
% footnote 2) on p. 56
\footnote[2]{Smlgn. herom E. \emph{Mach}: Über Umbildung und Anpassung im
naturwissenschaftlichen Denken. Wien 1884. --- \emph{Avenarius} har paa en
højst interessant Maade gennemført den ovenfor fremførte Tanke i sin
Afhandling: «Philosophie als Denken der Welt gemäss dem Princip des kleinsten
Kraftmasses». Leipzig 1876. --- Allerede \emph{Locke} har bemærket, at vi kun
danne almindelige Begreber og Sætninger og forme dem i Ord for at lette vor
Hukommelse og spare Tid baade, naar vi tænke, og naar vi ville meddele os til
andre. Essay III, 6, 30; IV, 12, 3.}. ---

Naturligvis kunne hverken Digteren eller Tænkeren ved Lighedsassociation
frembringe Forestillinger, de aldrig have havt før. Forberedende og supplerende
Processer spille her en vigtig Rolle, ikke mindre end ved den umiddelbare
Genkendelse (11). Al Association er en Reproduktion; men det ejendommelige ved
Lighedsassociationen er, at ved den reproduceres en Forestilling af en anden,
med hvilken den forud ikke har været i Forbindelse. Heri ligger det nye, det
«skabende». Muligheden for denne skabende Virk%
% --- p. 57 ---
\setcounter{footnote}{0}%
\opage{57}somhed er, som jeg har søgt at vise, given i Tendentsen til med saa
faa Variationer som muligt at bevare Bevidsthedens Kontinuitet.

Da der i det Foregaaende er lagt en stærk Vægt paa Følelsens Indflydelse, vil
jeg blot i al Korthed bemærke, at det ikke vilde være rigtigt \emph{altid} at
betragte Følelseselementet som Mellemled mellem de to Forestillinger, der søge
at associere sig formedelst Lighed. I nogle Tilfælde er det vistnok saa, f. Ex.
ved «Fornemmelsesanalogier», ved hvilke Forestillinger fra forskellige
Sanseomraader associeres, hvilket kun synes at kunne forklares ved, at de hver
for sig have Tendents til at vække samme Følelse. Naar vi sammenstille en høj
Tone med en aktiv Farve (Rødt, Orange eller Gult), kan det kun komme af den
beslægtede Følelsesvirkning%
% footnote 1) on p. 57
\footnote[1]{Psykologi p. 334. --- Den Mulighed kan dog ikke siges at være
ganske udelukket, at der kunde være en virkelig Lighed til Stede mellem
Fornemmelser fra forskellige Sanseomraader. \emph{Fick}: Die Lehre von der
Lichtempfindung (Hermann's physiol. Handbuch, III, 1) p. 166 hævder dette i
Modsætning til Helmholtz's skarpe Distinktion mellem de forskellige
«Modaliteter».}. Ved de Fænomener, jeg ovenfor omtalte, gives der derimod
(afset fra Følelseselementet) noget Ensartet og Analogt, og Følelseselementet
virker kun med, for saa vidt det driver til at søge beslægtede Forestillinger.
Den ene Forestilling (f. Ex. om Humlen, der kryber hen ad Jorden) vilde i og
for sig ikke vække just en saadan Stemning, som den, der i Exemplet er til
Stede; det er netop kun ved sit Slægtskabsforhold til den anden Forestilling,
at den kan tjene til Udtryk for Stemningen. --- Man vilde vel kunne paavise
Tilfælde, i hvilke begge Dele --- Forestillingsindholdets og
Følelsesvirkningens Slægtskab --- virkede sammen.

Man har ment, at der var forbundet en ganske særegen Vanskelighed ved at give
en fysiologisk Forklaring af Lighedsassociationen. Medens man ved
Berøringsassociationen henviser til Analogien med en Række indøvede Bevægelser,
der ved Gentagelse uvilkaarligt følge i samme Orden som før, mener man, at der
ved Lighedsassociationen ingen tilsvarende Proces kan paavises. Sagen synes dog
--- efter det ovenfor Udviklede --- ikke at frembyde nogen principiel
Vanskelighed. Tænke vi os, at der til enhver Forestilling paa en eller anden
Maade svarer visse Svingninger i Hjernens (eller en vis Hjernedels) Smaadele,
saa vil Hjernen, naar der i Bevidstheden er den ovenfor beskrevne Tendents til
Lighedsassociation, have en Tendents til at foretage en Række Svingninger, som
ligne dem, der svare til den Forestilling, der er Lighedsassociationens
Udgangspunkt. Samme Svingningsform eller Svingningsretning (hvorledes man nu
vil tænke sig Sagen) fastholdes med visse Ændringer. Man kunde maaske tænke sig
Sagen i Analogi med den Række Bølger af væsentlig samme Form og Retning, der
successivt danner sig, naar en Sten er kastet i Vandet. --- Denne
Forklaringsmaade er ingenlunde mere hypotetisk end den tilsvarende
hjernefysiologiske Forklaring af Berøringsassociationen.

% --- p. 58 ---
\setcounter{footnote}{0}%
\opage{58}25. Naar al Association er et Udtryk for Bevidsthedens Energi, maa vi
ogsaa vente i Associationen at genfinde Bevidsthedslivets almindelige
Grundform. Vi maa i Associationsprocessen kunne paavise den sammenfattende,
forenende Virksomhed, som er Bevidsthedens Kendemærke, og hvis Spor kunne
findes i ethvert Bevidsthedsfænomen. (Smlgn. Psykologi, II, 5; V B, 5.) Og
saaledes forholder det sig ogsaa. Association er jo, som Ordet allerede viser,
en forbindende Virksomhed. Hvad der sker i enhver Association, er egentligt
dette: ud fra det enkelte givne Bevidsthedselement konstruerer Bevidstheden
efter Lighedens og Berøringens samvirkende Principer en større Sammenhæng, i
hvilken det givne Bevidsthedselement udgør et enkelt Led. \emph{Associationen}
mellem Del og \emph{Helhed} bliver derfor den fundamentale Lov for
Forestillingsassociationen (se Psykologi, V B, 8 c). Jo mindre forskellige
Delene ere indbyrdes, des mere nærme vi os til en blot Lighedsassociation; jo
mere forskellige, des mere nærme vi os til en blot Berøringsassociation. Men i
Virkeligheden gøre begge Tendentser sig stedse gældende i Forening.
Virksomhedstrangen faar Luft baade ved Supplering af det brudstykkeagtigt Givne
og ved Bevægelse i Kredse udenom det. Baade gennem Lighedsforholdet og gennem
Berøringsforholdet faar det enkelte Element sin Plads i den hele sjælelige
Verdensorden.

Ikke blot fra den formale, ogsaa fra den reale Side set finde vi i Helhedsloven
den egentlige Grundlov for Forestillingsassociationen. Der finder ogsaa hvad
Bevidsthedens \emph{Indhold} angaar en Koncentration Sted, idet alle
Forestillinger saa vidt muligt bringes i Forhold til og ordnes efter deres
Slægtskab med den eller de Forestillinger, der ere inderligst knyttede til den
herskende Følelse og Villiesretning i Individet.

Det er Skotten \emph{William Hamilton}, som (efter tidligere Antydninger hos
Kant og Fries) har banet Vejen for denne Reduktion af de empiriske
Associationslove til én Grundlov.

Naar vi tale om Lighedsassociation og Berøringsassociation hver for sig, er det
altsaa Abstraktioner, vi operere med. De kunne ikke indbyrdes reduceres til
hinanden; men under nogle Forhold angiver den ene, under andre Forhold den
anden det vigtigste Synspunkt. Hvis man vil udstrække Begrebet
Lighedsassociation til ogsaa at omfatte den umiddelbare Genkendelse, saa viser
det sig, at Lighedsforholdet spiller en Rolle ved enhver Forestillingsovergang,
idet (som jeg i 2. Kap. har søgt at vise) al Berøringsassociation forudsætter
en Genkendelse. Vi have da her et psykologisk Udgangspunkt til Forstaaelse af
den Rolle, Identitetsprincipet spiller i den logiske Tænken. Den uvilkaarlige
og den vilkaarlige Tankevirksomhed bringes derved i naturlig Forbindelse med
hinanden.

Det var den ældre (engelske og herbartianske) Associationspsykologis
Ensidighed, at den betragtede de enkelte Fornemmelser og Forestillinger som
aldeles selvstændige Ting, der kun paa ganske udvortes Vis kom i Forbindelse
med hverandre. For den mest konsekvente Repræsentant for hele denne Retning,
David Hume, blev derfor selve Asso%
% --- p. 59 ---
\setcounter{footnote}{0}%
\opage{59}ciationsfænomenet en uløselig Gaade. Forestillingernes indbyrdes
Forskellighed og Uafhængighed var slaaet saa stærkt fast, at det nu blev
umuligt at forstaa deres Forbindelse. Man opfattede Bevidsthedens Enhed blot
som et \emph{Resultat} af Associationen, og saa ikke, at den netop er en
\emph{Forudsætning} for al Association.

\begin{center}\large IV. Forestillingsassociation og Sammenligning.\end{center}

26. Teorien om Forestillingernes Association hænger i sin historiske Udvikling
nøje sammen med den Maade, paa hvilken man til forskellige Tider og fra
forskellige Synspunkter har opfattet Sjælelivet i det Hele. De Forskere, som
særligt bestræbte sig for at paavise bestemte, naturlige Love paa Sjælelivets
Omraade, maatte lægge stor Vægt paa Forestillingernes Association, da dette
Fænomen aller tydeligst godtgør, at der hersker en Lovmæssighed i den indre saa
vel som i den ydre Natur. De kom derved i en oppositionel Stilling især til de
spiritualistiske og mystiske Anskuelser, for hvem det stod som en Nedværdigelse
for det sjælelige Liv at være bestemt Lovmæssighed underlagt. De blev ved denne
skarpe Modsætning til en Retning, der hævdede Sjælelivets Frihed for al
Naturorden, førte til ikke blot at opfatte Sjælelivet i dets nøje Sammenhæng
med materielle Processer, men ogsaa til at betragte det som mere eller mindre
analogt med saadanne, hvilket navnligt viste sig i, at man sammenstillede
Sjælelivets Elementer med de materielle Atomer. Det er nu en Gang saaledes, at
Sprogets Udtryk om sjælelige Fænomener ere Metaforer, Ord, som oprindeligt
brugtes om materielle Fænomener (smlgn. 21); og man maa derfor vel agte paa, at
man ikke lader sig lede af urigtige Analogier. De Distinktioner, vi foretage
hvad de sjælelige Fænomener angaar, og betegne ved Ordet Elementer, maa ikke
gøres til selvstændige reale Væsener. Forestillinger f. Ex. kunne med Rette
kaldes Bevidsthedselementer; men ikke i ganske samme Betydning, i hvilken vi
kalde Atomerne for Materiens Elementer. Enhver Forestilling er en Virksomhed,
en særlig Retning af Forestillingsvirksomheden, og denne igen en særlig Side af
den hele Bevidsthedsvirksomhed i et vist Øjeblik. De sjælelige Elementer have
altsaa deres materielle Paralleler ikke i materielle Atomer, men i materielle
Processer. Ved en urigtig Analogi kom nu den saakaldte «Associationspsykologi»
til at opfatte Associationen som en udvortes, mekanisk Proces, der forløber
mellem de hver for sig selvstændige og i Grunden af hverandre uafhængige
Forestillinger. Sit Højdepunkt naaede denne Ensidighed hos \emph{Hume}, der (se
25) gennemførte den saa vidt, at selve Associationsfænomenet blev ham
ubegribeligt: ti hvorledes --- spurgte han ganske konsekvent --- kan der være
en Forbindelse mellem Elementer, som ere helt uafhængige af hverandre? Man
havde hypostaseret Forestillingerne og overset, at de kun ere Retninger eller
Ytringer af sjælelig Aktivitet. Denne Aktivitet
% --- p. 60 ---
\setcounter{footnote}{0}%
\opage{60}blev skudt helt til Side, og hvor man, som \emph{Stuart Mill og Bain}
forsøgte, vilde raade Bod paa denne Ensidighed ved at anerkende en sjælelig
Aktivitet, der ytrede sig i Opmærksomhed og Sammenligning, dér blev der dog
ikke paavist nogen Forbindelse mellem den passive Proces, som gaar for sig i
Fornemmelse og Association, og den aktive Proces i Opmærksomhed og Sammenligning%
% footnote 1) on p. 60
\footnote[1]{Allerede i mit Skrift «Den engelske Filosofi i vor Tid» (1874) p.
37 og 87 har jeg gjort opmærksom herpaa. --- Det er ikke rigtigt at fremstille
den engelske Associationspsykologi som sluttet Enhed. En bestemt Ændring og
Udvikling med Forsøg paa Korrektion kan paavises, som jeg i anf. Skr. har søgt
at godtgøre; men fuld Klarhed og Konsekvents opnaaedes ikke af den engelske Skole.}.

Ikke mindre Vanskeligheder frembød der sig for den modstaaende, overvejende
spiritualistiske Retning i Psykologien. Man gik vel --- ofte modstræbende ---
ind paa Anerkendelsen af Associationslovene, men søgte at opfatte dem som
Udtryk for en lavere Side af Sjælelivet. Ligesom man i det Hele skelnede skarpt
mellem Sjæl og Legeme som to forskellige Væsener, saaledes spurgte man særligt
hvad hvert enkelt sjæleligt Fænomen angaar, om det skyldtes Legemets Indvirken
paa Sjælen eller Sjælens egen Virken. Associationsprocessen var man tilbøjelig
til at henregne til den Del af Sjælelivet, der umiddelbart bestemmes ved de
materielle Processer i Hjernen. Associationen finder ifølge Charles
\emph{Bonnet} egentligt Sted mellem Bevægelserne i Hjernens «Fibre», og til
denne Association mellem Aarsagerne i Hjernen, svarer
Forestillingsassociationen som en Association mellem de tilsvarende Virkninger.
Omvendt sætter ved Opmærksomhedsakten Sjælen paa aktiv Maade en eller anden
Hjernefiber i Bevægelse. I vore Dage har \emph{Lotze} stillet «den
sammenfattende og forholdssættende («beziehende») Bevidsthed» som en «ny»
Virksomhedsytring af Sjælen, en aandelig Virksomhed af højere Orden i
Modsætning til Reproduktion og Association. Hin højere Virken har vel den
lavere til Forudsætning, men den udvikler sig ikke af den%
% footnote 2) on p. 60
\footnote[2]{Ch. \emph{Bonnet}: Essai analytique p. 106. 255 ff. ---
\emph{Lotze}: Mikrokosmus. 2. Udg. I. p. 251 f. Drei Bücher der Metaphysik. p.
531. --- Medens \emph{Wundt's} tidligere Udtalelser om dette Punkt kunde
forstaas saaledes, at han betragtede Association og Apperception
(Bevidsthedsaktiviteten i Opmærksomheden) som to forskellige Processer, har han
nu i 3. Udgave af «Physiologische Psychologie» (II. p. 378) erklæret, at
Sondringen mellem Association og Apperception blot beror paa «en Abstraktion,
til hvilken Virkeligheden mere eller mindre kan nærme sig.»}.

Der er da al Grund til at gaa noget nærmere ind paa Forholdet mellem
Forestillingsassociation og Sammenligning.

27. Ligesom ved Genkendelsen og Associationen maa vi ogsaa ved Sammenligningen
skelne mellem forskellige Trin og Former. At sammenligne er at finde Lighed
eller Forskel eller begge Dele. Denne Virksomhed er os bedst bekendt i dens
højere Former; men den gaar uafbrudt for sig paa alle Trin af Bevidsthedslivets
Udvikling.

% --- p. 61 ---
\setcounter{footnote}{0}%
\opage{61}Der ligger to Genstande for mig paa Bordet saaledes, at jeg med ét
Blik kan se dem begge. Hvilken psykologisk Proces foregaar der nu, naar jeg
sammenligner dem? De to Genstande (\textit{A} og \textit{B}) ere ganske vist
begge indenfor min \emph{Synskreds}; men naar jeg sammenligner dem, bevæger min
\emph{Opmærksomhed} sig frem og tilbage mellem dem. Skønt der nemlig i hvert
Øjeblik kan være et vist Antal Elementer fremme i min Bevidsthed, fremtræde de
dog ikke alle i hvert Øjeblik med lige Liv og Tydelighed, selv om de ydre
Indtryk ere lige stærke. Snart staar det ene, snart det andet i Bevidsthedens
Midtpunkt, idet Opmærksomheden vendes imod det. I sin simpleste Form fremtræder
dette, naar Øjet er overmæt af et Farveindtryk og derfor søger et andet, som
danner Kontrast til det, eller ogsaa opsøger et mere neutralt Indtryk. Ved
samtidigt givne Genstande behøver nu den ene Genstand ikke at falde ud af
\emph{Synskredsen}, fordi \emph{Opmærksomheden} retter sig mod den anden. Men
vi mærke en ejendommelig Forandring som Følge af, at Opmærksomheden gaar fra
\textit{A} til \textit{B}. Denne Overgang er i Regelen forbunden med en vis
Fornemmelse af Anstrengelse, som i det hele haves ved opmærksom Sammenligning,
og som sandsynligvis er at forklare ved, at der sker en vis Anspændelse af
Sanseorganets Muskler, som er nødvendig, for at Indtrykket fra Genstanden skal
falde paa den «gule Plet»%
% footnote 1) on p. 61
\footnote[1]{Smlgn. Psykologi. 2. Udg. p. 136---138. 363---365.}. Resultatet af
denne Opmærksomhedsovergang, Opfattelsen af Forskellen eller Ligheden mellem
\textit{A} og \textit{B,} springer frem i Bevidstheden, som noget, vi vel have
forberedt, men hvis Karakter vi dog ikke ere Herrer over. Vi gøre ved
Opmærksomhedens Skiften et Experiment, men vi raade ikke for dets Udfald.

Det egentligt aktive Element ved Sammenligningen ligger altsaa i
Opmærksomheden, som skiftevis rettes mod de Genstande, der sammenlignes. Atter
her staa vi overfor to Begreber, der kun ved kunstig Abstraktion kunne holdes
ude fra hinanden, hvorfor man ogsaa undertiden har erklæret Opmærksomhed og
Sammenligning for ét og det samme%
% footnote 2) on p. 61
\footnote[2]{Saaledes \emph{Bonnet}: Essai analytique. p. 203. \emph{Lotze}:
Drei Bücher der Metaphysik. p. 540.}. Dog har det sin Betydning at agte paa
Forskellen mellem de to Elementer i vor erkendende Virksomhed, især da vi
ellers ikke forstaa, hvorfor en successiv Henvendelse af Opmærksomheden paa de
forskellige Genstande er nødvendig for at opnaa en klar Ligheds- eller
Forskelsopfattelse. Denne successive Henvendelse kan kun have den Betydning, at
Fornemmelsen forstærkes, ɔ: bliver tydeligere og livligere, og dette kan igen
kun forklares ved at Reproduktion og umiddelbar Genkendelse virker med. (Smlgn.
5.) Selv hvor Opmærksomheden er uvilkaarlig, d. v. s. ikke forudsætter nogen
forudgaaende Forestilling om det, som iagttages, ville Genkendelsesakter ikke
udeblive; ti da Opmærksomheden ikke virker kontinuerligt, men rytmisk, vil der
ske Reproduktioner ved alle de senere Henvendelser til samme Genstand.

% --- p. 62 ---
\setcounter{footnote}{0}%
\opage{62}Hvis de Genstande, der sammenlignes, ikke samtidigt ere til Stede for
Bevidstheden som hørende til en og samme Synskreds, men afløse hinanden, finder
en lignende Proces Sted som den nys skildrede. Fornemmer jeg først \textit{A,}
derpaa \textit{B,} maa jeg, hvis jeg skal kunne sammenligne dem, enten beholde
Forestillingen om \textit{A} i Bevidstheden eller reproducere \textit{A,}
efterat jeg er kommen til \textit{B.} Der tilvejebringes da ogsaa i dette
Tilfælde en Samtidighed mellem de Elementer, der sammenlignes. Ogsaa her finder
tillige (i hvert Tilfælde ved vanskeligere Sammenligningsakter) en successiv
Henvendelse af Opmærksomheden Sted fra det for Fornemmelsen givne \textit{B}
til Forestillingen om \textit{A} og omvendt. --- Det samme gælder naturligvis,
naar det er to Forestillinger, som skulle sammenlignes. Forestillinger kunne jo
genkendes lige saa vel som Fornemmelser (6).

Hvis jeg ved Opmærksomhedens Overgang fra \textit{A} til \textit{B} slet ingen
Forskel mærker (afset fra selve Anspændelsen ved Overgangen), eller kun delvis
eller i ringe Grad mærker nogen Forskel, opfatter jeg deres Forhold som
Identitet, eller som sammensat Lighed (ved hvilken nogle Elementer ere
identiske, andre forskellige), eller som Kvalitetslighed eller som Forholdslighed.

Den Art Sammenligning, som her er beskreven, kan kaldes fri og
\emph{uvilkaarlig}: fri, fordi Elementerne enten samtidigt eller successivt
fremtræde som selvstændige Led i Bevidstheden, --- uvilkaarlig, fordi ingen
Hensigt, ingen forudgaaende Forestilling er forudsat. Man ser her klart,
hvorledes den aktive og den passive Side samvirke: en Genstand tiltrækker sig
Opmærksomheden og vækker altsaa en Kraftanspændelse og en Reproduktion, og som
Virkning af denne aktive Henvendelse opstaar Ligheds- eller Forskelopfattelsen.
Det ses ligeledes klart, at Bevidsthedens almindelige sammenfattende
Grundvirksomhed er en nødvendig Forudsætning for Sammenligningen%
% footnote 1) on p. 62
\footnote[1]{I mit Skrift «Den engelske Filosofi i vor Tid» p. 37 staar
urigtigt, at Sammenligning er Betingelse for Sammenfatten. Forholdet er det
omvendte.}. Kun naar flere Elementer sammenholdes, er Sammenligning mulig.

Fra denne frie, uvilkaarlige Sammenligning skelne vi to Former eller Trin, som
ligge hver paa sin Side af den, idet den ene er af mere simpel og usammensat,
den anden af mere indviklet og sammensat Natur.

28. Den erkendende Bevidstheds mest primitive Elementer ere Fornemmelserne.
Lige saa lidt som de andre Bevidsthedselementer ere de aldeles selvstændige og
uafhængige i Forhold til hverandre. Deres Opstaaen og deres Beskaffenhed bero,
som især de nyere Undersøgelser have godtgjort (smlgn. Psykologi VA), paa det
indbyrdes Forhold imellem dem, eller nøjagtigere mellem de Betingelser, under
hvilke de blive til. Hvad der i den enkelte Fornemmelse kundgør sig for os, er
egentligt Forskellen mellem to Tilstande (
% --- p. 63 ---
\setcounter{footnote}{0}%
\opage{63}i Sjæleliv og Hjerne), idet et vist Forskels- eller
Modsætningsforhold mellem vore Tilstande er Betingelsen for Fornemmelsens
Tilbliven og bestemte Beskaffenhed. Dette kunne vi udtrykke ved, at enhver
Fornemmelse bliver til ved en \emph{Skelnen} eller en Forskelsopfattelse.

Det kunde synes uberettiget her at tale om en Sammenligning, da de to Led, der
staa i Forhold til hinanden, ikke ere Genstand for Bevidsthed, saaledes som ved
den egentlige, frie Sammenligning. Man har ogsaa ofte villet sondre mellem
selve Fornemmelsen og dens «Vurdering», d. v. s. dens Sammenligning med andre
Fornemmelser. Men en saadan Sondring er ikke berettiget. Ti om Fornemmelsen
bliver til, og med \emph{hvilken} Beskaffenhed, det beror netop paa Forholdet
mellem de forskellige Tilstande (eller --- ved samtidige Indtryk --- de
forskellige Dele af samme Tilstand). Fornemmelsen kan altsaa siges at være
vurderet, før den opstaar. Dens Bestemthed ved Forholdet til de andre Elementer
er ikke noget, som sker efter dens Tilbliven, men betinger netop denne. Naar
den er bleven til, kan den paa den i 27 beskrevne Maade gøres til Genstand for
fri Sammenligning. Men hvad vi her ville fremhæve, er den vigtige Omstændighed,
at den \emph{enkelte Fornemmelse springer frem} i \emph{Bevidstheden paa samme
Maade}, \emph{som Ligheds}- \emph{eller Forskelopfattelsen springer frem},
\emph{naar} vi ved fri \emph{Sammenligning sammenholde to Bevidsthedselementer}.

Vi befinde os her paa Grænsen af det ubevidste Sjæleliv (om dette Udtryk kan
tilstedes), til hvilken Selviagttagelsen ikke umiddelbart kan trænge frem, og
for hvilken Sproget ikke har dannet aldeles passende Udtryk. Men det synes at
være fuldkomment berettiget at sammenstille Grænsefænomenerne med, hvad den
klarere, mere udviklede Bevidsthed lærer os, naar vi blot ere opmærksomme paa
de af de forskellige Forhold følgende Forskelligheder. Det er en berettiget
Analogi at kalde Fornemmelsen en Skelnen; ti Grundforholdet er det samme, kun
at Leddene ved fri Sammenligning ere bevidste, medens de ved denne
\emph{elementære Sammenligning} ere ubevidste. Der er da her udfoldet en
Aktivitet, hvis Resultat fremtræder for Bevidstheden, skøndt den selv øves
ubevidst. Vi faa Slutningssætningen uden Præmisserne (for at bruge dette
Billede), medens vi ved fri Sammenligning have baade Præmisser og
Slutningssætning. Hvis man vilde tale om en ubevidst Opmærksomhed som
Betingelse for al Fornemmelse, vilde dette Udtryk ikke mangle Berettigelse; ti
al Opmærksomhed er en Tillempelse, ved hvilken en Opfattelse bliver mulig; og
Tillempelsen vil kunne foregaa ikke blot uvilkaarligt, men ogsaa ubevidst.
Ogsaa paa dette Punkt gribe aktive og passive Processer altsaa ind i hverandre.

En saadan umiddelbar Skelnen finder Sted ved enhver Sammenligning, der
konstaterer Forskel. Alt hvad jeg ved fri Sammenligning kan gøre er at
sammenholde de opmærksomt betragtede Led; Forskellen træder umiddelbart frem,
naar Leddene stilles Side om Side.

% --- p. 64 ---
\setcounter{footnote}{0}%
\opage{64}Den mest primitive Lighedsopfattelse have vi i den tidligere
beskrevne umiddelbare Genkendelse. Der finder her væsentligt den samme Proces
Sted som ved fri Sammenligning, der konstaterer Lighed, kun at Leddene her ikke
(lige saa lidt som ved den elementære Skelnen) træde selvstændigt frem i
Bevidstheden. Denne Proces kunde man kalde \emph{bunden Sammenligning}, fordi
de to Led her ere umiddelbart sammensmeltede. Ved fri Sammenligning springer
Ligheden lige saa umiddelbart frem som Genkendelsen i de ovenfor (3---4)
beskrevne Tilfælde.

29. Af den frie, men uvilkaarlige Sammenligning (27) udvikler sig den
\emph{frie}, \emph{vilkaarlige Sammenligning}, som forudsætter Evnen til at
sætte sig et Formaal, stille sig en Opgave, fastslaa en Maalestok%
% footnote 1) on p. 64
\footnote[1]{Jeg bruger Ordet Vilkaarlighed i en noget anden Betydning end f.
Ex. \emph{Wundt} (Zur Lehre vom Willen. Studien I, p. 352 kfr. Physiol.
Psychologie. 3. Aufl. II. p. 243, 500), der kun taler om vilkaarlig Virken,
hvor et Valg mellem flere Muligheder finder Sted. Jeg kalder enhver Virken
vilkaarlig, ved hvilken en Forestilling om Formaalet er bestemmende, selv om
kun én eneste Mulighed fremtræder. Ogsaa Driften (ikke blot Forsæt og
Beslutning) er efter min Sprogbrug en vilkaarlig Aktivitet.}. Der gøres her
Brug af Opmærksomheden --- ikke, som før, blot til at \emph{opfatte} givne
Fornemmelser eller Forestillinger, men først og fremmest --- til at
\emph{fastholde} en Forestilling saa kontinuerligt som muligt i Bevidsthedens
Midtpunkt for at sammenholde de følgende Fornemmelser eller Forestillinger med
den. Der hører hertil en større Energi og Udvikling end til den uvilkaarlige
Opmærksomhed og Sammenligning, der Skridt for Skridt følger det givne Stof.

Overgangen fra uvilkaarlig til vilkaarlig Sammenligning er analog med
Overgangen fra Reflexbevægelse og Instinkt til Drift. Den væsentlige Forandring
er, at der forud for Virksomheden gør sig en Forestilling om dens Formaal
gjældende. Og her er i Virkeligheden mere end en Analogi. Den vilkaarlige
Sammenligning, den egentlige Tænkning er en Villiesytring: den sætter et Maal
og søger Midlerne til at virkeliggøre det. Det er fra først af kun praktiske
Motiver, der sætte Tænkningen i Bevægelse; der søges de ydre Midler til
Afhjælpelse af en Trang. Allerede et praktisk Formaal giver
Forestillingsforbindelsen en vis Orden og Konsekvents, idet kun de
Forestillinger blive Genstand for Opmærksomhed, der stemme med Forestillingen
om Formaalet, idet de anvise Midler til dets Opnaaelse. \emph{Thomas Hobbes}
forklarede træffende Mangelen paa Sammenhæng i Drømmeforestillingerne deraf, at
vi i Søvne intet fast Formaal have%
% footnote 2) on p. 64
\footnote[2]{Qvia ordo et cohærentia omnis a freqventi ad finem respectatione,
id est, a consilio oritur, necesse est, amissà per somnum cogitatione finis, ut
phantasmata alia aliis succedant, non amplius eo ordine, qvi ad finem tendit,
sed ut contingit. \emph{Hobbes}: De corpore. XXV, 9.}; derfor bevæge
Forestillingerne sig i alle mulige Retninger og have intet Centrum, om hvilket
de kunne ordne sig. Enhver Forestilling, til hvilken derimod en stærk og varig
Følelse knytter sig, fast%
% --- p. 65 ---
\setcounter{footnote}{0}%
\opage{65}holdes derved, og alle Associationer, som udgaa fra den og begunstige
den herskende Følelse, fastholdes. Den bliver gjort til
\emph{Associationscentrum}. Foruden Ligheds- og Berøringsassociation virker her
særligt ogsaa Associationen mellem Følelser og Forestillinger (Interessens
Lov). Naar det til en Forestilling hørende Følelseselement er meget svagt, idet
der er den bestemte Forvisning om, at det Fænomen, Forestillingen svarer til,
ikke vil indtræde, vil Fænomenet kunne forblive ubemærket, selv om det
indtræder. (Negativ Opmærksomhed.) Saaledes naar man forbyder et hypnotiseret
Individ at se en vis Person eller Genstand, som dog virkeligt er til Stede%
% footnote 1) on p. 65
\footnote[1]{Smlgn. Revue philosophique. Mai 1887. p. 455. --- Udtrykket
«negativ Opmærksomhed» (som bruges af \emph{Féré}: Sensation et Mouvement.
Paris 1887. p. 18) er rigtigere end «negativ Hallucination» (som bruges af
Bernheim og Pierre Janet). Ti Individet \emph{faar} en Fornemmelse, hvilket ses
af, at der undertiden opstaar en Kontrastfornemmelse (saa at Sonnambulen ser en
grønagtig Plet dér, hvor man havde forbudt det at se noget Rødt). Det er altsaa
Opmærksomhedens Bortvendelse, der er den egentlige Aarsag.}.

Jo klarere og bestemtere jeg formaar at forestille mig det Forhold, i hvilket
de søgte Forestillinger skulle staa til de Forestillinger, jeg har, des sikrere
og nøjagtigere skrider Tænkningen frem. Naar jeg i Matematiken sætter en Opgave
i Ligning, vil det netop sige, at jeg præciserer de Fordringer, der stilles til
de søgte Forestillinger, --- at jeg forestiller mig det Maal, jeg vil naa, den
Maalestok, jeg vil anvende, med fuldkommen Nøjagtighed. Her kunne vi derfor ved
en metodisk Fremgangsmaade fremtvinge den søgte Forestilling. Paa andre
Omraader ere vi henviste til at prøve os frem. I sin simpleste Form sker dette,
naar vi «bore» en Forestilling frem ved energisk Koncentration af
Opmærksomheden om den givne Forestilling, med hvilken den søgte Forestilling
vides at staa i en eller anden Forbindelse. Ofte er det ved Lighedsassociation,
at den søgte Forestilling springer frem, som ved Opdagelser af Homologier og
Analogier mellem Organer hos forskellige Arter. Det andet Sammenligningsled kan
her springe frem under Forskerens Fordybelse i det enkelte givne Fænomen. Som
med et Slag gaar da Slægtskabet mellem det foreliggende Fænomen og andre, i Tid
og Rum derfra langt fjernede, op for Bevidstheden. Der virker her den tidligere
(24) omtalte Kredsen om en herskende Forestilling. ---

Ifølge \emph{Aristoteles}%
% footnote 2) on p. 65
\footnote[2]{De anima III, 10.} er den praktiske Fornuft forskellig fra den
teoretiske derved, at den gaar ud fra et Formaal. Ved denne Lære har han lagt
Grunden til en dualistisk Sondring mellem Teori og Praxis, mellem Erkendelse og
Villie, som har været af uheldig Indflydelse paa den psykologiske Opfattelse af
begge disse Sider af Sjælelivet. Navnlig oversaa man, at der ved al Tænken gjør
sig et Formaal, en Stræben, en Interesse gældende. Kun ved Tanken om et Formaal
kommer der Orden og Sammenhæng ind i vore Forestillingsforbindelser, saa at en
\emph{gyldig} Erkendelse gjennem dem kan naas. Det Formaal, vi stille os,
behøver ikke at være et ydre; det kan være det, at forstaa Naturen
% --- p. 66 ---
\setcounter{footnote}{0}%
\opage{66}eller at bevare Konsekventsen og at udelukke Modsigelser mellem vore
Forestillinger. Der finder ved al Tænken et Valg Sted, idet vi af alle de
Fornemmelser og Forestillinger, der frembyde sig, fastholde dem, som
fyldestgøre de opstillede Fordringer. Denne vælgende Fremgangsmaade lægger sig
for Dagen i den \emph{Tvedelingens} Lov, der, som \emph{Wundt} har paavist%
% footnote 1) on p. 66
\footnote[1]{Logik. Leipzig 1880. I. p. 557.}, er ejendommelig for de
Forestillingsforbindelser, der komme i Stand ved egentlig Tænken (vilkaarlig
Sammenligning), til Forskel fra de uvilkaarlige Forestillingsforbindelser, der
brede sig i det Ubestemte. Tvedelingen fremkommer ved, at der paa ethvert Punkt
af Tankebevægelsen rejses det Spørgsmaal: stemmer den nu opdukkende
Forestilling med Forestillingen om de Betingelser, den søgte Forestilling skal
opfylde? Her stilles i hvert Øjeblik et Enten---Eller, og Sammenligningen
foregaar i hvert enkelt Øjeblik mellem to Forestillinger ad Gangen (af hvilke
den ene er Maalestokken), fordi Opmærksomheden i hvert Øjeblik kun kan
koncentreres om én Forestilling. ---

Der er ved den vilkaarlige Sammenligning fire Momenter at sondre imellem, af
hvilke nogle have en mere passiv, andre en mere aktiv Karakter: 1)
Udgangspunktet er det Stød, som ved den energiske \emph{Koncentration} om en
enkelt Forestilling gives Associationsprocessen i en vis bestemt Retning. 2)
Selve \emph{Associationsprocessen} forløber uvilkaarligt, skønt enkelte Led af
den fra første Færd ville fremhæves formedelst deres Forhold til den herskende
Interesse. Dette andet Moment forholder sig til det første omtrent som Luftens
Indtrængen i Lungen forholder sig til vor aktive, vilkaarlige Udvidelse af
Brystkassen. 3) \emph{Genkendelse} og Fastholden af de Forestillinger, der
synes at stemme med Forestillingen om Formaalet eller Maalestokken. 4)
\emph{Verifikation} af denne Genkendelse gennem vilkaarlige Dækningsforsøg, der
føre til umiddelbar Lighedsopfattelse eller umiddelbar Skelnen, og dernæst
maaske til en nærmere Bestemmelse af Lighedens eller Forskellens Grad.

Hvor det er en Række af successivt indtrædende Fænomener, der ere Genstand for
den vilkaarlige Sammenligning, træder (i det andet Moment) Iagttagelsernes
Række i Stedet for Associationernes. Forskellen er kun den, at ved
Iagttagelserne ligger Aarsagen til de nye Leds Opdukken udenfor os, ved
Associationen ligger den i os selv. Hvis vi kunne have nogen Indflydelse paa
den Orden, i hvilken Iagttagelserne komme, ville Associationerne ogsaa i dette
Tilfælde spille en Rolle med.

Dersom denne Skildring af den egentlige Tænken eller den vilkaarlige
Sammenligning er rigtig, er der ingen Grund til at antage en fra
Associationsevnen helt forskellig Tænkeevne (Vurderéevne, Fornuft eller hvad
man nu vil kalde den). Opmærksomheden, Villieselementet, som gør Tænkningen til
egentlig Tænken, findes ogsaa (som uvilkaarlig) ved al Fornemmelse og al
Forestillingsassociation. I den egentlige Tænkning antager den
% --- p. 67 ---
\setcounter{footnote}{0}%
\opage{67}kun en særegen Form. Saa betydningsfuld Overgangen fra det
uvilkaarlige til det vilkaarlige nu end kan være, baade hvad Erkendelsen og
hvad Villien selv angaar, saa er det dog ikke berettiget at antage, at en helt
ny Evne her skulde træde i Kraft.

En af de ældre Psykologer, der har lagt megen Vægt paa Villieselementet
(Opmærksomheden) i Tænkningen (samtidigt med, at han med stor Dygtighed
udviklede Associationslæren), \emph{Charles Bonnet}, har sagt: «L'attention est
la mère du génie»%
% footnote 1) on p. 67
\footnote[1]{Essai analytique. p. 313.}. Han finder altsaa den intellektuelle
Genialitet særligt i Evnen til at koncentrere Opmærksomheden om en Genstand.
Men saa vigtigt dette Moment ogsaa er, saa er det dog kun Indledningen, Stødet.
Og ikke mindre vigtigt er det, at en livlig Proces udløses ved dette Stød,
ligesom det ikke er nok, at vi kunne udvide Brystkassen, naar Luften af en
eller anden Grund hindres i at strømme ind. Det kommer da an paa, om den
Forestilling, der er gjort til Associationscentrum, kan udløse en lang og vel
sammenhængende Række Associationer. Og i denne uvilkaarlige, instinktagtige
Side ved Tænkningen, ytrer Geniet sig maaske nok saa meget som ved det første
Moment. Der er da ingen Grund til at betragte Associationsprocessen som en
lavere Form for sjæleligt Liv. Netop det højeste sjælelige Indhold kan udvikles
ad denne Vej.

Og endeligt maa det stedse erindres, at vi i vore psykologiske Terminologier og
Teorier ikke kunne undgaa at fremhæve Forskellen mellem det uvilkaarlige og det
vilkaarlige langt stærkere, end den fremtræder i Virkeligheden. \emph{Selve
Overgangen fra uvilkaarlig til vilkaarlig Forestillingsforbindelse sker} i
\emph{Kraft} af den \emph{uvilkaarlige Association mellem en Følelse} og en
\emph{Forestilling}, \emph{hvorved en Drift}, en af \emph{Forestilling} om et
\emph{Maal ledet Stræben}, \emph{bliver mulig}. Det sjælelige Livs Kontinuitet
viser sig stedse paa ny for os, naar vi fra de Forskelligheder, foreløbig
Iagttagelse frembyder, trænge ind til Livets indre Gang.

\begin{center}\large V. Om Begrebet sjælelig Aktivitet.\end{center}

30. Det gaar med Begrebet Aktivitet paa det sjælelige Omraade, som med Begrebet
Bevægelse paa den ydre Naturs Omraade. Ligesom man længe har været hildet i
den, tilsyneladende af Iagttagelse støttede Antagelse, at der gaves absolut
Hvile, saaledes er man ogsaa tilbøjelig til at tænke sig, at visse Former og
Trin af Bevidsthedsliv ere absolut passive. Og ligesom det hist har vist sig,
at alt er i Bevægelse, eller rettere, at Bevægelse og Hvile ere relative
Begreber, saaledes vil det ogsaa her vise sig ved nærmere Betragtning, at det
sjælelige Livs Natur bestaar i Virksomhed, og at det kun er forskellige
% --- p. 68 ---
\setcounter{footnote}{0}%
\opage{68}Arter og Grader af Virksomhed, vi stille i Modsætning til hverandre,
naar vi skelne mellem passive og aktive sjælelige Fænomener.

Fra to Sider er i det Foregaaende Bevidsthedsaktiviteten bleven skildret. Fra
den formale Side ytrer den sig i den sammenfattende og sammenholdende
Virksomhed, af hvilken alle Bevidsthedsfænomener bære Præg, og hvis Resultater
de ere. Fra den reale Side ytrer den sig i den ved al Association og
Sammenligning virkende Opmærksomhed, som i sin Natur er Et med Villien. Dersom
man vil betegne al sjælelig Aktivitet som en Villen, kan man altsaa skelne
mellem en formal Villen, som udgør Bevidsthedens egentlige Væsen, saa vidt som
vi kende dette, og som er Udtryk for dens Selvopholdelsestrang, --- og en real
Villen, der lægger sig for Dagen i den særlige Retning, Bevidsthedslivet i det
enkelte Tidspunkt slaar ind paa%
% footnote 1) on p. 68
\footnote[1]{Smlgn. Psykologi. 2. Udg. p. 113, 159, 367.}.

31. Hvorvidt og hvorledes blive vi os nu denne sjælelige Aktivitet bevidst?

Medens Sensualismen og Associationspsykologien søgte at skyde Begrebet om
sjælelig Aktivitet ganske til Side, idet de negtede, at psykologisk Erfaring
viste os andet end passivt modtagne Fornemmelser, der passivt omdannedes ved
Association, have andre psykologiske Opfattelser tillagt os en
\emph{umiddelbar} Bevidsthed om Aktivitet som udgørende vor egen inderste
Natur. Det er efter denne Opfattelse et primitivt psykologisk Faktum, som
Selviagttagelsen godtgør, at vi \emph{ville}, ɔ: anspænde os, anvende Kraft.
Denne Fornemmelse af Anspændelse og Kraftudfoldelse (effort) skal da være den
psykologiske Erfarings Grundfaktum, der ikke kan betvivles, men heller ikke
nærmere forklares, og i hvilket vor egentlige Natur kundgør sig for os. Denne
Tanke har \emph{Maine de Biran} udviklet i sine forskellige Skrifter. Overfor
Hume's Lære, at der ikke i nogen som helst Virksomhed af legemlig eller
aandelig Art var noget, som kunde føre til Forestillingen om Kraft eller
nødvendig Forbindelse, stillede han Paastanden om en umiddelbar Bevidsthed om
Jeget selv som Aarsag%
% footnote 2) on p. 68
\footnote[2]{Kfr. f. Ex.: Oeuvres philosophiques. Paris 1841. IV. p. 75. Le moi
se manifeste à lui-même par le sens intime, dans l'effort ou le mouvement
volontaire que l'âme aperçoit intérieurement comme un produit de son activité,
comme un effet dont sa volonté est cause.}.

Begge de stridende Parter beraabe sig paa den umiddelbare psykologiske
Erfaring. Denne maa da enten paa den ene eller paa den anden Side være
sammenblandet med uberettigede Slutninger. Vi ville nærmere undersøge den
saakaldte Villiesbevidsthed, der i sin simpleste Form fremtræder i den
Fornemmelse af Anspændelse, det er Maine de Birans Fortjeneste at have henpeget
saa energisk paa, selv om han maaske ikke har analyseret den tilstrækkelig
nøjagtigt. Der er saa meget des mere Grund til at foretage en indgaaende
Analyse af dette sjælelige Fænomen, som den nyere Psykologi i Virkeligheden be%
% --- p. 69 ---
\setcounter{footnote}{0}%
\opage{69}kræfter Maine de Birans Paastand om, at vi her staa overfor et
sjæleligt Grundfaktum. Medens den gamle Evneteori nemlig skelnede skarpt mellem
Erkendelse og Villie, lægger den nyere Psykologi stærk Vægt paa, at der ved al
Fornemmelse, Association og Sammenligning finder en stadig Medvirken af Villien
Sted. Tænkningen opfattes som en Villiessag, ikke mindre end den ydre Handling.
Derved rykkes ikke blot Spørgsmaalet om Villiesbevidstheden frem i første Række
blandt de psykologiske Problemer, men bliver ogsaa lettere at belyse. Det
drejer sig væsentligt om den rette psykologiske Opfattelse og Forstaaelse af
\emph{Opmærksomheden}.

Opsøge vi de aller simpleste Fænomener, ved hvilke Opmærksomheden tydeligt
lægger sig for Dagen, saa viser det sig, at den ytrer sig gennem en Bevægelse,
ved hvilken Sanseorganet uvilkaarligt bringes i et hensigtsmæssigt Forhold til
den paavirkende Genstand. (Smlgn. Psykologi V A, 7.) Der er altsaa her ingen
Tvivl om, at Opmærksomheden kundgør sig i Bevidstheden ved
\emph{Bevægelsesfornemmelser}. Og her falder Forskellen mellem indre og ydre
Handling bort, eller begge staa i den nøjeste Forbindelse, idet den ydre
Handling (Sanseorganets Tillempelse) er Betingelsen for den indre (Opfattelsen
af Genstanden). Anderledes kunde det synes at forholde sig, naar Opmærksomheden
rettes mod en Forestilling, og ikke mod en udvortes Genstand. Hvorledes kunne
her Bevægelse og Bevægelsesfornemmelse spille en Rolle med?

Dog viser ogsaa her Selviagttagelsen tydeligt nok, at al energisk Fastholden
eller Fremdragen af Forestillinger er forbunden med en Fornemmelse af
Anspændelse eller Anstrengelse. Og det er ligeledes let at godtgøre, at der ved
enhver saadan aandelig Anstrengelse finder en vis Muskelspænding Sted. Allerede
det i Tanker fordybede Menneskes Udseende viser dette. Og Sammenhængen mellem
ivrig Forestillingsvirksomhed og Muskelspænding vil ikke være vanskelig at
paavise. For det Første ville alle Forestillinger, hvor abstrakte og ideale de
end maatte være, tilsidst bero paa Fornemmelser og Reproduktioner af
Fornemmelser. Naar nu Opmærksomheden ved Fornemmelsen ytrer sig ved Bevægelse,
maa det ogsaa være Tilfældet ved den til Forestillingen knyttede Opmærksomhed,
da al Forestilling tilsidst kun er reproduceret Fornemmelse. Ved Fremdragen og
Fastholden af Forestillingen vil Bevægelsesforestillingen (den reproducerede
Bevægelsesfornemmelse) derfor spille en stor Rolle. Der vil i noget svagere
Grad ske det Samme, som naar det er ydre Genstande, der tiltrække sig
Opmærksomheden. For det Andet er enhver Opmærksomhed betinget ved en Følelse af
Lyst og Ulyst, og enhver Følelsestilstand er i større eller ringere Grad
forbunden med Bevægelse. Da anspændt Opmærksomhed er en intensiv
Følelsestilstand, vil der altsaa stedse være Muskelspænding til Stede, og
følgelig ogsaa Bevægelsesfornemmelser som en væsentlig Bestanddel af den hele
Bevidsthedstilstand.

Opmærksomheden, og med den Villiesbevidstheden i det Hele (ti ogsaa ved ydre
Handling gælder det først og fremmest om Fastholden af Forestillinger:
Forestillinger om
% --- p. 70 ---
\setcounter{footnote}{0}%
\opage{70}Maal og om Midler), vilde da, synes det, kunne føres tilbage til
Bevægelsesfornemmelser, dels oprindelige, dels reproducerede. De to Begreber:
sjælelig Aktivitet og Bevægelsestendents vilde være identiske. Man vilde da
uden den Mystik, som ofte kommer frem hos Maine de Biran, kunne hævde den
Sætning, at vi umiddelbart blive os vor Aktivitet bevidst, idet Bevidsthed om
vor Aktivitet vilde blive det samme som Bevidsthed om Bevægelse eller
Bevægelsestendents. Saaledes bemærker \emph{Wundt}, at al Opmærksomhed (eller i
hans Terminologi: Apperception) er en bevidst Virksomhed, og at vi umiddelbart
fornemme den som en indre Handling%
% footnote 1) on p. 70
\footnote[1]{Physiologische Psychologie. 3. Aufl. I. p. 535. II. p. 466 ff.}.
Og han udtaler udtrykkeligt, at naar han kalder Apperceptionen en Virksomhed,
er det nærmest paa Grund af de Innervationer, der finde Sted ved den%
% footnote 2) on p. 70
\footnote[2]{«In diesen Innervationsempfindungen liegt der nächste sinnliche
Anlass dafür, dass wir die Apperception oder Aufmerksamkeit eine
\emph{Thätigkeit} nennen». Zur Theorie des Willens (Studien I) p. 347.}.

Imidlertid vilde denne Identifikation af Aktivitetsbevidsthed og
Bevægelsesbevidsthed kunne gøres til Genstand for kritiske Betænkeligheder. Den
maa nødvendigvis antage, at der gives Bevægelsesfornemmelser, som svare til
selve Bevægelsesimpulsens Udgang fra de motoriske Centre i Hjernen, eller til
selve Innervationen, Nerveprocessens Indledelse. Men dette er endnu et omstridt
Punkt. Naar vi holde os til Selviagttagelsen, synes der ganske vist at være god
Grund til at antage to Klasser af Bevægelsesfornemmelser, af hvilke den ene
kunde siges at have en mere aktiv Karakter, da den antages at tyde paa selve
Bevægelsens Indledning i Hjernecentret, medens den anden er af mere passiv Art,
idet vi ved den faa Underretning om Musklernes Tilstand som Følge af de
indledte eller udførte Sammentrækninger. Man kunde kalde den første Klasse
Kraftfornemmelser, den anden Muskelfornemmelser. Hin skulde altsaa svare til en
fra Hjernen udgaaende (efferent) Proces, denne til en til Hjernen kommende
(afferent) Proces. Det er klart, at hvis Bevidstheden om sjælelig Aktivitet
skal være Et med Bevægelsesfornemmelse, saa maa det være med «efferente»
Fornemmelser. Ti dersom der kun gaves «afferente» Fornemmelser, altsaa kun
«Muskelfornemmelser» og ikke «Kraftfornemmelser», vilde vi jo i
Bevægelsesfornemmelserne kun have Bevidsthed om en Virksomheds Resultat, ikke
om selve Virksomheden. Og vilde man alligevel fastholde, at der gives en
umiddelbar Bevidsthed om sjælelig Aktivitet, saa kunde denne Bevidsthed ikke
være Bevægelsesbevidsthed, da denne (hvis alle Bevægelsesfornemmelser vare
«afferente») ikke vilde være mere aktiv end alle andre Fornemmelser.

I sin Afhandling «The Feeling of Effort»%
% footnote 3) on p. 70
\footnote[3]{Anniversary Memoirs of the Boston Society of Natural History.
1880.} betragter den amerikanske Psykolog \emph{William James} den fysiske
Anstrengelsesfornemmelse som «en sammensat afferent Fornemmelse, der stammer
fra de spændte Muskler, de strammede Sener, de sammentrykkede Led, det spændte
Bryst, den lukkede Stemmeridse, de rynkede Bryn, de sammenpressede Kæber, o. s.
v.» Og fra denne Fornemmelse af fysisk Anstrengelse skelner han skarpt en
% --- p. 71 ---
\setcounter{footnote}{0}%
\opage{71}anden Art Anstrengelsesfornemmelse, der ikke er motorisk, men opstaar
ved «Anstrengelsen for at erindre, for at tage en Beslutning eller for at
hengive sig til et ubehageligt Arbejde». Denne sidste Art Anstrengelse, mener
han, er karakteristisk for den egentlige Villen, som er et rent indre,
sjæleligt Fænomen, der bestaar aldeles uafhængigt af, om der indtræder
Bevægelse eller ikke. Denne egentlige Villen skal ytre sig i det Bifald, den
Tilslutning, vi give en Tanke, i det \emph{Fiat}! med hvilket vi slaa fast, at
en Handling skal ske. Hvad der kræver Anstrengelse og Villiesanspændelse er
Fastholdelsen af en Tanke, selv om den medfører Ulyst eller Smerte. Denne med
al Tro og al Villen, med de fleste Ytringer af aandelig Virken forbundne
Anstrengelsesfornemmelse staar efter James i Modsætning til alle «afferente»
Fornemmelser (altsaa ogsaa til alle Bevægelsesfornemmelser). Om den
fysiologiske Proces, til hvilken den er knyttet, vil han ikke nærmere udtale
sig; kun er han sikker paa, at det er Forestillingens, ikke Bevægelsens Centra,
de ideationale, ikke de motoriske Centra, som her spille en Rolle.

Jeg skal her ikke gaa nærmere ind paa Drøftelsen af det Spørgsmaal, om alle
Bevægelsesfornemmelser ere periferiske (afferente), eller om der ogsaa gives
centrale (efferente) Bevægelsesfornemmelser (Innervationsfornemmelser)%
% footnote 1) on p. 71
\footnote[1]{I 3. Udgave af sin «Physiologische Psychologie» (I. p. 405), har
\emph{Wundt} repliceret til James's Kritik af Læren om
Innervationsfornemmelser. Som et experimentum crucis, der skulde stadfæste
Antagelsen af centrale (efferente) Innervationsfornemmelser havde man
(\emph{Helmholtz}, \emph{Gräfe}, \emph{Wundt}) fremført den Erfaring, at
Patienter, der lide af Parese af en vis Øjenmuskel (m. rectus externus),
forlægge Genstandene længere til Siden end de virkeligt ere, hvilket saa
forklares af den stærkere (forgæves) centrale Anspændelse for at bevæge Øjet.
Herimod bemærker \emph{James} (i Tilslutning til \emph{Ferrier}), at man har
overset, at det andet Øje, som ved Forsøget var bedækket, uvilkaarligt bevæger
sig i Retning af Genstanden, og det med saa meget des større Energi, som det
syge Øje ikke kan bevæges. Fra hint sunde Øje kommer der nu afferente
Paavirkninger, og Forlæggelsen til Siden kan altsaa meget godt være motiveret
ved periferiske Fornemmelser (On the Feeling of Effort. p. 12). I sin Replik
gør \emph{Wundt} opmærksom paa, at Forlæggelsen til Siden ved Parese kun
indtræder, naar det normale Øje er lukket; naar det er aabent og kan medvirke
ved Retningslokalisationen, indtræder Illusionen ikke, eller \emph{hvis} der
opstaaer Dobbeltbilleder, er det kun for det lammede Øjes Billedes Vedkommende.
--- Adhuc sub judice lis est.}. Ti jeg tror, at man fra begge Sider har lagt
for stor Vægt paa dette Spørgsmaal.

For det Første vilde \emph{Bevægelsesforestillinger} meget godt kunne gøre sig
gældende ved en Opmærksomheds- og Villiesakt, selv om alle
\emph{Bevægelsesfornemmelser} vare periferiske. Ti fordi Bevægelsesfornemmelsen
skyldes en Paavirkning fra Musklen til Hjernen, kan den meget vel som enhver
anden Fornemmelse reproduceres, og svarer da væsentligt til en central Proces,
der dog, som alle stærkere centrale Processer, har en Tendents til at brede sig
og altsaa til at genfrembringe den oprindelige, periferisk fremkaldte Proces
fuldstændigt. Ligesom enhver anden Forestilling har ogsaa
Bevægelsesforestillingen en Tendents til saa meget som muligt at nærme sig den
oprindelige Fornemmelse,
% --- p. 72 ---
\setcounter{footnote}{0}%
\opage{72}hvis Reproduktion den er. Det Faktum, at vi kunne have motoriske
Hallucinationer, ligesom vi kunne have Syns- eller Hørelseshallucinationer,
peger i denne Retning.

For det Andet vil Bevægelsesfornemmelsen, selv om den altid svarer til
periferiske Processer, meget godt kunne indtræde, før Bevægelsen er fuldendt,
og vil altsaa meget vel kunne svare til Bevægelsens \emph{Indledelse}. Musklens
Sammentrækning tager nogen Tid, og hvad vi psykologisk kalde Kraftfornemmelsen,
kunde antages at svare til Sammentrækningens aller første Stadium, hvor den
endnu ikke er synlig i det Ydre. Hvad vi psykologisk kalde Muskelfornemmelsen,
kunde da svare til den \emph{fuldendte} Sammentrækning, og vilde for saa vidt
med Rette kunne siges at have en mere passiv Karakter.

Men hvad enten man nu opfatter Sagen paa den ene eller paa den anden Maade,
ville Bevægelsesfornemmelserne ikke være tilstrækkelige til at begrunde den
inderlige og absolute Bevidsthed om Aktivitet, som Maine de Biran har søgt at
paavise. Hvad enten vor Villiesbevidsthed har sit fysiologiske Korrelat i
Processer, der umiddelbart udløses i selve Hjernen eller i saadanne, der kun
udløses der formedelst Indtryk fra Musklerne, saa vil den dog i begge Tilfælde
svare til Processer, der selv ere Virkninger af foregaaende Processer. Der vil
i sig selv ikke være nogen Grund til, at særlig disse Processer --- Udløsning
af Impulser til Bevægelsesnerver --- skulde kaldes aktive i højere Grad end
alle de andre Udløsningsprocesser, der gaa for sig i Nervesystemet.
Bevægelsesfornemmelser have deres eget ejendommelige Præg, ved hvilket de
skille sig fra andre Arter af Fornemmelse; men der er ikke dermed givet, at
dette Præg særligt skulde være at betegne som Aktivitet, og endnu mindre er det
givet, at al sjælelig Aktivitet skulde kundgøre sig gennem denne Art
Fornemmelse. \emph{Gennem Parallelismen med fysiologiske Processer} vil man
derfor neppe kunne oplyse, hvorfor visse sjælelige Tilstande synes os at have
en særligt aktiv Karakter.

Betragte vi nu Sagen fra den \emph{rent psykologiske Side}, rejser det
Spørgsmaal sig, om vi i det Hele umiddelbart kunne blive os Aktivitet bevidst.
Alt hvad vi ved Selviagttagelse umiddelbart finde i vor Bevidsthed --- ogsaa i
de af Maine de Biran og James saa stærkt fremhævede Anspændelsesfænomener ---
er dog kun Fornemmelser, Forestillinger og Følelser af forskellige Kvaliteter
og i forskellig Grad af Sammensætning. Der er tillige stor Forskel paa, i
hvilken Grad de enkelte Elementer komme til at raade i Bevidstheden, hvilken
Magt de have i Forhold til andre Elementer. Men kunne vi \emph{umiddelbart}
mærke denne Magt? I saa Fald maatte Virksomhedsbevidstheden være ganske simpel
og usammensat, som vore simple Sansefornemmelser ere det; den vilde fremtræde
med en ejendommelig Kvalitet, der var lige saa lidt til at tage fejl af, som vi
kunne tage fejl af vore Farvefornemmelsers Kvalitet. \emph{Spørgsmaalet} er, om
\emph{Aktivitetskvaliteten virkeligt er saaledes knyttet til ganske simple
Bevidsthedstilstande}.

% --- p. 73 ---
\setcounter{footnote}{0}%
\opage{73}32. Den ejendommelige Tilstand, i hvilken vi ere ved enhver
Beslutning, enhver Afgørelse, til hvilken vi mene at være komne, eller selv
hvor Afgørelsen ikke er sket, men hvor vi søge at fortrænge nogle Elementer i
vor Bevidsthed og at fastholde og begunstige andre, kan ikke beskrives
anderledes end som en Koncentration af Forestilling og Følelse saa vidt muligt
paa ét Punkt, nemlig om Forestillingen om den Handling, vi ville udføre, eller
den Antagelse, vi anerkende. Der har dannet sig et Associationscentrum (smlgn.
29), som vi betragte som vort Jegs væsentligste Indhold i Øjeblikket, og gøre
os saa vidt muligt til Et med. Selv om andre Forestillinger og Følelser rejse
sig og lægge Beslag paa en Del af vor Bevidsthedsenergi, staa de dog som mere
fremmede overfor os; vi genkende ikke os selv i dem med den Klarhed, som i hint
Centrum; de staa for os som noget nyt og udefra paatvunget. Vi ere
tilnærmelsesvis i en Tilstand af Monoïdeïsme: én Forestilling (med tilhørende
Følelse) behersker os. Vi have paa en Maade hypnotiseret os selv.

Grunden til, at vi betegne saadanne Tilstande som aktive, ligger vistnok i to
Omstændigheder. \emph{Dels} maa Grunden søges i selve den koncentrerede,
tilspidsede Form, vort Bevidsthedsliv faar, og som udelukker, eller tenderer
til at udelukke enhver indre Modstand fra andre Bevidsthedselementers Side. Saa
snart derimod flere forskellige Elementer gøre sig gældende og lægge Beslag paa
vor Energi, uden at kunne sammensmelte mere eller mindre inderligt med én
herskende Forestilling have vi en Fornemmelse af Delthed, Splid eller i hvert
Tilfælde af Neutralitet eller Ligevægt. \emph{Dels} er Grunden at søge i, at
denne Koncentration efterfølges af indre eller ydre Forandringer, som med
usædvanlig Tydelighed kunne henføres til den koncentrerede Tilstand som Aarsag.
Enhver Tilstand, vi ere i, afføder indre og ydre Forandringer som sine
Virkninger. Men jo mere mangfoldigt og forskelligartet dens Indhold er, des
mindre tydeligt og slaaende fremtræder Kausalitetsforholdet mellem den og det,
der følger efter. Det er da intet Under, at vi særligt fornemme os som aktive,
som Aarsag, naar vor Bevidsthed er stærkt koncentreret.

Dersom Bevidstheden om Aktivitet var aldeles umiddelbar, maatte den første
Grund være tilstrækkelig. Den anden Grund forudsætter, at vi foruden vor
umiddelbare Bevidsthed ogsaa tage Hensyn til hvad der følger efter. Naar vi
finde Kriteriet for Aktivitet i det Følgende, er Bevidstheden ikke ganske
umiddelbar; en Sammenligning og Reflexion, en Art Slutning gør sig gældende.
Den anden Grund forholder sig til den første som Verifikationen til Hypotesen.
Det er urigtigt, naar man tillægger os en umiddelbar Bevidsthed om at være
Aarsag til noget i os eller udenfor os; en saadan Bevidsthed kan efter Sagens
Natur aldrig være umiddelbar. \emph{Kausalitet} kan kun \emph{opdages ved
Slutning}, \emph{ikke ved Intuition}. Hume beholder Ret: Iagttagelse (og hin
umiddelbare Bevidsthed er en Selviagttagelse) kan kun vise os et
Successionsforhold, aldrig et Kausalitetsforhold. Det er derfor Mystik, naar
Maine de Biran og andre Psykologer tillægge os en umiddelbar
Aktivitetsbevidsthed i den Forstand, at vi umiddelbart skulde mærke, at vi ere
Aarsag.

% --- p. 74 ---
\setcounter{footnote}{0}%
\opage{74}Lad os da holde os til den første Grund og undersøge, om den er
tilstrækkelig.

Det viser sig let, at det er eller kan være overordentlig vanskeligt at afgøre,
om den omtalte Koncentration har fundet Sted eller ikke. En vis Koncentration
er altid til Stede i Bevidstheden og er en nødvendig Betingelse for
Bevidsthedens Bestaaen. Det mere eller mindre mangfoldige Bevidsthedsindhold
maa stedse sammenfattes og sammenholdes; Bevidsthedens Karakter som Syntese kan
aldrig ganske fornegte sig, saa længe Bevidsthedslivet endnu bestaar. Der er
utallige Grader af den Energi, med hvilken Syntesen virker; men en vis
Koncentration følger nødvendigvis af Bevidsthedens almindelige Natur. Og afset
fra denne formale Sammenhæng vil der ogsaa stedse være en vis real Sammenhæng
til Stede: Følelse og Opmærksomhed ville i hvert Øjeblik være rettede mod
bestemte Forestillinger mere end mod andre. Paa \emph{hvilket Punkt} af den
\emph{Gradeskala af Koncentration}, som vi her \emph{kunne opstille},
\emph{skulle} vi nu \emph{tænke} os, at den \emph{Tilstand ligger}, \emph{som
særligt skulde have Krav} paa at \emph{kaldes Aktivitetsbevidsthed}? Skulde der
gives en egen «Klangfarve», ved hvilken netop denne bestemte Grad af
Koncentration kundgjorde sig for Bevidstheden selv? --- Dersom der gaves et
saadant umiddelbart og usvigeligt Kriterium, vilde nogle af de største
Vanskeligheder ved den praktiske Selverkendelse falde bort. Paa den ene Side
kan der være Tanker og Stemninger, som vi lege med og betragte som ufarlige, og
som dog senere kunne vise sig at have slaaet Rod og at være blevne til en
Villen i os. Paa den anden Side kan det være os aldeles umuligt at komme til
Klarhed over vor Villies Beskaffenhed, hvor den ikke faar Lejlighed til at at %
% PRINTED AS IS: „at at“ (p. 74)
 vise sig i Handling. Omvendt kan man plage sig selv ved at forvexle et
Øjebliks uædle Stemning eller Ønske med en afsluttet Villen og derved dømme sig
selv uretfærdigt. Til alle Tider have dybere gaaende moralske og religiøse
Erfaringer godtgjort, at her ligger et af de største Problemer paa det
sjælelige Livs Omraade. I den psykologiske Teori er det let nok in abstracto at
skelne mellem Passivitet og Aktivitet; den virkelige Erfaring finder det ikke
saa let, og den advarer netop bestemt mod at stole paa den umiddelbare
Bevidsthed: Noli credere affectui tuo, qvi nunc est; cito mutabitur in aliud
(De imitatione Christi. III, 33). --- Der foreligger ogsaa tydelig Erfaring nok
for, hvor vanskeligt det kan være at angive det bestemte Tidspunkt, da den
indre Afgørelse blev truffet%
% footnote 1) on p. 74
\footnote[1]{Smlgn. det i min Psykologi p. 396 analyserede Tilfælde.}.

En Omstændighed, som særligt volder Vanskelighed her for Selviagttagelsen, skal
jeg endnu paapege. Hvad der med største Heftighed og Voldsomhed gør sig
gældende i Bevidstheden og bevirker den største Koncentration paa et enkelt
Punkt, behøver derfor ikke at være det i Virkeligheden Stærkeste og Sejrende.
Der gives nemlig to Arter af sjælelig Styrke; vi kunne betegne dem som
Voldsomhedens og Inderlighedens Styrke. Afset fra enkelte pludselige
Svingninger vil det sjælelige Livs hele Retning overvejende be%
% --- p. 75 ---
\setcounter{footnote}{0}%
\opage{75}stemmes ved den sidstnævnte Art af Styrke. Men denne virker snarere
fordelt paa mange Bevidsthedselementer og mange Bevidsthedsøjeblikke, end
koncentreret i et enkelt Element og et enkelt Øjeblik.. Selviagttagelsen, der
holder sig til de mest fremspringende Punkter, har vanskeligere ved at blive
klar over Styrkeforhold af den mere inderlige eller fordelte Art. Og dog kan
vor \emph{Aktivitet}, \emph{vor Villen være stærkest} i \emph{saadanne Tider}
og paa \emph{saadanne Punkter}, \emph{hvor vi mærke mindst} til den i vor
\emph{Selviagttagelse}. Og naar det kommer dertil, at vi sige: «Jeg vil!»
(eller med W. James: «Fiat!») saa ligger det egentlige Vendepunkt oftest slet
ikke her, men langt tidligere og paa et helt andet Punkt. Hvad der sker ved vor
klart bevidste Beslutning, er da kun en Konstatering, en saa at sige officiel
Afslutning af hvad der i Realiteten er afgjort%
% footnote 1) on p. 75
\footnote[1]{Smlgn. herom Psykologi p. 108, 326 ff., 380, 397 f. --- Se ogsaa
\emph{Ehrenfels}: Ueber Fühlen und Wollen. Wien 1887. p. 76 og \emph{Ribot}:
Les maladies de la volonté. Paris 1883. p. 175.}.

Dertil kommer endnu, at den højeste Grad af Anspændelse og Koncentration, om
hvilken Selviagttagelsen bærer Vidnesbyrd, for Subjektet selv kommer til at
staa som en Hengivelsens, Lidelsens eller Passivitetens Tilstand. Det er ikke
tilfældigt, at to saa modsatte Tilstande som den hypnotiske, i hvilken
Individet har mistet Herredømmet over sig selv, og den mystiske, i hvilken
Individet med den største Selvanstrengelse og Selvvirksomhed fordyber sig i én
eneste Tanke, psykologisk set ere saa nær beslægtede. De have begge Karakteren
af Monoïdeïsme, og ved al Opgaaen i én Tanke eller i det Hele i ét
Bevidsthedselement gaa Aktivitet og Passivitet over i hinanden. Anspændelsen
danner Indledningen til Betagetheden, og hvad Mystikerne kalde Ekstasen
betegner det Punkt, hvor den mere aktive Tilstand giver Plads for den mere
passive. Den Anvisning, Mystikerne give til at naa dette Punkt, gaar simpelt
hen ud paa en stigende Ophæven af al Mangfoldighed og Forandring i Forestilling
og Erindring. Villien («Evnen til at elske») fastholder Genstanden og tillader
ingen Forestillingsvirksomhed, der fører bort fra denne; enhver opstigende
Erindring om andre Ting hæmmes. At denne ekstatiske Tilstand alligevel er mere
aktiv end den hypnotiske, ses af, at medens det hypnotiserede Individ slaar ind
paa enhver Forestilling, der svarer til de Stillinger, man lader det indtage,
eller de Gestus, man lader det udføre, fastholder den ekstatiske Mystiker den
Tanke, om hvilken han har koncentreret sig, og antager de Stillinger, viser de
Gestus, der svare til den%
% footnote 2) on p. 75
\footnote[2]{Se \emph{Maury}: Le sommeil et les rêves. Paris 1865. p. 247.}.
Hin bestemmes stedse udefra indad, denne indefra udad. Og dog er Mystikeren
fuldt overbevist om, at det højeste Trin af Koncentration ikke naas ved
vilkaarlig Villen, men ved en «Naadesimpuls», ved en Indskydelse, i Forhold til
hvilken han forholder sig modtagende. Han griber ikke mere, men bliver greben%
% footnote 3) on p. 75
\footnote[3]{Hos \emph{Ribot}: (Les maladies de la volonté. Paris 1883.
123---35 og La psychologie de l'attention. Paris 1889. p. 138---50) findes
interessante Bidrag til den mystiske Ekstases Psykologi, især byggede paa St.
Theresa's Skrifter.}.

% --- p. 76 ---
\setcounter{footnote}{0}%
\opage{76}Hvad der saaledes finder Sted ved mystisk Ekstase, der er
Højdepunktet af indre Koncentration, finder i lavere Grad Sted ved enhver
Beslutning. Enhver, som har fattet en afgørende og betydningsfuld Beslutning,
vil have havt en umiddelbar Fornemmelse af, at der var noget, som drev ham
frem, en Kraft, der betog ham, og at hans indre Tilstand kun ufuldkomment
udtryktes ved Ordet Aktivitet. Det Udtryk, som ved en alvorlig Villiesafgørelse
saa ofte kommer igen: «Jeg kan ikke andet!» (smlgn. f. Ex. Luther i Worms) er i
denne Henseende karakteristisk. Man beraaber sig ikke paa sin egen
Selvvirksomhed, men paa en objektiv Umulighed af at ville anderledes. Vi ikke
blot gribe det, som besluttes, men gribes selv af det. Hvor det Aktive ender,
og hvor det Passive begynder, kan ingen Selviagttagelse vise os.

Hvad følger nu af alt dette? Det kunde synes at være den naturlige Følge, at vi
maatte opgive den sædvanlige Tredeling i Psykologien, idet Villien ikke længere
synes at kunne anerkendes som en selvstændig Side ved Bevidsthedslivet. Men en
saadan Slutning kan kun drages paa Grund af en Misforstaaelse af de
psykologiske Distinktioner. Disse maa ikke forvexles med reale Sondringer
mellem forskellige, for umiddelbar Iagttagelse givne Fænomener. Saa vilde vi
heller ikke med Rette kunne skelne mellem Erkendelse og Følelse, ti denne
Distinktion grunder sig heller ikke paa umiddelbar Iagttagelse, men dels paa
det Faktum, at samme Fornemmelser og Forestillinger under forskellige Forhold
kunne være forbundne med forskellig Følelse, og omvendt, dels paa, at de to
Sider ved vore Bevidsthedstilstande vise sig at være forskellige Vilkaar
underlagte, hvad Udvikling, Reproduktion og Gentagelse angaar. Ligesom de
kemiske Elementer ikke kunne konstateres ved umiddelbar ydre Iagttagelse, men
opdages ved Analyse, saaledes maa de sjælelige Elementer, de indbyrdes
irreduktible Beskaffenheder ved vore Bevidsthedstilstande, ogsaa opdages ved
Analyse, naturligvis paa Grundlag af Selviagttagelse. Hvad den aktive
(villende) Side ved Bevidstheden angaar, følger det nu af sig selv, at der her
ikke kan være Tale om umiddelbar Iagttagelse af den. Ti Aktivitet og Kausalitet
opdage vi ved Slutning ud fra de for Iagttagelsen givne Successionsforhold; i
en enkelt Iagttagelse, en umiddelbar Tilstand kunne de ikke være givne. Dette
gælder paa det sjælelige Omraade lige saa vel som paa det fysiske Omraade. I de
for Selviagttagelsen givne Bevidsthedsfænomener have vi stedse Aktivitetens
Resultater, ikke Aktiviteten selv.

Jo inderligere Aktiviteten er knyttet til Sjælelivet, des mere indlysende er
det, at den ikke kan være Genstand for en enkelt og umiddelbar Bevidsthedsakt.
Netop fordi Sjælelivet er en uafbrudt Villen, kan det blive vanskeligt at finde
Villien i en speciel Tilstand. Lige saa lidt som vi kunne bygge vort Begreb om
Jeget eller Selvet paa en enkelt Fornemmelse eller Forestilling, lige saa lidt
er dette muligt for Villiens Vedkommende. Det er altsaa saa langt fra, at
Negtelsen af en speciel Aktivitetsbevidsthed fører til at negte den sjælelige
Aktivitets Realitet, at den netop grunder sig paa Overbevisningen om, at
% --- p. 77 ---
\setcounter{footnote}{0}%
\opage{77}Sjælelivet bestaar under kontinuerlig Aktivitet, af hvis Resultater
Selviagttagelsen kun formaar at gribe enkelte fremspringende Punkter, ud fra
hvilke vi saa ad Konstruktionens Vej danne os et Begreb om den Aktivitet, der
ligger til Grund. Vor Erkendelse gaar her som overalt fra foreliggende, mere
eller mindre sporadiske og forskelligartede Erfaringer til Opfattelsen af den
indre, kontinuerlige Sammenhæng.

33. Der synes dog stadigt at hefte en vis Dunkelhed ved Begrebet Aktivitet. Vi
have fremhævet, at Sjælelivet helt igennem er aktivt. Denne Aktivitet kan
fremtræde mere eller mindre koncentreret, saa at (i Overensstemmelse med
Fechners Regel) den lavere, mere elementære Grad af Aktivitet, en saadan f.
Ex., som lægger sig for Dagen i Sansefornemmelse og Reflexbevægelse, i Forhold
til en højere Grad af Aktivitet, som i egentlig Tænkning og bevidst Beslutning,
kan fremtræde som en Passivitet. Men hvad forstaa vi da i det Hele taget ved
Aktivitet? Vi have jo hidtil brugt dette Ord uden at definere det.

Dersom man vilde sige, at dette Begreb ikke kan defineres, kunde meget synes at
tale for en saadan Paastand. Ti det er ikke blot det mest fundamentale Begreb i
Psykologien; det er ogsaa Grundbegrebet i enhver gennemtænkt Verdensanskuelse,
idet al passiv eller hvilende Væren ved nærmere Undersøgelse viser sig som en
Virksomhed. Newton advarede i sin Tid mod den sanselige Opfattelse af de ydre
Fænomener, der skelner mellem Hvile og Bevægelse, da «det jo kunde være, at
intet Legeme i Virkeligheden er i Hvile». Denne Sætning vilde vi kunne udvide
til at gælde ikke blot Bevægelse, men al Virksomhed. Dog dette vilde her føre
os for vidt, hvis det skulde tilstrækkeligt udvikles. Min Mening her er kun, at
man med lige saa god Ret i Filosofien kunde erklære Begrebet Aktivitet for
udefinerbart, som man i Fysiken kan erklære Begrebet Bevægelse for udefinerbart%
% footnote 1) on p. 77
\footnote[1]{I «Drei Bücher der Metaphysik» (1879) p. 534 kalder \emph{Lotze}
«den Begriff der Thätigkeit» «einen Begriff, von dem wir glauben werden, dass
er etwas Eigenthümliches und in der Welt wirklich Befindliches bezeichne,
\emph{obwohl} wir es \emph{ganz unmöglich finden}, das, was wir mit ihm im
\emph{Gegensatze zum blossen Geschehen meinen}, \emph{auf irgend eine
mechanischer Construction sich annähernder Weise zu bezeichnen}». ---
\emph{James Ward} siger i Mind. 1887. p. 569: «I doubt if activity can be
defined in terms that do not already imply it».}.

Dog vilde man kunne give en Definition af Begrebet Aktivitet ved at fastholde
dets nøje Sammenhæng med Begrebet Kausalitet. \emph{Spinoza} har i Henhold
hertil defineret Aktivitet paa følgende Maade: «Da siger jeg, at vi virke
(\emph{agere}), naar der sker noget i os eller udenfor os, hvis fuldstændige
\emph{Aarsag} vi ere, d. v. s. naar der af vor Natur følger noget i os eller
udenfor os, som klart og tydeligt kan forstaas af den alene. Derimod siger jeg,
at vi \emph{lide} (pati), naar der sker noget i os, eller følger noget af vor
Natur, til hvilket vi kun for en Del er Aarsag»%
% footnote 2) on p. 77
\footnote[2]{Ethica III. def. 2. --- En lignende Definition findes hos
\emph{Charles Bonnet}: Essai analytique. p. 93.}. Denne Definition har den
Fortjeneste at sondre
% --- p. 78 ---
\setcounter{footnote}{0}%
\opage{78}mellem Fænomenerne, det Givne, «det, der sker i os eller udenfor os»,
og dets Forklaring. Den appellerer altsaa ikke til nogen speciel Art umiddelbar
Bevidsthed, men gaar ud fra, at enhver saadan behøver en Forklaring. Den
antyder tillige, at en Del af Aarsagen til, hvad der sker i os eller følger af
vor Natur, altid ligger i os selv, saa at vi aldrig ere absolut passive; men
den vil indskrænke Aktivitetsbegrebets Anvendelse til saadanne Tilfælde, i
hvilke hele Aarsagen ligger i os selv.

Spørgsmaalet er, om den nogensinde fuldstændigt kan ligge i os selv, om nogen
Bevidsthedstilstand og nogen Handling kan forstaas ud fra vort foregaaende
Bevidsthedsliv alene, og om der ikke stedse er andre medvirkende Forhold. Vi
ændre derfor den anførte Definition saaledes: Jo \emph{større} en Del af
\emph{Aarsagen} til hvad der sker i os eller udenfor os, der ligger i os selv,
des \emph{mere} aktive ere vi. Og dertil maa udtrykkeligt føjes, at selv hvor
vi forholde os passive, bestemmes dog Virkningens Art og Grad stedse ogsaa ved
vor egen Natur, ikke blot ved det Paavirkendes Beskaffenhed. --- I Henhold til,
hvad der i det Foregaaende er udviklet, er det et Spørgsmaal for sig, om vi
blive os vor Aktivitet bevidst; og den Bevidsthed, vi kunne have om den, er i
hvert Tilfælde ikke altid tilforladelig.

Men denne Definition fremkalder da igen det Spørgsmaal: \emph{hvad menes} der
med «os \emph{selv}» \emph{eller} «\emph{vor Natur}»? --- Kan dette Spørgsmaal
finde nogen Besvarelse, og hvilken?

En hyppig Anskuelse giver følgende Svar. Det ligger i Sagens Natur, at enhver
Virksomhed maa udøves af nogen eller noget. Vi maa skelne mellem det, som
virker, og selve \emph{Virksomheden}. Særligt maa vi paa det psykologiske
Omraade skelne mellem det, som tænker, føler og vil, og selve Tænkningen,
Følelsen og Villien. --- Efter denne Anskuelse vilde vi altsaa være des mere
aktive (i Henhold til den givne Definition), jo mere de sjælelige Fænomener
udspringe af selve \emph{Bæreren} af alle vore Tanker, Følelser og
Beslutninger. Om dette noget, der «er» bevidst eller «har» Bevidstheden, vide
vi rigtignok intet; det erklæres for et \textit{x. ---} Dette er den
\emph{monadologisk}-\emph{spiritualistiske} Teori. Jeg formaar ikke at indse
hverken dens Berettigelse eller dens Mulighed.

For det Første kan den nævnte Teoris Opstaaen simpelt hen forklares ved
uberettiget Indflydelse dels af sproglige Udtryk, dels af materielle Analogier.

Sproget skelner mellem Subjekt og Prædikat og betegner i Regelen hvert af dem
med et eget Ord. Der dannes da et Ord for det virkende Væsen, et andet for dets
Virksomhed. Selv hvor det er vanskeligt for Anskuelsen at gøre denne Sondring
mellem det Virkende og Virksomheden, anvender Sproget dog et formelt Subjekt;
saaledes i upersonlige Sætninger («det regner»). Der ytrer sig her en Trang til
at analysere de umiddelbart givne Fænomener og holde de forskellige Elementer i
dem ude fra hinanden. Men sondrede sproglige Udtryk og logiske Distinktioner
maa ikke uden videre gøres til Et med reale Forskelle. Et er de Fornødenheder,
vi føle, naar vi ville tænke over Fænomenerne og
% --- p. 79 ---
\setcounter{footnote}{0}%
\opage{79}udtrykke vore Tanker, et ganske Andet er selve den reale Tilværelse
og dens Forskelligheder. --- Hin Distinktion mellem det Virkende og
Virksomheden, --- det, der «har» Bevidstheden, og selve Bevidstheden, kan
maaske finde sin Forklaring ved Sprogbrugen.

Ved materielle Fænomener (som have faaet den aller vigtigste Indflydelse paa
Sprogets Udtryksmaade) skelne vi med Lethed mellem Virksomhedens (her:
Bevægelsens) Udgangspunkt og selve Virksomheden (Bevægelsen). Vi søge stedse at
føre Bevægelsens Udspring tilbage til et Legeme, som befinder sig paa det Sted
i Rummet, fra hvilket Bevægelsen antages at stamme. Vi se ofte dette Legeme
forblive tilsyneladende roligt paa sin Plads, hvad enten den Bevægelse, det kan
volde, virkeligt sker eller ikke. Hvad enten Stenen falder ned fra Taget eller
ikke, forbliver Jorden tilsyneladende paa sin Plads. Naar den ene Billardkugle
sætter den anden i Bevægelse ved at støde til den, kalde vi den første for
aktiv i dette Forhold, og vi ere da her ikke i Forlegenhed med Aktivitetens
Udspring; vi kunne jo se den og forestille os den, aldeles afset fra, om den i
Øjeblikket fremkalder nogen Forandring eller ikke. Paa analog Maade tænker man
sig da ogsaa Sjælen sondret fra den Virksomhed, der antages at udgaa fra den.
Man frygter ligesom for, at Sjælelivet skal forflygtiges, naar det ikke hænges
op paa en Sjælesubstants som paa en forsvarlig Krog, skønt man ikke er saa
konsekvent at undersøge, hvor denne Krog er slaaet fast. Sjælesubstantsen synes
lige saa vel at kunne behøve en «Bærer», som de sjælelige Virksomheder. Det er
den anskuende Fantasi, ikke den egentlige Tænkning, som her fører Ordet. Det er
rumlige Billeder og Symboler, som forvexles med psykologiske Begreber. Trangen
til faste, hvilende Billeder gør, at man ikke kan fastholde Begrebet Virksomhed
i dets Ejendommelighed og Konsekvents. Den spiritualistiske Teori kan altsaa
forklares ved Fantasiens uvilkaarlige --- \emph{materialistiske} Tendents. Det
er ikke blot (smlgn. 26) den ensidige Associationspsykologi, som har ladet sig
beherske for meget af Analogien med de materielle Atomer; det er ogsaa
Tilfældet for den monadologisk-spiritualistiske Teoris Vedkommende. Forskellen
er kun, at hin betragter de enkelte Forestillinger og Følelser som Atomer,
medens denne gør selve det bevidste Væsen til en Slags Atom.

For det Andet faar Aktivitetsbegrebet kun en højst mystisk Anvendelse efter den
nævnte Teori, idet denne selv indrømmer, at vi ikke kende Bevidsthedens Bærer.
Vi kunne da ikke udlede noget indre eller ydre Fænomen af den, saaledes at vi
derved vinde en klar Forstaaelse af det. En videnskabelig Forklaring faa vi
kun, hvor Fænomenerne udledes af noget, hvis Natur vi i Forvejen kende af
Erfaring. Men Erfaringer kan jo aldrig vise os hint \textit{x.} Dette
\textit{x} kunde da i og for sig meget godt være noget andet end «os selv», et
andet Væsen, eller andre Væsener, af hvilke Fænomenerne udspringe. Vi vide
efter den nævnte Teori jo egentligt kun saa meget, at visse Fænomener
\emph{ikke} kunne finde deres Forklaring ved ydre, materielle Aarsager; at de
udledes af \textit{x,} vil jo kun sige, at deres
% --- p. 80 ---
\setcounter{footnote}{0}%
\opage{80}Udspring er ubekendt. Og selv om vi kendte det, vilde det staa som en
uløselig Gaade, hvorledes Virksomheden kunde udspringe fra noget, som selv ikke
var virkende.

Det er fuldstændig rigtigt, at vi ikke kunne gennemskue den sjælelige
Aktivitets Udspring, Forløb, Former og Love, ej heller dens Forhold til andre
Arter af Aktivitet. Men ingen af disse Gaader løses i mindste Maade ved at
antage et særligt Væsen bag ved Aktiviteten. --- Hvis man ved at antage, at der
ligger noget bag Bevidstheden, blot vil udtrykke netop dette, at der er saare
meget ved den sjælelige Virksomhed, vi ikke forstaa, saa har jeg intet mod
denne Sprogbrug at indvende, og jeg anvender den selv. Men metafysiske
Anskuelser bør ikke bygges paa et saadant Grundlag.

For det Tredie vilde vi meget vel kunne sige, hvad vi forstaa ved «os selv»
eller «vor Natur», uden at gaa ind paa den monadologisk-spiritualistiske
Hypotese. Idet vi holde os til Erfaringen, se vi Sjælelivet udfolde sig som en
Virksomhed, der gennemløber forskellige Stadier og Former. Vi søge at beskrive
disse Stadier og Former og at forklare Overgangene mellem dem. Paa ethvert
givet Udviklingstrin have vi da Sjælelivet for os med visse Egenskaber og
Dispositioner, og det gælder da at forstaa, hvad der, under visse Betingelser,
kan udvikle sig af dem. «Vi \emph{selv}» \emph{eller} «\emph{vor Natur}» er paa
\emph{ethvert givet Trin de Egenskaber og Dispositioner}, \emph{som vore
Bevidsthedsfænomener vise} sig i \emph{Besiddelse} af. Dette er et aldeles
klart og forstaaeligt Begreb, og vi behøve ikke at gaa tilbage til mystiske
Dybder. Hvor vi nødes til at antage ubevidste Mellemled og Dispositioner,
tydeliggøre vi os dem ved Analogien med de bevidste Virksomheder, ligesom vi
paa det materielle Omraade tydeliggøre os Ligevægtstilstanden som en Tilstand,
i hvilken modsatte Bevægelsestendentser gensidigt hæmme hverandre. Begrebet
Egenskab eller Disposition er ganske vist paa det sjælelige Omraade
vanskeligere at tydeliggøre og anvende end paa det materielle Omraade%
% footnote 1) on p. 80
\footnote[1]{Se herom Psykologi p. 162 kfr. p. 76 og 96.}; og der er saare
mange Dispositioner, som undgaa vor Opmærksomhed. Derfor vil vor psykologiske
Viden altid forblive mangelfuld. Men saa meget des større Betydning har det da
at bestemme de psykologiske Grundbegreber i saa nøje Tilslutning til Erfaring
som muligt.

Jo mere nu den fortsatte sjælelige Udvikling betinges ved de givne Egenskaber
og Dispositioner, des mere Aktivitet tillægge vi Sjælelivet; jo mere den
betinges ved udefra kommende (uafhængigt af Sjælelivets Fortid opstaaende)
Indflydelser, des større Passivitet tillægge vi det. Dette Forhold mellem et
Indre og et Ydre strækker sig gennem alle Livets Trin og gør sig gældende lige
saa tidligt, som vi i det Hele finde Anledning til at antage et Sjæleliv. Naar
Sjælelivet i det Hele kaldes aktivt, saa er det først og fremmest paa Grund af
den Enhedsform, med hvilken alt dets Indhold præges, og som ikke kan forklares
ved ydre Paavirkning. Men aktivt er Sjælelivet ogsaa, fordi enhver Fornemmelse
% --- p. 81 ---
\setcounter{footnote}{0}%
\opage{81}optræder med en Selvstændighed og en Kvalitet, der er betinget ikke
blot ved det enkelte, øjeblikkelige Indtryk, men ogsaa ved Bevidsthedens
samtidige og forudgaaende Tilstand, --- fordi det i Associationsprocessen paa
Grundlag af en ofte forsvindende lille Antydning udspinder en lang Række af
Forestillinger, --- fordi det i den egentlige Tænken, der kan defineres som en
med Opmærksomhed fremskridende Association, forfølger en enkelt Sammenhæng i
alle dens Konsekventser, --- fordi det som villende (i snevrere Betydning)
koncentrerer sin hele Energi i Tanken om et Formaals Virkeliggørelse. Ved alle
saadanne sjælelige Ytringer ligger Aarsagen \emph{overvejende} i en
forudgaaende sjælelig Aktivitet og de Dispositioner, den har frembragt.

Ligesom Kloderne i vort Verdenssystem ifølge den Kant-Laplaceske Hypotese havde
en vis Begyndelseshurtighed og Begyndelsesretning, da de løsrev sig fra den
store Taageklode, saaledes begynder vort Bevidsthedsliv --- paa det Punkt, hvor
det vaagner fra Ubevidsthedens Nat --- med at være aktivt efter de samme Love,
som det senere følger. Det begynder med Dispositioner, som ikke blot betinge,
hvad det senere kan blive, men ogsaa giver det en vis Selvstændighed overfor
ydre Indflydelser. Udviklingen begynder ikke fra et absolut Nulpunkt, men
skrider frem fra mere elementære til mere sammensatte Former af Aktivitet. De
individuelle Forskelligheder bero paa Aktivitetens Grad og paa den Retning, i
hvilken den fra først af bevæger sig.

\begin{center}\large VI. Sjælelig og fysisk Aktivitet.\end{center}

34. Det Spørgsmaal, hvorledes sjælelig Aktivitet forholder sig til fysisk
Aktivitet, har jeg udførligt behandlet i min Psykologi (2. Udg. p. 62---80,
92---98), hvor jeg har søgt at begrunde \emph{Identitetshypotesen}, efter
hvilken fysisk og sjælelig Aktivitet ikke i deres inderste Grund, men kun i
deres fænomenale Fremtræden ere forskellige. Jeg skal her kun vende tilbage til
dette Spørgsmaal for at besvare de Indvendinger, Prof. \emph{Kroman} i 2.
Udgave af sin «Tænke- og Sjælelære» har fremsat mod denne Hypotese, hvilken han
med et lidet heldigt Navn kalder Duplicismen. (Lidet heldigt er dette Navn af
den Grund, at det afleder Opmærksomheden fra, hvad der er det Væsentlige ved
Hypotesen, nemlig Hævdelsen af, at det er ét og \emph{samme} Væsen, en og
\emph{samme} Aktivitet, der fremtræder for os under de to Former.) De fremførte
Indvendinger ere væsentligt tre: a) Der skal ikke være nogen naturvidenskabelig
Grund til at opfatte Forholdet mellem det Sjælelige og det Materielle paa den
Maade, Identitetshypotesen udsiger. b) Det skal, naar Identitetshypotesen
lægges til Grund, være uforklarligt, hvorledes vi kunne vide noget om den
materielle Omverden. c) Det skal være uforklarligt, hvorledes der kan være
Identitet mellem to Ting, der ere saa forskellige som den ydre Mangfoldighed
(hvori Legemet bestaar) og
% --- p. 82 ---
\setcounter{footnote}{0}%
\opage{82}den indre Enhed (som er Sjælelivets Karaktermærke). Enhver af disse
Indvendinger prøve vi for sig. Denne Prøvelse vil give Anledning til at belyse
Problemet fra nye Synspunkter.

35. a) Naar vi i det Hele ville danne os en Hypotese om et Forhold, der i den
Grad ligger paa Grænsen af vor Erkendelse, som Forholdet mellem Bevidsthed og
Materie, saa kan en saadan Hypotese kun have nogen Betydning, naar den saa
strengt som muligt holder sig til de Principer og Love, der i Forvejen staa
fast i Videnskaben. En Hypotese, som begynder med at slaa en Streg over alle
Synspunkter og Metoder, ved hvilke vi ellers orientere os i vor Forsken, har
aldeles ingen Værdi, ligegyldigt om den ligefrem kan gendrives eller ikke. Den
videnskabelige Betydning, en Hypotese kan have paa et Punkt, hvor --- som her
--- egentlig Verifikation er umulig, kan kun være den at vise, hvilke logiske
Konsekventser vi med Hensyn til dette Punkt kunne drage af det, som i Forvejen
staar fast for os. En Hypotese om Forholdet mellem Sjæl og Legeme, der vilde
begynde med at tilsidesætte Inertiens Lov, den første Sætning i den
Naturvidenskab, vi hidtil besidde, en saadan Hypotese vilde ikke have nogen
Værdi. \emph{Hvis} det er \emph{berettiget ved Problemet om Forholdet mellem
Sjæl og Legeme ganske} at se \emph{bort fra Inertiens} Lov, saa \emph{existerer
der aldeles ikke noget Problem} af den \emph{angivne} Art; man har da i de
mange Diskussioner om denne Sag plaget sig med indbildte Vanskeligheder, og det
bliver inkonsekvent, hvis man vil vedblive at drøfte de forskellige Hypoteser.
Sagen frembyder i saa Fald ingen større Vanskeligheder, end der er forbundet
med at forstaa, hvorledes den ene Billardkugles Bevægelse kan volde den andens
Bevægelse.

Men --- siges der --- Naturvidenskaben opstiller kun Inertiens Lov som gældende
for Forholdet mellem Legemer indbyrdes. Denne Lov finder da naturligvis ingen
Anvendelse her, hvor det gælder Forholdet mellem Sjæl og Legeme. --- Denne
Opfattelse af Inertiens Lov er ikke forsvarlig. Denne Lov er opstillet som
Princip for Forklaringen af Legemers Bevægelse i det Hele og indeholder den
Fordring, at naar et Legeme forandrer sin Tilstand af Hvile eller retlinet
Bevægelse, skal der søges en \emph{ydre} Aarsag. Den angiver altsaa,
\emph{hvor} vi skulle søge Aarsagen, naar en materiel Forandring (en
Bevægelsesændring) sker. Dette var det væsentligste Fremskridt, som blev gjort
ved Lovens Opstilling. \emph{Galilæi}, Lovens første Opdager, vilde netop
afvise \emph{indre} Aarsager til et Legemes Bevægelsesændringer%
% footnote 1) on p. 82
\footnote[1]{Smlgn. \emph{Wohlwill}: Die Entdeckung des Beharrungsgesetzes
(Zeitschrift für Völkerpsychologie. XV). p. 114.}. \emph{Newton}, til hvis
Formulering af Inertiloven jeg især holder mig, siger i sin «første
Bevægelseslov»: «Et Legeme forbliver i den Tilstand, i hvilken det er, Hvile
eller jevn Fremskriden, saa længe ingen \emph{ydre Kræfter} tvinge det til at
forandre sin
% --- p. 83 ---
\setcounter{footnote}{0}%
\opage{83}Tilstand», og Begrebet ydre Kraft (vis impressa) forklarer han
nærmere paa følgende Maade (i den fjerde Definition): «Ydre Kræfter kunne være
af forskellig Oprindelse; de kunne hidrøre fra Stød, Tryk eller Tiltrækning.»
En nyere Fysiker, \emph{Maxwell} siger (i sit Skrift «Matter and Motion», § 41)
efterat have fremstillet Inertiloven: «Det experimentale Bevis for denne Lovs
Sandhed ligger i, at vi, hver Gang vi møde en Forandring i et Legemes
Bevægelsestilstand, kunne føre denne Forandring tilbage til en \emph{Virken
mellem dette Legeme} og et \emph{andet}, d. v. s. til en \emph{ydre Kraft}».

Inertiens Lov angiver den Metode, som Naturvidenskaben har fulgt i al den Tid,
den har bestaaet, og som ligger til Grund ved \emph{alle} de Resultater, den
har naaet. Det er da ikke nogen forhastet Paastand, at vi ikke kende nogen
Forklaring af et materielt Fænomen, som ikke bestaar i en Paavisning af et
andet materielt Fænomen som dets Aarsag. Nu er det en anerkendt Forskningsregel
(som Newton har udtalt i sin tredie og fjerde «regula philosophandi»), at de
Love, Erfaringen klart har godtgjort, maa vi antage som gyldige, indtil nye
Erfaringer indskrænke dem. Altsaa maa vi ogsaa antage Inertiens Lov som gyldig
for de materielle Fænomener, til hvilke Bevidsthedslivet er knyttet. Da vi nu
ikke kunne forklare nogen materiel Forandring i Hjernen af Sjælens Indgriben
uden at komme i skarp Strid med Inertiloven, er det en uundgaaelig Nødvendighed
at se os om efter en Hypotese, der ikke medfører en saadan Indgriben og altsaa
ikke bryder Inertiloven. Identitetshypotesen fyldestgør denne Betingelse, idet
den opfatter de materielle Bevægelser i Hjernen som ydre, sanseligt anskuelige
Former for hvad der foregaar i Sjælen.

Prof. Kroman udtrykker nu ganske vist Inertisætningen anderledes end Newton og
Maxwell, idet han blot lader den udsige, at «enhver Bevægelsesforandring
skyldes en Kraft»%
% footnote 1) on p. 83
\footnote[1]{Tænke- og Sjælelære. 2. Udg. p. 63.}. Han udelader altsaa Ordet
«ydre». Definitioner ere enhver Mands egen Sag. Men denne Ændring hænger sammen
med, at han opfatter Inertisætningen som et simpelt Korollar til den
almindelige Aarsagssætning%
% footnote 2) on p. 83
\footnote[2]{Vor Naturerkendelse. Kbhvn. 1883. p. 292 f.}. Den udsiger da kun,
at alle materielle Fænomener have deres Aarsag. Den særegne Betydning,
Inertisætningen har havt under den nyere Naturvidenskabs Udviklingsgang ---
særligt hver Gang en ny Hypotese opstilledes eller en ny Opdagelse gjordes, ---
den Betydning nemlig, at opfordre Forskerne til at spejde efter bestemte
\emph{materielle} Aarsager, den forstaar man ikke efter Kromans Opfattelse.
Dersom Inertisætningen i den Form, i hvilken Galilæi og Newton opfattede den,
ikke havde været den stadige (mere eller mindre klart bevidste) Forudsætning,
vilde Kant, Laplace og Darwin ikke være blevne førte til deres Hypoteser. ---
Men det er, som jeg tror, ogsaa med Rette gjort gældende mod Kromans Opfattelse
af Inertiloven som en blot Omskrivning af den almindelige Aarsagslov, at den
strider mod «saa temmeligt alle Fysikeres og Erkendelses%
% --- p. 84 ---
\setcounter{footnote}{0}%
\opage{84}teorikeres Overbevisning i vore Dage»%
% footnote 1) on p. 84
\footnote[1]{Vierteljahrsschrift für wissenschaftliche Philosophie. IX. p.
366.}. Efter Autoriteter kan naturligvis dette Spørgsmaal ikke afgøres. Hvad
jeg vil slaa fast, er kun, at den Opfattelse af Inertiloven, jeg har holdt mig
til, ikke blot er den naturlige, men ogsaa den, der har mest Betydning for de
specielle Problemer, og at det særligt er den eneste, der lader forstaa,
hvorledes Forholdet mellem Sjæl og Legeme i det Hele har kunnet blive et Problem. ---

Naar jeg her har holdt mig til Inertiloven, saa er det, fordi jeg mere og mere
bliver overbevist om, at det er i den, det afgørende Punkt ligger ved det hele
Problem. Denne Lov er --- naar den forstaas, som Galilæi og Newton forstod den
--- det typiske Udtryk for Naturvidenskabens Metode, og den indbefatter i sig
som i en Hovedsum alt hvad Naturvidenskaben har lært os om den materielle
Verden. Den maa da især komme i Betragtning, naar det gælder at fastslaa
Forholdet mellem Naturvidenskab og Aandsvidenskab.

I den nyeste Tid er det mest Loven om Energiens Bestaaen, man (første Gang
\emph{Fechner} i hans «Elemente der Psychophysik» 1860) er gaaet ud fra ved
Drøftelsen af Problemet om Forholdet mellem Sjæl og Legeme. Det er jo ogsaa
klart, at denne Lov ikke kan gælde for de materielle Fænomener, der antages at
bevirkes af eller at bevirke sjælelige Tilstande. Her vilde man imidlertid (som
jeg allerede har gjort opmærksom paa i min Psykologi. 2. Udg. p. 64) kunne
henvise til, at fordi den \emph{Retning}, i hvilken en Bevægelse gaar, ændres,
forøges derved ikke Energiens Sum. Diskussionen maa da her vende tilbage til
Inertiloven, som udtrykkeligt fordrer en ydre Aarsag ogsaa til Ændringer i
Bevægelsens Retning.

Imidlertid har naturligvis Loven om Energiens Bestaaen en stor Interesse ved
Behandlingen af vort Problem, og det er intet Tilfælde, at Diskussionen i de
sidste Decennier i saa høj Grad har drejet sig om den. Ved Forklaringen af et
Naturfænomen gælder det jo (ifølge Newtons regulæ philosophandi) at holde sig
til saadanne Aarsager, der ere paaviste i virkelig Erfaring og ere
\emph{tilstrækkelige} til at forklare Fænomenet. Det er da berettiget at gaa ud
fra, at de materielle Forandringer i Hjernen, ligesom alle de materielle
Forandringer, der ere blevne nøjagtigst undersøgte, underligge Loven om
Energiens Bestaaen. Vi kunne naturligvis ikke paastaa, at det er saa, førend
det i det Enkelte er paavist; men hvis vi i det Hele ville danne os en
Hypotese, er dette dog den nærmest liggende. Bevisbyrden synes at maatte
paahvile den, der her vilde statuere en Undtagelse. Det er en uvidenskabelig
Hypotese, som opbygges paa lutter saadanne Undtagelser.

Det vilde nu ganske vist være nødvendigt at statuere Undtagelser, dersom
virkelig Erfaring nødte os dertil. For en virkelig Iagttagelse kan ingen Teori,
ingen Lov og intet Princip bestaa. Men Sagen er jo netop den, at Iagttagelse
her er aldeles umulig. Ingen
% --- p. 85 ---
\setcounter{footnote}{0}%
\opage{85}kan \emph{iagttage}, hvorledes Forholdet er mellem det, der foregaar
i hans Hjerne, og det, der foregaar i hans Bevidsthed, og lige saa lidt kan
nogen \emph{iagttage} dette Forhold hos andre. De «Erfaringer», man her
beraaber sig paa, ere endnu mere raa end de, paa hvilke man beraabte sig, da
man paastod, at Erfaring lærte, at Solen gik rundt om Jorden. ---

Det synes klart, at enhver Opfattelse, hvilken Hypotese den saa end slutteligt
havner i, eller selv om den slet ingen vil opstille, maa tillægge de her
fremførte Betragtninger stor Vægt, da det dog i hvert Tilfælde er i højeste
Grad mærkeligt, at vi ere stillede saaledes, at vi ud fra den Viden, vi besidde
om den materielle Natur, ikke blot aldeles ikke kunne forstaa
Bevidsthedsfænomenernes Fremkomst, men endog ledes til den Konsekvents, at de
helt kunde undværes i Tilværelsen, da de \emph{tilstrækkelige} Aarsager til
alle materielle Fænomener maa findes indenfor selve den materielle Natur,
\emph{dersom} Naturvidenskaben bliver ved at skride frem paa samme Maade som
hidtil. At kaste et Slør over denne ejendommelige Omstændighed synes mig ikke
at være god Filosofi, da det er en af Filosofiens vigtigste Opgaver at bringe
os Konsekventsen af vor Viden til fuld og klar Bevidsthed. --- Den Mulighed, at
de hidtil uopklarede materielle Fænomener (til hvilke først og fremmest saa
mange af det organiske Livs vigtigste Fænomener høre) skulde nødvendiggøre helt
andre Principer og Love end dem, der hidtil udelukkende have staaet deres Prøve
som ledende Tanker, kan naturligvis ikke udelukkes. Men denne Mulighed er for
luftig en Grund at bygge nogen Hypotese paa. I de Huller, Videnskaben til
enhver Tid maa lade staa aabne, kan man bygge, hvad det skal være; der kan
Fantasien opføre sine Luftkasteller uden at frygte for at komme i Strid med
Videnskaben --- foreløbigt. --- Paa den anden Side gør den Omstændighed, at
disse Huller ere der, at vi paa dette Punkt ikke kunne naa mere end en
Hypotese, og vel aldrig kunne naa mere. Det gælder kun om ved vor
Hypotesedannelse at statuere saa faa Undtagelser som muligt fra de ellers
gældende Love, og allerhelst slet ingen at statuere. ---

Af den her saa skarpt fremhævede \emph{Forskel} mellem de sjælelige og de
materielle Fænomener følger nu ingenlunde, at de ere «\emph{absolut} uensartede»%
% footnote 1) on p. 85
\footnote[1]{\emph{Kroman}: Tænke- og Sjælelære. 2. Udg. p. 116. Efter først at
have gjort Identitetshypotesen til «Duplicisme», lader han den her avancere til
Dualisme.}. En for os irreduktibel Forskel mellem to Egenskaber eller
Ytringsformer behøver ikke at være en Forskel i det inderste Væsen: det er jo
netop dette, Identitetshypotesen hævder. Dersom vi kendte det Væsen, den
Grundvirksomhed, der paa saa forskellig Maade aabenbarer sig for vor indre og
for vor ydre Sans, vilde vi vel ogsaa kunne forstaa, hvorledes denne dobbelte
Fremtræden er nødvendig. En saadan Viden besidde vi ikke, og ville sikkert for
stedse være afskaarne fra at opnaa den. Men derfor have vi ingen Grund til at
mene, at der ligger
% --- p. 86 ---
\setcounter{footnote}{0}%
\opage{86}to forskellige Væsener til Grund. Svovlets Tilstandsform og Farve
variere efter Varmegraden; vi kunne ikke udlede Formen af Farven eller omvendt,
men vi kunne indse, at begge afhænge af Varmegraden. Ved Forholdet mellem de
sjælelige og de legemlige Fænomener kende vi ikke den sidste Aarsag til begges
Sammenhæng; men vi finde, at de variere paa tilsvarende Maade, uagtet de ikke
kunne udledes af hinanden, og slutte derfor, at de ere Ytringer af et og samme
Væsen. Den sædvanlige (monadologisk-spiritualistiske) Opfattelse gør
uberettiget Forskellen saa stor, at den betragter Sjæl og Legeme som to
forskellige Væsener. Den fremhæver Modsætningen mellem Sjæl og Legeme, hvor der
ingen Grund er til det, og oversér den dér, hvor den store Vanskelighed netop ligger.

Identitetshypotesen gør Front til to Sider. Den hævder overfor Materialismen,
at man ikke kan udlede Bevidsthedsfænomenerne af Materien efter det Begreb,
Naturvidenskaben giver os om denne. Den hævder overfor Spiritualismen
(Descartes og Lotze), at der i denne Irreduktibilitet ikke ligger nogen Grund
til at sætte en Grænse for de naturvidenskabelige Principers Gyldighed for alle
materielle Fænomener, den menneskelige og dyriske Hjerne indbefattet. For at
forstaa, hvorledes Bevidsthedslivet kan blive til paa et vist Trin af materiel
Udvikling, maa man da antage, at der ved al materiel Virksomhed, foruden de
Egenskaber og Ytringer, Naturvidenskaben konstaterer, ogsaa medvirker en for
ydre Sansning ikke fremtrædende Egenskab, og at det er af denne, Bevidstheden
bliver til. Man kan da tænke sig den i Analogi med Bevidstheden, eller som et
Sjæleliv af forsvindende Intensitet i Sammenligning med det Sjæleliv, vi kende
fra vore egne Bevidstheds Fænomener. Naar Stenkastet fremkalder Fornemmelse og
Smerte, saa er det altsaa egentligt \textit{a} + \textit{α}, som fremkalder
\textit{b + β,} idet \textit{a} betyder den fysiske Bevægelse ved Stenkastet,
\textit{α} det dertil svarende psykiske Analogon, \textit{b} den ved Kastet
frembragte Hjerneproces og \textit{β} den til denne svarende
Bevidsthedstilstand (Fornemmelse og Smerte). --- Ved at opfatte Sagen paa denne
Maade bevare vi Kontinuiteten i vor Naturopfattelse, samtidigt med, at vi
anerkende de sjælelige Fænomeners Ejendommelighed overfor de materielle.

Naturvidenskabens Principer og Love lære os efter denne Hypotese kun
Tilværelsens ydre, for sanselig Anskuelse aabne Side at kende, men Tilværelsen
har ogsaa en indre Side, der kundgør sig for os i Selviagttagelsen, og som vi
overalt, hvor Selviagttagelsen ikke naar hen, kunne tænke os i Analogi med hvad
den viser os. Vi maa udvide vort Materiebegreb, lægge mere ind deri end
Naturvidenskaben, der stedse kun bygger paa den ydre Sans, har Anledning til.
En saadan Udvidelse er nødvendig, hvis de sjælelige Fænomeners Opstaaen skal
blive forklarlig. Denne Udvidelse af Materiebegrebet rokker aldeles ikke ved
dettes Gyldighed og Anvendelse og synes at være mere berettiget end den
Antagelse, at en materiel Bevægelse pludseligt --- i Strid med alt hvad der kan
sluttes ud fra de hidtil faststaaende Principer og Love --- skulde afbrydes og
afløses af en sjælelig Proces.

% --- p. 87 ---
\setcounter{footnote}{0}%
\opage{87}36. b) Ved den sidste Betragtning have vi allerede forberedt Svaret
paa den anden Indvending mod Identitetshypotesen. Denne Indvending gik ud paa,
at vor Viden om den legemlige Verden blev uforklarlig, naar der ikke findes
Vexelvirkning mellem Sjæl og Legeme.

Denne Indvending har undret mig, da det forekommer mig, at selve Hypotesens
Ordlyd, eller dens blotte Navn (naar man ikke omdøber den) indeholder et Svar.
Naar der skal være Identitet af Bevidsthed og Hjernevirksomhed, saaledes at
Hjernevirksomhed er det ydre, sanselige Udtryk for hvad der samtidigt gaar for
sig i Bevidstheden, saa er det klart, at naar en vis Hjernevirksomhed
(\textit{b}) foraarsages ved en vis fysisk Paavirkning (\textit{a}), opstaar
der i samme Øjeblik foruden \textit{b} ogsaa en Bevidsthedstilstand
(\textit{β}), og at hvis \textit{b} svarer til \textit{a,} maa det i sin
inderste Grund med \textit{b} identiske eller sammenhørende \textit{β} ogsaa
svare dertil. Enten jeg, med den sædvanlige spiritualistisk-dualistiske
Opfattelse siger, at \textit{a} frembringer \textit{b,} og at \textit{b} derpaa
frembringer \textit{β}, eller jeg, med Identitetshypotesen, siger, at
\textit{a} frembringer \textit{b,} som er i sin Grund identisk eller%
% footnote 1) on p. 87
\footnote[1]{Hvilken Udtryksmaade man vil bruge, er ligegyldigt. \textit{b} og
\textit{β} antages at være \emph{fænomenalt forskellige}, men i \emph{deres
Grund identiske}. Efter det første Synspunkt kunne vi bruge Formlen \textit{b}
+ \textit{β}, efter det andet Formlen \textit{b = β.}} sammenhørende med
\textit{β}, maa i denne Sammenhæng komme ud paa Et. \emph{Fælles for begge
Opfattelser} er, at \textit{b} ikke \emph{fremkommer}, \emph{uden at} \textit{β
ogsaa fremkommer}. Om man forklarer Forholdet mellem \textit{b} og\textit{ β}
som et Kausalitetsforhold eller som et Identitetsforhold, gør i denne
Sammenhæng aldeles ingen Forskel. Identitetshypotesen har dog ogsaa her et
Fortrin; ti medens Spiritualismen maa statuere to Overgange: fra \textit{a} til
\textit{b,} og fra \textit{b} til \textit{β}, behøver Identitetshypotesen kun
én Overgang: fra \textit{a} til \textit{(b + β). ---} Den i 35 (mod Slutningen)
udviklede Opfattelse viser os, hvorledes vi kunne forstaa, at \textit{a} kan
frembringe ikke blot \textit{b,} men (\textit{b + β}); det bliver nemlig kun
muligt under den Forudsætning, at \textit{a} i Virkeligheden betyder
(\textit{a} + α)\textit{. ---} Det er altsaa for Identitetshypotesen meget vel
muligt at fastholde Kontinuiteten mellem den materielle Omverdens Fænomener og
de Bevidsthedsfænomener, der svare til dem. Indvendingen falder bort, saasnart
man tilstrækkeligt har tænkt sig ind i Hypotesen. ---

Men foruden denne \emph{psykologisk}-\emph{fysiologiske} (psykofysiske)
Betragtning giver Indvendingen ogsaa Anledning til en
\emph{erkendelsesteoretisk} Betragtning, som er af Betydning i denne
Sammenhæng. --- Alt hvad vi vide om den materielle Verden, vide vi fra vore
Fornemmelser. At disse Fornemmelser igen skrive sig fra den materielle Verden,
er en Paastand, der vel er plausibel fra et populært-dogmatisk Standpunkt, men
som fra et erkendelsesteoretisk Synspunkt betragtet grunder sig paa en
Cirkelslutning. At vi --- saa fremt vi hylde Aarsagsprincipet --- ikke kunne
mene, at vore Fornemmelser opstaa «af sig selv», eller at vi selv frembringe
dem af intet (d. v. s. uden \emph{given} Anledning), følger af sig selv. Men
det er, erkendelsesteoretisk set, selvmodsigende at erklære den materielle
Verden, som
% --- p. 88 ---
\setcounter{footnote}{0}%
\opage{88}vi kun kende fra vore Sansefornemmelser, for Aarsag til disse. Vore
Erkendelsesobjekter ere Bestanddele af selve vor Erkendelse, af vor erkendende
Bevidsthed, og Berettigelsen til at antage, at der existerer Ting med samme
Kvaliteter, som de fremtræde med i vor Bevidsthed, denne Berettigelse skal
først godtgøres. Under Indflydelse af de Paavirkninger, det erkendende Subjekt
modtager, opstaar der hos det Fornemmelser og Forestillinger, som efter
Bevidsthedens Love bearbejdes saaledes, at Billedet af den udstrakte, farvede
og tonende Verden danner sig. Men i dette Billede have vi ingenlunde Ret til
uden videre at se et adækvat Billede af de Aarsager, hvis Indflydelse paa
Subjektet sætter og holder dettes Erkendelsesvirksomhed i Gang.

Den omtalte Indvending synes uden videre at gøre Erkendelsens \emph{Genstand}
(Objektet) til Et med dens \emph{Aarsag}. Vore Objekter ere de Billeder, der
danne sig i Bevidstheden under Indflydelse af visse virkende Aarsager; men
umiddelbart have Erkendelsens Genstande og Aarsager intet at gøre med
hverandre. For at bruge et (allerede af Lotze anvendt) Billede: vor Erkendelse
af den materielle Verden kan sammenlignes med Gnisten, der slaar ud, naar
Staalet slaar mod Flintestenen. Men Gnisten ligner ikke Staalet, og det vilde
dog være højst urimeligt at paastaa, at man ikke kunde forstaa Gnistens
Fremkomst af Flinten, naar det ikke var en Gnist, der paavirkede denne! ---
Hele Spørgsmaalet, som det er stillet i den nævnte Indvending, forekommer mig
at være urigtigt stillet. Det beror paa en Forvexling af det
erkendelsesteoretiske Forhold mellem Subjekt og Objekt og det psykofysiske
Forhold mellem Bevidsthed og Materie. Disse to Forhold maa vel holdes ude fra
hinanden%
% footnote 1) on p. 88
\footnote[1]{I mit Skrift «Spinozas Liv og Lære» (Kbhvn. 1877) p. 100 f. har
jeg paavist, hvorledes denne Forvexling fremtræder hos Spinoza. Ogsaa
\emph{Windelband} (Geschichte der neueren Philosophie. I. Leipzig. 1878. p.
211) og \emph{Tönnies} (Studie zur Entwickelungsgeschichte der Spinoza.
Vierteljahrsschrift für wissenschaftliche Philosophie. 1883. p. 176 f.) have
gjort denne Bemærkning. --- Hos Schelling og Hegel spiller denne Forvexling en
skæbnesvanger Rolle.}. Det første Forhold vedrører al vor Erkendelse, ikke blot
Erkendelsen af den materielle Verden; det andet Forhold er væsentligt et
Grænseforhold mellem Psykologi og Fysiologi. Dette sidste er altsaa langt mere
specielt end hint, og desuden af en helt anden Karakter.

I vor specielle Forsken og vor praktiske Bevidsthed holde vi os til
\emph{Erkendelsens Genstande} og undersøge deres indbyrdes Sammenhæng. Vi danne
da paa Grundlag af vore Fornemmelser et saa nøjagtigt Billede af den materielle
Verden som muligt. Men ogsaa de sjælelige Tilstande hos os selv og andre ere
Erkendelsesgenstande, skønt de ikke fremtræde for ydre Sanseánskuelse eller i
Rummets Form. Naar vi nu ville afgøre, hvorledes vi indenfor vort objektive
Verdensbillede skulle opfatte Forholdet mellem det Materielle og det Sjælelige,
kunne vi kun bygge paa hvad paa den ene Side Naturvidenskaben, paa den anden
Side Psykologien lærer os, og vi naa til vor Hypotese ved at drage
% --- p. 89 ---
\setcounter{footnote}{0}%
\opage{89}de sidste Konsekventser af disse Lærdomme. Ifølge Identitetshypotesen
f. Ex. opfattes Forholdet saaledes, at de sjælelige Tilstande forholde sig til
de tilsvarende Hjerneprocesser paa samme Maade som en Cirkelbues konkave Side
forholder sig til dens konvexe. Tilværelsen kommer altsaa til at staa for os
som en Bue, der kan ses fra to Sider. --- Dette er den psykofysiske Betragtning.

\emph{Erkendelsesteoretisk} spørge vi derimod om, paa hvilket givet Grundlag og
ved Hjælp af hvilke Principer vi have dannet dette Verdensbillede, som er vor
Erkendelses Genstande, ordnede i en bestemt Sammenhæng. Det givne Grundlag er
tilsidst vore Fornemmelser, hvorimod disses Aarsag, og med den Aarsagen til
vort hele psykofysiske Verdensbillede, aldrig selv er \emph{given}, men kun
findes ved Hjælp af en metafysisk Hypotese. Ved Bestemmelsen af denne Aarsag
træde de forskellige metafysiske Systemer i Modsætning til hverandre, idet
ethvert af dem bestemmer den paa sin Maade. I Psykologi og Fysiologi se vi bort
fra dette Problem og indordne vore Fornemmelser som Led i den
Verdenssammenhæng, vi opbygge paa Grundlag af selve disse Fornemmelser. Men
erkendelsesteoretisk er vor Fornemmelse (lige saa vel som vor Tænknings
Principer) et Sidste, en oprindelig Forudsætning, der naturligvis ikke kan
forklares ved det, der opbygges paa Grundlag af selve denne Forudsætning. ---

Til Afvisning af den omtalte Indvending er det tilstrækkeligt at holde sig til
den ovenanførte psykofysiske Betragtning. Men da Spørgsmaalet kan have en
større Rækkevidde og let kan give Anledning til Misforstaaelse, er jeg gaaet
ind paa den erkendelsesteoretiske Side af Sagen.

37. c) Problemet om Forholdet mellem det Sjælelige og det Legemlige vilde ikke
høre til de største og mest centrale Problemer for menneskelig Forsken, hvis
nogen Hypotese allerede nu kunde antages at have givet en endelig Løsning. Vi
kunne ikke naa videre end at lodse os frem i dette Farvand, som er opfyldt af
saa mange Skær, idet vi som Sømærker bruge de Principer og Resultater, der
hidtil have staaet deres Prøve i vor Erkendelse. Hvad særligt
Identitetshypotesen angaar, saa var det kun den dogmatisk-metafysiske Tankegang
hos dens store Grundlægger, der kunde føre til den Tro, at Problemet ved den
fuldstændigt var løst. Der bliver stadigt Stof nok tilbage for den Forundring,
som sætter Tanken i Bevægelse. For mit Vedkommende har jeg ganske vist den ---
som jeg tror, ikke ubegrundede --- Overbevisning, at der her er et skarpt
Hjørne, en Station, som vor Verdensanskuelse maa passere. Den maa her, som paa
saa mange andre Punkter, i sin Opfattelse af Tilværelsen forlade den
tilsyneladende umiddelbart givne Betragtningsmaade. Men tilbage staar stedse
den store Mærkelighed, at Tilværelsen frembyder to saa forskellige Sider: paa
den ene Side den \emph{inderlige Enhed} i \emph{Bevidstheden}, paa den anden
Side den \emph{ydre Mangfoldighed} i \emph{Materien}.

% --- p. 90 ---
\setcounter{footnote}{0}%
\opage{90}Det maa her først og fremmest mærkes, at den Vanskelighed, som heri
kan ligge, er fælles for alle Teorier, og ikke mindst gør sig gældende i den
sædvanlige spiritualistiske Teori. Ti hvorledes vil man kunne forbinde nogen
anskuelig Forestilling med den Vexelvirkning, der antages at finde Sted mellem
det udstrakte Mangfoldige og den ikke-rumlige Enhed? Gaaden bliver her i det
Mindste lige stor, hvad enten man saa antager et Aarsagsforhold eller et
Identitetsforhold. Og Gaaden bliver egentligt i begge Tilfælde den samme. Al
Kausalitet er Kontinuitet. Hvad vi kalde Aarsag og Virkning, staar for den
videnskabelige Erkendelse ikke som to forskellige Ting, men som to Led i én og
samme kontinuerlig Proces. Alle videnskabelige Teorier gaa ud paa at formulere
en saadan Kontinuitet mellem Fænomener, der for den sanselige Opfattelse staar
som forskellige Ting eller Begivenheder. Atomteorien f. Ex. gør det muligt at
se bort fra de kvalitative Forandringer og den dermed følgende Diskontinuitet i
vor sanselige Erfaring om Stoffernes Forandringer og danner Forestillingen om
en Omsætnings- og Omlejringsproces, ved hvilken materielle Smaadele
kontinuerligt bevæge sig fra det ene Punkt i Rummet til det andet. Dersom der
nu bestod et Kausalitetsforhold mellem legemlige og aandelige Fænomener, maatte
en og samme kontinuerlig Proces føre fra Materiens Mangfoldighed til
Bevidsthedens Enhed. Jeg vil nu ikke stille den Fordring til den populære
Spiritualismes Talsmænd, at de skulle vise os, hvorledes denne Overgang gaar
for sig i Virkeligheden. Jeg vil blot spørge, om vi kunne danne os nogen som
helst Forestilling om dens \emph{Mulighed}. Et Brud paa Kontinuiteten maa der i
hvert Tilfælde komme dér, hvor den fysiologiske Proces hører op for at afløses
(fortsættes kan man ikke kalde det) af den psykiske. Lad os tænke os
Resultanten af de fysiske og kemiske Processer i Hjernen som en ret Linie.
Denne Linie maa da pludseligt afbrydes; en Iagttager vilde fra et vist Øjeblik
af intet se mere. Den rumlige Bevægelse vilde pludseligt være «absorberet». ---
Et saadant Brud paa Kontinuiteten medfører Identitetshypotesen ikke.

Den Vanskelighed, der bliver tilbage for Identitetshypotesen, mener jeg
allerede i min tidligere Fremstilling at have fremhævet saa stærkt som muligt%
% footnote 1) on p. 90
\footnote[1]{Smlgn. Psykologi. 2. Udg. p. 76, 96, 98, 409.} Den materielle
Verden fremtræder for os (i hvert Tilfælde i vore Hypoteser) med en ganske
anden Kontinuitet, end vi formaa at tænke os for den sjælelige Verdens
Vedkommende. Paa den anden Side fremtræder til Gengæld Individualiteten, den
sluttede Enhed og Helhed ganske anderledes udpræget paa Sjælelivets Omraade,
selv om det vilde være Dogmatisme her at tale om absolut Enhed, da Enheden i
vor Erfaring fremtræder i alle mulige Grader lige fra den højeste Koncentration
til saadanne Fænomener, hvor Enheden næsten synes tabt i de rent sporadisk
opdukkende Bevidsthedselementer. Inden for hver enkelt individuel Enhed er det
igen forbundet med Vanskelighed at hævde Kontinuiteten; der svinges mellem
% --- p. 91 ---
\setcounter{footnote}{0}%
\opage{91}Bevidsthed og Ubevidsthed, og de psykiske Elementer synes at
forsvinde, uden at Ækvivalenter kunne paavises. Vi have uovervindelige
Vanskeligheder ved at tænke os en Lov om Energiens Bestaaen for den sjælelige
Tilværelses Vedkommende. Baade indenfor den enkelte Bevidsthed og mellem de
forskellige Bevidstheder synes Kontinuiteten brudt, og det ovenfor (35 og 36)
antydede hypotetiske Supplement naar ikke videre end til en højst vag Antydning
af noget med de psykiske Elementer Analogt, der kunde udfylde Lakunerne. Hvis
vi i vor Erkendelse af Tilværelsen fuldstændigt formaaede at forbinde disse to
Synspunkter: Fænomenernes Kontinuitet og deres Individualitet, saa vilde
egentligt alle Gaader være løste.

Det forekommer mig at være et Fortrin ved Identitetshypotesen, at den stiller
baade Sammenhængen og Modsætningen mellem den materielle og den sjælelige Natur
i saa klar og simpel Belysning. Den slutter sig til den barnlige og praktiske
Opfattelse, der betragter Legemet som det ydre Udtryk for Menneskets
Personlighed, medens Legemet efter Spiritualismen er et fra Menneskets
Personlighed forskelligt Væsen. I videnskabelig Henseende tillader den
Fysiologien og Psykologien at gennemføre hver sin Metode og dog at arbejde
Haand i Haand med hinanden. Den foregriber intet Resultat. Den paastaar aldeles
ikke, at alt i Hjernen er forklaret efter Inertiens Lov og Loven om Stoffets og
Kraftens Bestaaen. Men den paastaar, at intet videnskabeligt Resultat hidtil er
naat, som ikke bekræftede disse Love, og at en Hypotese, der begynder med at
afvise dem, derfor ikke kan have nogen Berettigelse eller Betydning. ---

Den sjælelige Aktivitet mister ikke sin reale Betydning i Verden, fordi den i
sit Væsen er identisk med en fysisk Aktivitet. Den faar lige saa meget at sige
som den Nerveproces, der er dens for ydre Sansning tilgængelige Udtryk.
Udviklingshistorien viser os, at Nervesystemet er en af de Dele af Organismen,
der først med Tydelighed fremtræde i Anlægget, og at det fra først af indtager
en central Stilling i Forhold til de andre Dele af Organismen%
% footnote 1) on p. 91
\footnote[1]{Se herom: \emph{Kölliker}: Entwickelungsgeschichte des Menschen
und der höheren Thiere. Leipzig 1861. Vorlesung VII, VIII, XXV. --- R. S.
\emph{Berg}: Forelæsninger over den almindelige Udviklingshistorie. København
1887. p. 102 --- 104.}. Der er her et oprindeligt virkende Anlæg, som viser, at
Nervesystemets Funktioner, hvis Korrelata de sjælelige Virksomheder ere, ikke
ere blotte Virkninger af ydre Indflydelse, men bestemte i deres Natur og
Retning ved Individets første Tilblivelse.

Om vi i nogen af de to Ytringsformer have Tilværelsens absolute Natur given, er
et Spørgsmaal, som ikke kan besvares. Aand og Materie staa ganske vist som
Hovedforskelligheder i vor Erfaring; vi kende intet Tredie foruden dem. Men det
er en blot faktisk Tohed, vi her staa overfor. Der er intet kontradiktorisk
Forhold mellem dem, saa at en tredie eller en fjerde Ytringsform skulde være en
logisk Umulighed. Heller ikke er der, som ved Rummets tre Dimensioner, nogen
Anskuelsesumulighed til Stede for Antagelsen af
% --- p. 92 ---
\setcounter{footnote}{0}%
\opage{92}flere Former. Det er derfor ingen løs Paastand, men en Bemærkning,
der kaster et Lys over vor Erkendelses faktiske Begrænsning, naar vi henvise
til, at Løsningen af Tilværelsesproblemet kan ligge langt dybere end baade de
spiritualistiske og de materialistiske Teorier tænke sig det. Det vilde være
kortsynet Dogmatisme, naar man her paa dette Grænsepunkt ikke havde et aabent
Øje for de forskellige Muligheder, der frembyde sig. Kun langsomt og
modstræbende har den menneskelige Verdensanskuelse forsonet sig med Tanken om
det fysiske Univers's Uendelighed; den vægrer sig paa lignende Maade ved at
indrømme, hvor lidet vi formaa at gennemskue Tilværelsens indre Sammenhæng, men
er tilbøjelig til at mene sig tilstrækkeligt udstyret til at kunne fastslaa en
afsluttende Erkendelse.

\begin{center}\large VII. Den sjælelige Aktivitets Lovmæssighed.\end{center}

38. For Psykologien som for enhver real Videnskab er Aarsagsprincipets
Anerkendelse en Livsbetingelse. Kun naar det anerkendes, at de sjælelige
Fænomener have deres Aarsager, blive de psykologiske Problemer til. I hvert
Tilfælde vilde Psykologien aldrig kunne konstatere, at Aarsagsloven ikke
gjaldt; den vilde kun kunne slaa fast, at vi endnu ingen Aarsager kunde finde
til visse Fænomener.

Dog vilde det kunne tænkes, at den psykologiske Forsken stødte paa Fænomener,
hvis Virkelighed den maatte anerkende, men hvis ejendommelige Natur netop
betingedes ved, at Aarsagsloven ikke gjaldt, eller ikke fuldstændigt gjaldt,
saaledes at man fik Valget imellem at forkaste saadanne Fænomeners Virkelighed
og at anerkende Aarsagslovens Begrænsning. Da nu et Faktum altid har Forrangen
for et Princip, kunde Valget ikke blive tvivlsomt.

Et Fænomen af denne Art anses af mange de moralske Følelser (særligt Anger-
eller Skyldfølelsen) for at være, og man stiller da det Valg enten at forkaste
saadanne Følelser som grundede paa Illusion eller at antage Aarsagsløshed. Der
skal nemlig med Angerfølelsen være forbundet en uafviselig Forestilling om, at
man i det givne Tilfælde \emph{kunde have} villet anderledes end man vilde, og
denne Forestillings Gyldighed skal være den Betingelse, til hvilken
Angerfølelsens Bestaaen er knyttet.

En Paastand som denne forvexler Psykologiens Standpunkt med Etikens. Psykologen
spørger blot om Bevidsthedsfænomenernes faktiske Opstaaen og Sammenhæng. Han
undersøger da ogsaa hin Følelses Natur nærmere. Først undersøger han, om den i
den angivne Form, bunden til Forestillingen om at have kunnet ville anderledes,
\emph{findes} hos alle Mennesker. Da der nu faktisk findes dem, som vel kende
Anger, men aldrig have havt den klare Bevidsthed om, at de havde kunnet ville
anderledes naar hele deres Sinds%
% --- p. 93 ---
\setcounter{footnote}{0}%
\opage{93}tilstand ikke var bleven ændret, --- saa slutter Psykologen, at hvor
hin Følelse findes i den angivne Form, skylder den rimeligvis ganske særegne
Betingelser sin Opstaaen, og han vil da søge en Forklaring af, hvorledes en saa
mærkelig Forestilling, som den omtalte, kan have sat sig fast i Bevidstheden og
forbundet sig med Angerfølelsen. Ja, selv om den omtalte Form for Angerfølelse
fandtes hos alle Mennesker, vilde Psykologen som teoretisk Forsker ikke kunne
stille sig anderledes. De Forestillinger, der forbinde sig med Følelserne og
give disse deres specielle Karakter, behøve nemlig ingenlunde at indeholde den
\emph{virkelige} Aarsag til Følelserne. Det er tilstrækkeligt, naar Individet,
som har Følelsen, betragter den dertil knyttede Forestilling som gyldig; saa
længe han gør det, vil Følelsen bevare sin Karakter. I Følelsernes Psykologi
prøves ikke Gyldigheden af de Forestillinger, der bestemme Følelsens Karakter
og Retning; naar der blot indgaas en Association af Forestilling og Følelse, er
det rent psykologisk ikke af afgørende Betydning, om Forestillingen er gyldig
eller ikke.

39. Ad flere forskellige Veje vilde man kunne vise, hvorledes Forestillingen om
en Villiesakts Aarsagsløshed kan opstaa.

For det Første kan der henvises til, at jo mere en sjælelig Tilstand lægger
Beslag paa vor Interesse og Opmærksomhed, des mindre lægge vi Mærke til de
forudgaaende og betingende Omstændigheder. I en stærkt bevæget Følelsestilstand
tænke vi ikke paa, at den paa mange Maader er forberedt og indledet ikke blot
ved bestemte ydre Erfaringer, men ogsaa ved forudgaaende Bevidsthedstilstande
dels af beslægtet, dels af forskelligartet, maaske endogsaa modsat
Beskaffenhed. Den hele Tilstand kommer til at staa for os som pludseligt, af
sig selv indtraadt, tilbleven af intet, og vi ere utilbøjelige til at gaa ind
paa nogen som helst genetisk Forklaring af den. Aller mest gælder dette ved en
alvorlig Villiesafgørelse. Naar efter en lang indre Debat den afgørende
Beslutning skyder sig frem og i ét Øjeblik fejer alle modstaaende
Betænkeligheder bort, da synes der at foreligge et Brud paa Kontinuiteten. Den
nye Tilstand, i hvilken vort hele Sind er samlet i Tanken om Handlingen, staar
i en saadan Modsætning til Overvejelsens indre Strid og Splid, at selv om vi
bevare en levende Erindring om denne, have vi dog Vanskelighed ved at fastholde
den Erkendelse, at der er et Kausalitetsforhold mellem de to Tilstande.
Achilleus's pludselige og afgørende Selvbeherskelse i det heftige Optrin
overfor Agamemnon (i Iliadens første Bog) kunde de gamle Grækere kun forklare
sig ved, at Atene kom ned fra Olymp og --- kun synlig for den vrede Helt ---
trak ham i hans Lokker og holdt ham tilbage, da han allerede var begyndt at
drage Sværd af Skede. Den gammelgræske Bevidsthed appellerede i saadanne
Tilfælde til en guddommelig Indgriben; en moderne Indeterminist vilde her have
appelleret til «den frie Villie».

For det Andet kunde man henvise til, at Villiesafgørelsen, Beslutningen, medfører
% --- p. 94 ---
\setcounter{footnote}{0}%
\opage{94}en ejendommelig Fornemmelse af Uhæmmethed og Kraft paa Grund af, at
Sindet næsten fuldstændigt gaar op i én Tanke. Al indre Modstand, al Tvang er
falden bort. Alle Tanker og Følelser gaa i samme Retning. Ingen Hensyn eller
Skranker agtes. Der sker i Beslutningen et Gennembrud af vort reale Selv, af
den Kreds af Bevidsthedselementer, der danner den relativt konstante Del af vor
Bevidstheds Indhold, og vi have derved samme Fornemmelse som ved et fuldt og
dybt Aandedrag, en Fornemmelse af, at vi nu kom ret til os selv, en potentseret
Selvfølelse, der skyder al Tanke om, at denne Tilstand dog selv kun blev til
som Følge af andre Tilstande, langt i Baggrunden. Vi føle os suveræne,
ubetingede. Men al Foraarsagethed er Betingethed, Begrænsning, Bundethed. Intet
Under, at vi ikke i samme Øjeblik kunne gaa helt op i den nærværende Tilstand
og klart fastholde Forestillingen om dens Afhængighed af forudgaaende Tilstande!

For det Tredie vil Erindringen om Overvejelsens Tilstand --- \emph{hvis} den
kan gøre sig gældende --- kunne føre til, at man antager den objektive Mulighed
af en anden Beslutning end den virkeligt fattede. Der fremstillede sig under
Overvejelsen forskellige mulige Handlinger for vort Blik. Vi satte os ind i
enhver af dem og reves kun bort fra dem hver for sig, fordi Forestillingen om
de andre gjorde sig gældende. Men Forestillingen om en forkastet Mulighed kan
være traadt saa levende frem, at den næsten har faaet en sjælelig Virkeligheds
Præg for os. (Smlgn. hermed 32). Vi drage os jo ofte selv til Ansvar, blot
fordi vi have kunnet \emph{tænke} paa at foretage en vis Handling. Og i den
levende Forestilling om en Handling ere vi --- da Aktiviteten ikke begynder paa
et vist bestemt Punkt --- altid et Stykke paa Vej til at udføre Handlingen. En
saadan Handling kommer da i den levende Erindring ikke til at staa som en blot
subjektiv, bleg Mulighed; vi tillægge den en vis Objektivitet --- og dermed er
Illusionen om en real Evne til i det enkelte Øjeblik at ville anderledes, end
vi virkeligt ville, stærkt forberedt. Jo klarere og intensivere den forkastede
Mulighed staar Side om Side med den foretrukne, des større Fristelse er der til
at tillægge dem omtrent samme Gyldighed. Det Billede, der naturligt skyder sig
frem, viser os to Veje; vi gaa kun ad den ene; men den Vej, vi ikke gaa ad,
bliver dog ved at staa for os som en objektivt given Vej, en Realitet. Dette
Landevejsbillede overføre vi paa de subjektive Muligheder, vi have kæmpet med.

For det Fjerde. Selv om Angeren ikke altid betinges af Forestillingen om at
\emph{have kunnet ville} anderledes, saa indeholder den dog, hvor den rører sig
med Styrke, et levende Ønske om, at man \emph{havde villet} anderledes. Den
Tilstand, i hvilken man i Angerens Øjeblik ser tilbage paa sin Handling, er
højst forskellig fra den Tilstand, i hvilken man fattede Beslutningen om at
udføre Handlingen. Man har Møje med at identificere de to Jeger, der ytre sig i
de to Tilstande. Selvgenkendelsen er vanskelig. Identifikationen sker da let
paa den Maade, at man overfører noget fra den nuværende Tilstand til den
tidligere. Man tillægger da det tidligere Jeg paa én Gang den Kraft, der
% --- p. 95 ---
\setcounter{footnote}{0}%
\opage{95}førte til den forkastelige Beslutning, og den Kraft, der i det
nærværende Øjeblik fører til at forkaste den. Hvad der i det tidligere Øjeblik
var et Enten---Eller, gør man ved en ejendommelig Sammenknytning til et
\emph{Baade}---Og. Medens Angeren ikke nødvendigvis \emph{forudsætter}
Forestillingen om at have kunnet ville anderledes, kan det altsaa vises, at den
kan \emph{frembringe} en saadan Forestilling ved en ganske naturlig psykologisk
Proces. Ti saadanne Forskydninger foregaa uafbrudt i vort Bevidsthedsliv.

Endeligt, for det Femte, kan en lignende Illusion som den sidst omtalte gøre
sig gældende, naar man under en alvorlig Overvejelse tænker paa Fremtiden%
% footnote 1) on p. 95
\footnote[1]{Smlgn. herom \emph{Darlu} i Revue philosophique. Juin 1887. p.
574---576.}. Naar jeg saaledes ser fremad, staar det klart mig, at Fremtiden
--- for det Omraade, min Villie har Indflydelse paa --- bliver en anden, alt
efter som jeg vælger den ene eller den anden af de Handlinger, jeg drøfter.
Fremtiden kommer altsaa til at staa for mig i et dobbelt Lys: dels saaledes som
den vilde være som Følge af en vis bestemt Handling, dels saaledes som den
vilde være uden denne. Der finder da her let en Sammenskydning, en Art
Interferents Sted af de to Fremtidsbilleder. Det heraf resulterende
Fremtidsbillede kommer til at staa ubestemt, skiftende, vaklende for mig. Det
faar Udseendet af, at Fremtidens Begivenheder i sig selv ere indeterminerede,
fordi vi intet klart bestemt Billede af dem kunde fastholde, idet andre
Billeder gribe forstyrrende ind. Her gør sig en stadig Oscilleren gældende. ---
Denne Oscilleren, som hos mange ytrer sig med Hensyn til den Del af Fremtiden,
deres Villen kan bestemme, ytrede sig for Oldtidens Tænkere, selv den største
af dem, ogsaa hvad Fremtiden i det Hele angaar. \emph{Aristoteles} lærte, at
der til vore problematiske Domme svare objektive Muligheder i Naturen, ligesom
der til vore apodiktiske Domme svare nødvendige Forhold i Naturen. Verdensløbet
var for ham ubestemt, hvad Universets lavere Sfærer (den sublunariske Verden)
angaar. I Skriftet de interpretatione (cap. 9), som i hvert Tilfælde er blevet
\emph{tillagt} Aristoteles, læres det, at den Sætning, at af to kontradiktorisk
modsatte Domme maa den ene være sand, den anden falsk, kun gælder for Fortiden
og Nutiden, ikke for Fremtiden. Naar vi sige: enten vil der staa et Søslag
eller ikke, saa er det --- efter denne Lære --- \emph{hverken} nødvendigt, at
der vil staa et Søslag, \emph{eller} at der intet vil staa. Ingen af Delene kan
i Forvejen være sand. Hvad Fortiden og Nutiden angaar, ere Mulighederne
udtømte, men hvad Fremtiden angaar, staa de, ikke blot for den forskende
Tænkning, men i selve den reale Tilværelse (hvorledes man saa nærmere har tænkt
sig dette!) endnu aabne!%
% footnote 2) on p. 95
\footnote[2]{Smlgn. om denne aristoteliske Lære \emph{George Grote}: Aristotle.
I. p. 183 ff. --- I den nyeste Tid har W. \emph{Wundt} hævdet, at Principet om
Ækvivalents af Aarsag og Virkning ikke gælder for den sjælelige Kausalitets
Vedkommende, men kun for den fysiske Kausalitet. Han opstiller i Modsætning til
Ækvivalentsprincipet Principet om den stedse voxende sjælelige Energi. Han
sætter denne Antagelse i Forbindelse med den Paastand, at vor Kausalforklaring
paa det sjælelige Omraade kun angaar Fortiden, ikke Fremtiden. (Ethik. Leipzig
1886. p. 399 f. --- Physiologische Psychologie. 3. Aufl. II. p. 183 f.) Det kan
naturligvis ikke negtes, at Aandsvidenskaben i langt ringere Grad end
Naturvidenskaben formaar at forudsige Fremtiden. Men er det ikke det
Naturligste at udlede denne Omstændighed af vor højst mangelfulde og lidet
exakte Viden paa det aandelige Omraade? Fordi meteorologiske Forudsigelser
langtfra kunne maale sig med astronomiske Forudsigelser, mene vi jo dog ikke,
at der paa det meteorologiske Omraade gør sig en anden Art Kausalitet gældende
end paa det astronomiske Omraade. Vor relative Uvidenhed paa det psykologiske
Omraade har desuden sin Parallel i en ikke mindre Uvidenhed paa
Hjernefysiologiens Omraade, saa at hvad der gælder for de aandelige Fænomener,
maatte (efter Identitetshypotesen, som Wundt hylder) ogsaa gælde for
Hjernefunktionerne. --- Naar der synes at blive ny Energi til paa det aandelige
Omraade, er det at forklare i Analogi med den Maade, paa hvilken Organismen
under sin Væxt kommer til at raade over større Energi. Det er den fremskridende
Differentiation af Organer og Funktioner, som gør en alsidigere Optagen og
Anvendelse af Energi mulig. Men herved brydes Ækvivalentsprincipet ikke. ---
Iøvrigt har jeg, som jeg mener, med tilstrækkeligt Eftertryk paavist, hvilke
Vanskeligheder en Gennemførelse af Ækvivalentsprincipet paa det aandelige
Omraade lider under. --- Det maa bemærkes, at Wundt (Ethik p. 412) indrømmer,
at hvor der har dannet sig en virkelig \emph{Karakter}, ere vi i Stand til ikke
blot bag efter at forklare, men ogsaa at forudsige et Menneskes Handlinger
Altsaa netop for den højeste Form for sjælelig Kausalitet gælder hin Paastand
ikke.} --- Det er naturligt, at Overbevisningen om
% --- p. 96 ---
\setcounter{footnote}{0}%
\opage{96}den faste Lovmæssighed danner sig senere for det aandelige Livs end
for den ydre Naturs Vedkommende. Og dog er det mest manglende Opmærksomhed, det
skyldes, naar det aandelige Liv frembyder Billedet af uregelmæssig Skiften, og
i Praxis stole vi i hvert Tilfælde nok saa meget paa den menneskelige Naturs
Lovmæssighed som paa den ydre Naturs. Vi ere jo ikke mindre afhængige af
Mennesker end af Naturen, og kunde ikke leve mellem Mennesker, naar vi ikke til
en vis Grad kunde forudsige deres Handlemaade.

40. Den psykologiske Forsken kan ikke stille sig anderledes overfor Paastanden
om Undtagelser fra Aarsagsloven for visse aandelige Fænomeners Vedkommende. Som
enhver real Forsken staar og falder den med Aarsagsprincipet. Det er da ogsaa
nu temmelig almindeligt indrømmet, at der rent teoretisk ingen Grund vilde være
til at hævde Indeterminismen. Derimod paastaas det endnu ofte, at moralske
Motiver skulle føre til paa visse Punkter at holde vor Forskerlyst i Tømme og
tage Forstanden fangen (om man ellers kan bruge disse Udtryk her; ti Opgivelsen
af Aarsagsprincipet skildres undertiden, som om den var den letteste Sag af
Verden og ingen Kamp kostede). \emph{Hermed flyttes}, hvad vel maa bemærkes,
det \emph{hele Spørgsmaal fra Psykologien over} i \emph{Etiken}.

Jeg maa her tillade mig for mit eget Vedkommende at henvise til, hvorledes jeg
i alle mine tidligere Udtalelser om det deterministiske Problem har hævdet
dette Synspunkt, at det er et etisk, ikke et psykologisk Spørgsmaal. At det er
et etisk Spørgsmaal, vil nu naturligvis ikke sige, at vi kunne appellere til de
gængse moralske Forestillinger, der ligesaa lidt i etisk som i psykologisk
Henseende ere tilstrækkelig klare og begrundede. Dersom Spørgsmaalet virkeligt
skal drøftes, maa man gaa tilbage til de Principer, der ligge til Grund for al
etisk Vurdering, og paa hvilke derfor al Tale om etisk Betydning og
% --- p. 97 ---
\setcounter{footnote}{0}%
\opage{97}Berettigelse maa hvile. Der maa spørges: hvilken etisk Betydning og
Berettigelse have de Følelser (særligt Anger- eller Skyldfølelsen), paa hvilke
man ofte grunder Nødvendigheden af at antage Villiens Aarsagsløshed? Og det maa
da undersøges, om Aarsagsløshedspostulatet virkeligt er en uundgaaelig logisk
Konsekvents af det, der udgør hine Følelsers etiske Værdi, eller om den
indeterministiske Paastand ikke meget godt kan bortfalde, uden at denne etiske
Værdi lider noget Afbræk. Hvis man ikke foretager en saadan Undersøgelse, er
det uden al Betydning at appellere til moralske Motiver.

I min første Afhandling om dette Punkt (Kapitlet om Villiens Frihed i
«Grundlaget for den humane Etik». 1876) søgte jeg allerede at vise, hvorledes
«Angeren bliver forstaaelig og ses i sin store etiske Betydning, selv om vi
ikke antage nogen absolut Frihed», og hvorledes den vilde være aldeles
overflødig, «dersom den skulde forudsætte, at vi ogsaa havde kunnet handle
anderledes under de givne Forhold» (p. 124). Dette Standpunkt har jeg siden
stedse lagt til Grund, senest i den nye Undersøgelse af dette Spørgsmaal i min
Etik (1887). Min Grundtanke er den: at den Forklaring af Angerfølelsen, som ud
fra Determinismen er mulig, ganske stemmer med den Betydning, som fra et etisk
Standpunkt maa tillægges Angeren. Efter min Opfattelse mødes Psykologi og Etik
paa den skønneste Maade i samme Resultat. Jeg \emph{forsvarer} derfor ikke
Determinismen. Men jeg \emph{angriber} visse gængse og overleverede
Forestillinger, der gaa ud fra et etisk uholdbart Begreb om Anger og
Tilregnelighed og derved (foruden andre Ulemper) frembringe Disharmoni mellem
den psykologiske og den etiske Betragtningsmaade.

Anger og Skyldfølelse have, som jeg har søgt at vise, kun Betydning ved at
vække Driften til Forandring og Fremgang i Villien. Som blot Rugen over, hvad
der ikke kan gøres om, ere de uetiske, idet de slappe Villien og kue Modet.
Angeren frembringes ved den skarpe Modstrid mellem hvad min egen Overbevisning
anerkender for ret, og hvad Erindringen foreholder mig om min egen Villens
Beskaffenhed. Naar man har ment, at ifølge Determinismen den Smerte, der føles
over at have villet uret, maa være væsentligt af samme Art som den, der føles
ved ikke at være smukkere eller bedre begavet, oversér man dels, at i Villien
finder Menneskets reale Selv sit mest potentserede Udtryk, saa at en
Selvfordømmelse her virker stærkest og dybest, dels, at medens man ikke kan
hjælpe synderligt paa sin Begavelse og sin Skønhed (især ikke paa sin Skønhed),
er derimod en Ændring af Villien ikke umulig, netop i Kraft af den psykologiske
Lovmæssighed. Derved faar Ulystfølelsen her en særligt aktiv Karakter, som den
æstetiske eller intellektuelle Misfornøjelse med sig selv savner. Angerens
Smerte og Uro skriver sig fra de to anførte Momenters Samvirken, med den
levende Følelse af det Rette som Baggrund. Angeren kan beskrives og forklares
lige saa naturligt som i det Hele nogen Følelse kan det med de Midler, der staa
til Psykologiens Raadighed. Og med selve denne Forklaring er egentligt den
etiske Betydning allerede given.

% --- p. 98 ---
\setcounter{footnote}{0}%
\opage{98}Siger nu den, om hvem en etisk forkastende Dom fældes: «Saaledes var
nu én Gang min Villie! Jeg kunde i det givne Øjeblik ikke ville anderledes!»
--- saa bliver Svaret følgende: «Derom er ikke Talen. Vi konstatere blot, at
den Villie, du lagde for Dagen, var slet. Og vi sige det til dig, for energisk
at henvende din Opmærksomhed derpaa, saa at din Villie kan blive en anden for
Fremtiden.» --- Vi have i det Hele ikke Ret til at udtale etiske Domme om
andre, uden for at vække en bedre Villen hos dem. \emph{Etiske Domme maa selv
være etisk motiverede}. Dette kunde synes at følge af sig selv, men bliver
(under den store Lyst til at moralisere og operere med moralske Argumenter) alt
for meget overset. Etiske Domme ere ikke berettigede blot som Udbrud, i hvilke
man skaffer sig Luft. Selv om det ikke er farisæisk Selvfølelse, der paa denne
Maade skaffer sig Afløb, saa overses dog, at det at være Genstand for etisk
Misbilligelse er en Lidelse, der --- som al Lidelse --- ikke maa paaføres nogen
uden som nødvendigt Middel til Fremgang. Dette følger ligefrem af de første
etiske Grundsætninger%
% footnote 1) on p. 98
\footnote[1]{Se min Etik (1887). p. 73---75.}. Enhver, som udtaler en etisk
Dom, paadrager sig derved et etisk Ansvar. Han opererer med Kræfter, som hører
til de stærkest og dybest bevægende i Verden, og som derfor ikke bør bruges i
Blinde. --- Overfor Sindssyge og Bevidstløse udtale vi derfor heller ingen
etiske Domme. Og den, der er kuet ved Følelsen af sin Afmagt til at holde sig
oppe i etisk Henseende, kue vi ikke endnu mere ved at rette etiske Domme til
ham; vi søge derimod netop at styrke hans Selvtillid og hans Selvagtelse. De
etiske Dommes teleologiske eller pædagogiske Karakter bliver her aabenbar.

Naar vi sige til et Menneske: «Saaledes burde du ikke have villet!» da hviler
denne Udtalelse ingenlunde paa indeterministiske Forudsætninger. Vi foreholde
ham kun Modsætningen mellem hans virkelige Villen og den Villen, han selv maa
anerkende som den rette. I denne skarpe Modsætning ligger en Brod, en stimulus
for Villien. ---

Denne Betragtning belyser klart den Paastand, at Determinisme og Etik ere
uforenelige. Den etiske Lov lever i Bevidstheden som Motiv, og etiske Domme
udtales --- paa etisk Maade --- for at blive Motiver. Og Determinismen er Læren
om Villiens fuldstændige Motiverethed. Det er da ikke godt at se, hvorledes der
her skulde være nogen Uforenelighed.

Problemet om «Villiens Frihed» er derfor aldeles ikke uløseligt. Det finder sin
klare og bestemte Afgørelse ved Fastsættelsen af de etiske Principer. Af disse
maa Indeterminismen enten følge eller ikke følge. Jeg undervurderer ikke
Vanskelighederne ved Fastsættelsen af disse Principer; jeg tør sige, at jeg har
søgt at stille disse Vanskeligheder i saa klart Lys som muligt; men hvis man i
det Hele vil lægge bestemte Principer til Grund ved sine etiske Reflexioner,
maa man dog kunne blive klar over, om Indeterminismen
% --- p. 99 ---
\setcounter{footnote}{0}%
\opage{99}er en Konsekvents af dem eller ikke. Og hvilken Værdi en Appel til
det Etiske kan have, som ikke lægger bestemte Principer til Grund, forstaar jeg ikke.

41. I den nye Behandling af Spørgsmaalet om Villiens Frihed, som Prof.
\emph{Kroman} har givet i 2. Udgave af sin Tænke- og Sjælelære, er nu ogsaa en
lignende Betragtningsmaade antydet. Han har her udtalt (p. 310), at
Determinismen og Indeterminismen hver gaa ud fra sit Skyldbegreb, og at det
Spørgsmaal staar tilbage: hvilket Skyldbegreb er det rette? --- Med dette
Spørgsmaal burde den hele Debat være begyndt. Den vilde da have været
frugtbarere. Men denne Betragtningsmaade betegner i hvert Tilfælde et stort
Fremskridt i Forhold til Behandlingen af samme Spørgsmaal i Forfatterens Skrift
«Vor Naturerkendelse» (Kbhvn. 1883), hvor Determinismens Mulighed kun
indrømmedes, hvis man forkaster Ansvarsfølelsen. Der gjordes i dette Skrift
ikke noget Forsøg paa at fremdrage den Mulighed, at det eneste etisk holdbare
Begreb om Ansvar og Anger kunde være et saadant, som netop Determinismen baade
kunde og maatte lægge til Grund. Ved «\emph{virkeligt} Ansvar, \emph{virkelig}
Skyld og Anger» (Vor Naturerkendelse. p. 246) forstodes i det tidligere Skrift
saadanne Former for disse Følelser, som kun Indeterminismen kunde anerkende.
--- Endnu i sin nye Fremstilling yder Forfatteren dog ikke lige Ret og Skel for
begge Parter. I det Mindste kunne hans Udtryk forstaas saaledes, som om det med
Determinismen forenelige Skyldbegreb var mere overfladisk end Indeterminismens.
Har man engang erklæret det for et etisk aabent Spørgsmaal, hvilket Skyldbegreb
der er det rette, saa maa man ogsaa tage alle Konsekventserne. Det vil da
navnligt være klart, at \emph{dersom} der skulde være etiske Betænkeligheder
ved Indeterminismens Skyldbegreb, saa blive disse Betænkeligheder des større,
jo «dybere» dette Begreb er.

Det har sikkert højligt undret de fleste Læsere af «Vor Naturerkendelse», at
Forfatteren i sin sidste Fremstilling protesterer saa bestemt mod at blive
betegnet som Indeterminist%
% footnote 1) on p. 99
\footnote[1]{I samme Note, i hvilken min ærede Kollega gør Indsigelse imod, at
jeg har regnet ham til Indeterministerne, replicerer han til nogle
Bemærkninger, jeg i min Etik har gjort om Indeterminismens Forhold til den
religiøse (særligt den teistiske) Verdensanskuelse. Jeg søgte at vise, at
Indeterminismen er uforenelig med Teismen, altsaa med den gængse religiøse
Opfattelse (til hvilken der var henvist), fordi den ophæver Guds Almagt og
Alvidenhed. Kromans Svar gaar ud paa følgende: 1) «Betegnelsen «almægtig»
betyder jo ikke, at Guddommen er \emph{eneste} Aarsag, men at alle andre
Aarsager ere satte og atter kunne hæves af ham». 2) Indeterministen kan kun da
ikke tro paa en alvidende Gud, «saafremt han antager Guddommens Viden
\emph{ensartet} med den menneskelige. Men en saadan Opfattelse næres nu om
Stunder vel neppe af ret mange». (p. 312.) Hertil maa jeg bemærke: 1) Almagt
kan hverken sætte eller ophæve \emph{noget}, \emph{som ingen Aarsag} har. Ti
baade at sætte og at ophæve er at være Aarsag, og Magtens Omraade strækker sig
kun saa langt, som Aarsager virke; gives der noget, der ingen Aarsag har, er
det eo ipso unddraget Magten, og denne er altsaa ikke Almagt. 2) Hvilken
Forestilling kunne vi danne os om en Viden, der ikke bestaar i Indsigt i Sagens
\emph{begrundede Sammenhæng}? Skal den guddommelige Viden være saa forskellig
fra den menneskelige, at den ikke bygger paa den tilstrækkelige Grunds Princip,
saa er det kun et Ord, vi bruge, ikke en Tanke, naar vi kalde den for Viden.
(Og vi befinde os jo ogsaa netop paa det Omraade, hvor Ordene pleje at træde i
Begrebernes Sted.) Jeg havde sagt: «Hvis man hævder, at vi her staa overfor et
«Mysterium», saa maa man se til, hvorledes man kan skelne mellem Mysterium og
Selvmodsigelse». Det forekommer mig ikke, at der ved Modbemærkningen er givet
noget Svar herpaa. Ogsaa min Protest imod, at Determinismen særligt skulde være
Udslag for en antireligiøs Retning, maa jeg opretholde. Kirkens Historie viser,
at de egentligt religiøse Naturer stedse have hældet til Determinismen, medens
Naturer, der førtes mere i etisk end i religiøs Retning, (ved en
Misforstaaelse) hævdede Indeterminismen. Ikke al Determinisme er religiøs; men
al sand Religiøsitet maa være deterministisk, idet den maa hævde den inderlige
Enhed mellem Mennesket og dets Gud, saa at Mennesket føler sig gennemstraalet
af Guddommen og vedkender sig kun at kunne ville og udrette, hvad den
guddommelige Virken fuldbyrder i dets Indre. ---}. Hele Forsvaret for
Indeterminismen, og ikke mindst den Belysning, i hvilken
% --- p. 100 ---
\setcounter{footnote}{0}%
\opage{100}Determinismen stilles (se især «Vor Naturerk.» p. 250---253), er af
den Beskaffenhed, at det i hvert Tilfælde er undskyldeligt, om man har ment, at
Forfatteren stod i et langt inderligere Forhold til Indeterminismen end det at
være dens Advokat. Imidlertid maa han jo være sin egen bedste Fortolker. Og
Hovedsagen er, at den nye Fremstilling i hvert Tilfælde fremhæver noget, der
ikke var udtalt i den tidligere. --- Der er nogle Punkter af almindelig
Interesse, som jeg her skal gaa noget nærmere ind paa.

I den nye Fremstilling udtales der, at den største Del af Villiens Psykologi
kan fremstilles, uden at man behøver at komme ind paa det deterministiske
Problem. (Se p. 304, 312.) Tidligere synes Forfatteren ikke at have været saa
sikker herpaa. I «Vor Naturerkendelse» (p. 244) var han ikke uden Betænkelighed
ved, at en Determinist brugte Udtryk som f. Ex. at «naa». Dersom ethvert
Udtryk, der betegner Drift og Higen, skulde indeholde hele Problemet i sig, saa
maatte Forfatterens egen Fremstilling af Villiens Psykologi have faaet en helt
anden Karakter. Det er til afgjort Fordel for hans Psykologi, at han har
overvundet denne Betænkelighed. --- Jeg fremdrager her dette Punkt kun for
dertil at knytte en Bemærkning. Prof. Kroman sætter den Omstændighed, at saa
stor en Del af Psykologien kan behandles, uden at Spørgsmaalet om «Friheden»
behøver at komme frem, i Forbindelse med den psykologiske Forskens ringe
Exakthed. Men naar man reducerer Indeterminismen, eller rettere
Indeterminationen, til en saa forsvindende lille Størrelse, som Kroman gør, saa
vil man --- uden Fare for at blive gendreven --- kunne indføre den selv i de
mest exakte Erfaringsvidenskaber. Ingen Kemiker vil nogensinde kunne modbevise
den Paastand, at naar Vand dannes af Ilt og Brint, skabes en Billiontedel af
intet, eller at denne Vanddannelse ingenlunde er en kontinuerlig Proces, men at
i et forsvindende lille Øjeblik (f. Ex. en Billiontedel af et Sekund)
tilintetgøres Ilten og Brinten, og der skabes strax Vand i Stedet for. Men den
Slags Paastande har enhver Erfaringsvidenskab, ogsaa Psykologien, Ret til
roligt at lade ligge, overladende Bevisbyrden til dem, der ønske at hævde dem.

% --- p. 101 ---
\setcounter{footnote}{0}%
\opage{101}Det er Selvmodsigelsen i min ærede Kollegas Standpunkt i dette
Spørgsmaal, at han har reduceret Indeterminismen til noget saa Forsvindende, at
den hverken har nogen teoretisk eller praktisk Betydning længere. Imidlertid
søger han i sin nye Fremstilling at vise, at den saare ringe Del af
Villiesakten, der efter hans Teori skal være (eller kunne antages at være)
aarsagsløs, dog ikke er uden Betydning. --- Den Del af Villiesakten, som ikke
er aarsagsbestemt, er (hedder det p. 308) dog ikke «opdukket af intet»: Jeget
anvender blot vilkaarligt en vis Del af den Energi, som det i det Hele sidder
inde med. --- Hertil maa jeg svare følgende.

Man kan kun ved at benytte et Billedsprog skelne mellem Jeget og dets Energi.
Jeget selv ytrer sig i en mere eller mindre kraftig Virksomhed, og det er netop
denne Virksomheds Vækkelse i det rette Øjeblik og i den rette Retning, det
kommer an paa. Hvis man nu vil sige, at der maa skelnes mellem Jeget selv og
dets Virksomhed, saa kan Jeget kun betyde den potentielle Energi, der kan gaa
over til aktuel Energi; men selve denne Overgang kan ikke foregaa, uden at en
Udløsning finder Sted. Hvis man negter, at der behøves en Impuls, for at denne
Udløsning skal finde Sted, paastaar man eo ipso, at noget «opdukker af intet»,
nemlig selve Energiens Udløsning i dette Øjeblik og i denne Retning. Selv om
man lægger nok saa stor Vægt paa den menneskelige Villies Energi (og det kan
Deterministen helt vel gøre), maa man dog ikke glemme, at i det Mindste «et
Differentiale af et Stød» stedse behøves, for at denne Energi skal blive
anvendt. Uden saadanne Impulser kan vor Villie lige saa lidt virke, som vort
organiske Liv kan bestaa uden Aandedræt. Impulserne behøve ikke i hvert Øjeblik
at komme udefra. I Kraft af de almindelige psykologiske Love dels om Forholdet
mellem Forestillingerne indbyrdes, dels om Forholdet mellem Forestillinger og
Følelser, ville udløsende Motiver kunne udvikle sig af selve Jegets Indre. I
denne, men ogsaa kun i denne Betydning kan man sige, at Jeget udløser sin egen
Energi. Men det maa da ikke glemmes, at al Forestillingsvirksomhed og
Følelsesbevægelse selv er en Anvendelse af Energi og ikke forløber uafhængigt
af en Villen. Vi komme, saa længe vi tænke klart og bestemt paa psykologisk
Vis, aldrig udenfor et System af Vexelvirkninger dels mellem indre og ydre
Betingelser, dels mellem indre Betingelser indbyrdes. Kun ved at holde
Forestillingen om Jeget hen i mystisk Ubestemthed kan dette undgaas.

Naar det siges, at Jeget «\emph{vilkaarligt}» anvender en Del af sin Energi,
hvad forstaas da ved Ordet «vilkaarligt»? Dersom det tages i den i Psykologien
sædvanlige Betydning (smlgn. 29), betyder det, at der i Bevidstheden forud for
Akten gør sig en Drift eller en Beslutning gældende. Men hele Spørgsmaalet
vender jo saa tilbage, idet der maa spørges, om denne Drift eller Beslutning
opstaar efter psykologiske Love eller ikke. ---

En etisk Betænkelighed vækker den indeterministiske Teori paa dette Punkt, idet
den frister til at stole paa de ubestemte Muligheder, som vi ikke forud kunne
gøre os til
% --- p. 102 ---
\setcounter{footnote}{0}%
\opage{102}Herrer over. Om hin ubegribelige Minimaltilgift til den naturligt
udviklede Villen vil indtræde i et givet Øjeblik, raade vi jo forud ikke for.
Hvis vi kunde bestemme vor egen fremtidige Villen, vilde denne jo blive
determineret! Determinismen opfordrer os derimod alvorligt til at danne en god
og fast Grundvold, idet den viser os, at enhver af vore Villiesakter, ja enhver
af vore Tanker og Følelser lægger en Sten til Bygningen af vor fremtidige
Karakter, og at vi ikke i hvert Øjeblik kunne begynde forfra. Medens
Indeterminismen lærer, at vi kunne høste uden at saa (og uden at der i det Hele
taget er saat), indskærper Determinismen, at som vi have saat, ville vi ogsaa høste.

I den nye Udgave af «Tænke- og Sjælelære» gøres som Svar paa denne Indvending
følgende Bemærkning. «Hver Gang Subjektet handler paa en vis Maade, vil dets
totale Beskaffenhed ændres en Kende i samme Retning, og det virker saaledes
daglig selv med til Dannelsen af den Karakter, hvis øjeblikkelige Beskaffenhed
giver den enkelte Handling dens væsentligste Særpræg» (p. 308). Vi staa da her
ved det Spørgsmaal, om Indeterministen kan anerkende Øvelsens Lov for Villiens
Vedkommende. Jeg kan ikke indse, hvorledes det er muligt uden Selvmodsigelse.
Enhver aarsagsløs Villen (eller aarsagsløs Del af en Villen) er ifølge Sagens
Natur ikke bestemt ved noget som helst Forudgaaende. Selv om en hel vanemæssig
Tendents paa en eller anden Maade er frembragt, kunne de indeterminerede
Tilgifter i hvert Øjeblik slaa ind i nye Retninger. Det er et \emph{rent
Tilfælde}, om en saadan minimal Tilgift virker i samme Retning som en
tidligere. \emph{Mellem de indeterminerede Tilgifter indbyrdes} kan der
\emph{intet Kausalitetsforhold} finde \emph{Sted}. Hvis min «frie» Villen om en
Time lettere gik i en vis bestemt Retning, fordi min «frie» Villen i det
nærværende Øjeblik gaar i denne Retning, saa vilde den ikke være «fri».

Lige saa lidt kan jeg indse, hvorledes man ud fra indeterministiske
Forudsætninger kan hævde en vis Ansvarlighed for de indeterminerede Tilgifters
Virkninger. Disse Virkninger ere jo dog nødvendige, og dersom Ansvar og
Aarsagsløshed skulle staa og falde med hinanden, saa kan Ansvaret heller ikke
strække sig længere end selve den aarsagsløse Villen gaar. Saa snart
Aarsagsloven træder i Kraft efter det indeterministiske Intermezzo, og det gør
den, saa snart Villiesaktens Virkninger udvikle sig, maa Ansvaret ogsaa holde op.

Man vil maaske hertil svare, at den deterministiske Etik jo dog gør os
ansvarlige for mere end vi direkte have villet, ogsaa for vore Beslutningers og
Handlingers Virkninger, som vi jo dog ikke behøve at have villet, skønt vi
vilde selve Handlingen. Denne Replik vilde oversé, at hvad der (efter den
deterministiske Etik) lægges os til Last, og hvad der indskærpes os med
Eftertryk gennem Straffen, naar Talen er om Virkninger, der ikke vare
tilsigtede, er, at vi have handlet kortsynet og letsindigt. Kun saadanne
Virkninger bør vi derfor drages til Ansvar for, som vi muligvis kunne komme til
at forudse og til ved grundigere Overvejelse at undgaa. Det er derfor aldeles
konsekvent, naar Ansvaret i den
% --- p. 103 ---
\setcounter{footnote}{0}%
\opage{103}deterministiske Etik udstrækkes til at omfatte ogsaa fjernere
Virkninger af Handlingen. Men for Indeterminismen er denne Mulighed afskaaren,
da den én Gang for alle har knyttet sit Skyldbegreb til Postulatet om
Aarsagsløsheden. Skal det afgørende Moment i Villiesafgørelsen være
aarsagsløst, saa kan det ikke nogensinde bestemmes ved nogen Tænken over
Villiesaktens mulige Virkninger. Enhver Berettigelse til at gøre Ansvar
gældende for disse, kommer altsaa til at mangle.

\begin{center}\rule{2cm}{0.4pt}\end{center}

42. Den Række af Spørgsmaal, som ere drøftede i denne Afhandling, ere ikke
vilkaarligt sammenstillede. Skridt for Skridt har det ene Spørgsmaal ført os
til det andet. Den ledende Traad har været Sjælelivets lovmæssige Sammenhæng og
de Former, under hvilke den kundgør sig paa Forestillingernes og Villiens
Omraade. Den Mulighed er, som jeg ogsaa har gjort opmærksom paa, ikke
udelukket, at den, der er enig med mig i et Punkt, derfor ikke er enig med mig
i alle andre Punkter. Man kan f. Ex. give mig Ret hvad Teorien om Genkendelse
angaar, og dog være uenig med mig med Hensyn til Forklaringen af
Lighedsassociationen. Men dermed bortfalder ikke Berettigelsen af at forbinde
de foregaaende Undersøgelser til et Hele, selv om jeg maa indrømme, at den ydre
Anledning til, at netop disse Spørgsmaal her behandledes sammen, var den, at
Forskere, for hvem jeg nærede Agtelse, med Hensyn til dem vare komne til andre
Resultater end jeg. Der er herved kommet et dialogisk Moment ind i
Behandlingen, som ved Drøftelsen af denne Slags Spørgsmaal ikke er til Skade. I
det Mindste indrømmer jeg villigt, at flere af de Indvendinger, der ere gjorte
mod mine Synsmaader, have givet mig Anledning til at udvikle dem til større
Klarhed og Nøjagtighed, eller dog at forsøge derpaa.

\begin{center}\rule{2cm}{0.4pt}\end{center}

\end{document}
